% Séries statistiques regroupées en classes de même amplitude — Mathématiques 1re Tech — Cours signé OnBuch+
% Programme officiel MINESEC. Compiler avec : tectonic N05-series-statistiques-regroupees-classes-meme-amplitude.tex
\newcommand{\DOCMATIERE}{Mathématiques}
\newcommand{\DOCNIVEAU}{1re Tech}
\newcommand{\DOCMODULE}{Mathématiques – classe de Première technique}
\newcommand{\DOCLECON}{N05}
\newcommand{\DOCTITRE}{Séries statistiques regroupées en classes de même amplitude}
% OnBuch+ — préambule « cours signés » ENRICHI (compilé avec tectonic / XeLaTeX)
% Système visuel « Pop » d'OnBuch+ : fond crème (jamais blanc), bordures encre
% épaisses, ombres « dures » décalées, titres Archivo Black en capitales.
% Version MATHÉMATIQUES : graphes (pgfplots/tikz), boîtes pédagogiques
% (définition, propriété, méthode, attention, le savais-tu, exercice résolu,
% exercice, TP, bilan), page de garde et signature OnBuch+ ancrée en bas.
%
% Macros attendues dans main.tex AVANT \input{preamble} :
%   \DOCMATIERE  \DOCNIVEAU  \DOCMODULE  \DOCLECON (numéro)  \DOCTITRE
\documentclass[11pt]{article}

\usepackage{fontspec}
\usepackage[french]{babel}
\usepackage[a4paper,margin=2cm,top=3.4cm,bottom=2.8cm,headheight=1.4cm,footskip=1.3cm]{geometry}
\usepackage{amsmath,amssymb,amsfonts}
\usepackage{mathtools} % \xmapsto, \coloneqq... souvent utilisés par les modèles
\usepackage{cancel} % simplifications barrées
% \abs{z} (module, valeur absolue) et \norm{u} (norme), utilisés par les modèles
\providecommand{\abs}{}\let\abs\relax\DeclarePairedDelimiter{\abs}{\lvert}{\rvert}
\providecommand{\norm}{}\let\norm\relax\DeclarePairedDelimiter{\norm}{\lVert}{\rVert}
\usepackage[table]{xcolor} % [table] : fournit aussi \rowcolors (alternance de lignes)
\usepackage{pagecolor}
\usepackage{tikz}
\usetikzlibrary{arrows.meta,positioning,calc,shapes.geometric,shapes.misc,shapes.arrows,decorations.pathmorphing,decorations.pathreplacing,decorations.markings,patterns,shadows,mindmap,trees,fit,backgrounds,matrix,intersections,angles,quotes}
\usepackage{pgfplots}
\usepackage{pgf-pie} % diagrammes circulaires \pie (statistiques)
\pgfplotsset{compat=1.18}
\usepgfplotslibrary{fillbetween,groupplots} % aires entre courbes (calcul intégral)
\usepackage{enumitem}
\usepackage{fancyhdr}
% [most] charge toutes les bibliothèques tcolorbox (listings, minted, raster,
% documentation...) et gonfle inutilement la mémoire TeX sur un document long
% (~300 boîtes) : on ne charge que ce qui est réellement utilisé.
\usepackage[skins,breakable]{tcolorbox}
\usepackage{graphicx}
\usepackage{colortbl}
\usepackage{booktabs}
\usepackage{tabularx}
\usepackage{diagbox} % case d'en-tête diagonale des tableaux à double entrée
% Colonnes extensibles centrée / gauche / droite (C, L, R), souvent utilisées par les modèles
\newcolumntype{C}{>{\centering\arraybackslash}X}
\newcolumntype{L}{>{\raggedright\arraybackslash}X}
\newcolumntype{R}{>{\raggedleft\arraybackslash}X}
\usepackage{array}
\usepackage{multicol}
\usepackage{siunitx}
% parse-numbers=false : accepte aussi \SI{2,5 \times 10^{-3}}{...}
\sisetup{locale=FR,output-decimal-marker={,},group-minimum-digits=4,parse-numbers=false}
\DeclareSIUnit\ohm{\ensuremath{\symup{\Omega}}} % U+2126 absent de la police
\usepackage{qrcode}
% \degree : commande fréquemment utilisée par les modèles pour un angle ou une
% température en dehors de \SI{}{} (non fournie par les paquets chargés).
\providecommand{\degree}{^{\circ}}
% \qed (fin de démonstration), utilisé par les modèles sans amsthm
\providecommand{\qed}{\hfill$\square$}
% \barre : double barre des valeurs interdites dans un tableau de variations (tkz-tab)
\providecommand{\barre}{\ensuremath{\|}}
% environnement proof (amsthm non chargé) utilisé par les modèles
\ifdefined\proof\else\newenvironment{proof}[1][Démonstration]{\par\noindent\textit{#1.}\ }{\qed\par}\fi
\usepackage{lastpage}
\usepackage{hyperref}
\hypersetup{hidelinks}

% ============================================================
% Palette « Pop » OnBuch+ — mêmes hex que la classe P (plus_ui.dart)
% ============================================================
\definecolor{poporange}{HTML}{F59321}
\definecolor{popdark}{HTML}{E07A0C}
\definecolor{popcream}{HTML}{FFF6E8}
\definecolor{popink}{HTML}{1C1714}
\definecolor{poppurple}{HTML}{7A5AE0}
\definecolor{popgreen}{HTML}{2E9E6A}
\definecolor{poppink}{HTML}{E0457B}
\definecolor{popblue}{HTML}{2F7DE1}
\definecolor{popgold}{HTML}{C98A12}
\definecolor{popmuted}{HTML}{8A7F76}
\definecolor{popline}{HTML}{EADFD2}
\definecolor{popgray}{HTML}{8A7F76}
\colorlet{popgrey}{popgray}
% Alias tolérants (le modèle nomme parfois les couleurs différemment)
\colorlet{popwhite}{white}
\colorlet{popblack}{popink}
\colorlet{popred}{poppink}
\colorlet{popyellow}{popgold}
% Teintes claires pour fonds de tableaux / zones
\colorlet{popblueL}{popblue!10!white}
\colorlet{poporangeL}{poporange!14!white}
\colorlet{popgreenL}{popgreen!12!white}
\colorlet{poppurpleL}{poppurple!12!white}
\colorlet{poppinkL}{poppink!12!white}
\colorlet{popgoldL}{popgold!14!white}

\pagecolor{popcream}

% ============================================================
% Typographie
% ============================================================
\defaultfontfeatures{Ligatures=TeX}
\setmainfont{PlusJakartaSans-Medium.ttf}[Path=../../../fonts/,BoldFont=PlusJakartaSans-Bold.ttf]
% Maths en sans-serif (Fira Math) avec chiffres et lettres droites de Plus
% Jakarta, pour que les formules se fondent dans le texte.
\usepackage[math-style=ISO,bold-style=ISO]{unicode-math}
\setmathfont{FiraMath-Regular.otf}[Path=../../../fonts/]
\setmathfont{PlusJakartaSans-Medium.ttf}[Path=../../../fonts/,range={up/num,up/{Latin,latin},\mathup}]
\usepackage{icomma} % virgule décimale sans espace en mode math
% Caractères mathématiques Unicode absents des polices de texte (Plus Jakarta,
% Archivo Black) : rendus par leur équivalent mathématique, en texte comme en
% formule — sinon ils s'affichent en carré vide (ex. « Arithmétique dans ℤ »).
\usepackage{newunicodechar}
\newunicodechar{ℝ}{\ensuremath{\mathbb{R}}}
\newunicodechar{ℤ}{\ensuremath{\mathbb{Z}}}
\newunicodechar{ℂ}{\ensuremath{\mathbb{C}}}
\newunicodechar{ℕ}{\ensuremath{\mathbb{N}}}
\newunicodechar{ℚ}{\ensuremath{\mathbb{Q}}}
\newunicodechar{𝔻}{\ensuremath{\mathbb{D}}}
\newunicodechar{α}{\ensuremath{\alpha}}
\newunicodechar{β}{\ensuremath{\beta}}
\newunicodechar{γ}{\ensuremath{\gamma}}
\newunicodechar{δ}{\ensuremath{\delta}}
\newunicodechar{ε}{\ensuremath{\varepsilon}}
\newunicodechar{θ}{\ensuremath{\theta}}
\newunicodechar{λ}{\ensuremath{\lambda}}
\newunicodechar{μ}{\ensuremath{\mu}}
\newunicodechar{σ}{\ensuremath{\sigma}}
\newunicodechar{τ}{\ensuremath{\tau}}
\newunicodechar{φ}{\ensuremath{\varphi}}
\newunicodechar{ψ}{\ensuremath{\psi}}
\newunicodechar{ω}{\ensuremath{\omega}}
\newunicodechar{Σ}{\ensuremath{\Sigma}}
% majuscules produites par \MakeUppercase dans les titres (α-aminés -> Α)
\newunicodechar{Α}{\ensuremath{\alpha}}
\newunicodechar{Β}{\ensuremath{\beta}}
\newunicodechar{Γ}{\ensuremath{\Gamma}}
\newunicodechar{Δ}{\ensuremath{\Delta}}
\newunicodechar{Ε}{\ensuremath{\varepsilon}}
\newunicodechar{Θ}{\ensuremath{\Theta}}
\newunicodechar{Λ}{\ensuremath{\Lambda}}
\newunicodechar{Μ}{\ensuremath{\mu}}
\newunicodechar{Τ}{\ensuremath{\tau}}
\newunicodechar{Φ}{\ensuremath{\Phi}}
\newunicodechar{Ψ}{\ensuremath{\Psi}}
\newunicodechar{Ω}{\ensuremath{\Omega}}
\newunicodechar{↦}{\ensuremath{\mapsto}}
\newunicodechar{∈}{\ensuremath{\in}}
\newunicodechar{∉}{\ensuremath{\notin}}
\newunicodechar{∧}{\ensuremath{\wedge}}
\newunicodechar{⇔}{\ensuremath{\Leftrightarrow}}
\newunicodechar{⟺}{\ensuremath{\Longleftrightarrow}}
\newunicodechar{⇒}{\ensuremath{\Rightarrow}}
\newunicodechar{⟹}{\ensuremath{\Longrightarrow}}
\newunicodechar{∘}{\ensuremath{\circ}}
\newunicodechar{∪}{\ensuremath{\cup}}
\newunicodechar{∩}{\ensuremath{\cap}}
\newunicodechar{≡}{\ensuremath{\equiv}}
\newunicodechar{⊂}{\ensuremath{\subset}}
\newunicodechar{⊥}{\ensuremath{\perp}}
\newunicodechar{∃}{\ensuremath{\exists}}
\newunicodechar{∀}{\ensuremath{\forall}}
\newunicodechar{∛}{\ensuremath{\sqrt[3]{\,}}}
\newunicodechar{∜}{\ensuremath{\sqrt[4]{\,}}}
\newunicodechar{⃗}{\ensuremath{{}^{\to}}}
\newunicodechar{ⁿ}{\ifmmode^{n}\else\textsuperscript{\ensuremath{n}}\fi}
\newunicodechar{ˣ}{\ifmmode^{x}\else\textsuperscript{\ensuremath{x}}\fi}
\newunicodechar{ᵇ}{\ifmmode^{b}\else\textsuperscript{\ensuremath{b}}\fi}
\newunicodechar{ʳ}{\ifmmode^{r}\else\textsuperscript{\ensuremath{r}}\fi}
\newunicodechar{ᵘ}{\ifmmode^{u}\else\textsuperscript{\ensuremath{u}}\fi}
\newunicodechar{ʸ}{\ifmmode^{y}\else\textsuperscript{\ensuremath{y}}\fi}
\newunicodechar{ᵃ}{\ifmmode^{a}\else\textsuperscript{\ensuremath{a}}\fi}
\newunicodechar{ᵉ}{\ifmmode^{e}\else\textsuperscript{\ensuremath{e}}\fi}
\newunicodechar{ᵏ}{\ifmmode^{k}\else\textsuperscript{\ensuremath{k}}\fi}
\newunicodechar{ᵖ}{\ifmmode^{p}\else\textsuperscript{\ensuremath{p}}\fi}
\newunicodechar{ᴺ}{\ifmmode^{N}\else\textsuperscript{\ensuremath{N}}\fi}
\newunicodechar{ᶜ}{\ifmmode^{c}\else\textsuperscript{\ensuremath{c}}\fi}
\newunicodechar{ʰ}{\ifmmode^{h}\else\textsuperscript{\ensuremath{h}}\fi}
\newunicodechar{ᵠ}{\ifmmode^{q}\else\textsuperscript{\ensuremath{q}}\fi}
\newunicodechar{ₙ}{\ifmmode_{n}\else\textsubscript{\ensuremath{n}}\fi}
\newunicodechar{ₐ}{\ifmmode_{a}\else\textsubscript{\ensuremath{a}}\fi}
\newunicodechar{ᵢ}{\ifmmode_{i}\else\textsubscript{\ensuremath{i}}\fi}
\newunicodechar{ᵣ}{\ifmmode_{r}\else\textsubscript{\ensuremath{r}}\fi}
\newunicodechar{ₖ}{\ifmmode_{k}\else\textsubscript{\ensuremath{k}}\fi}
\newunicodechar{ᵦ}{\ifmmode_{\beta}\else\textsubscript{\ensuremath{\beta}}\fi}
\newunicodechar{ₓ}{\ifmmode_{x}\else\textsubscript{\ensuremath{x}}\fi}
\newfontfamily\popdisplay{ArchivoBlack-Regular.ttf}[Path=../../../fonts/]
\newfontfamily\poplogo{SpaceGrotesk-Bold.ttf}[Path=../../../fonts/]
\newfontfamily\popsemi{PlusJakartaSans-SemiBold.ttf}[Path=../../../fonts/]
\newfontfamily\popxbold{PlusJakartaSans-ExtraBold.ttf}[Path=../../../fonts/]
\newfontfamily\popxboldls{PlusJakartaSans-ExtraBold.ttf}[Path=../../../fonts/,LetterSpace=3.0]

\setlist[enumerate]{leftmargin=1.7em,itemsep=5pt,topsep=4pt}
\setlist[itemize]{leftmargin=1.7em,itemsep=4pt,topsep=4pt,label={\color{poporange}\textbullet}}
\setlength{\parindent}{0pt}
\setlength{\parskip}{5pt}
\renewcommand{\arraystretch}{1.25}

% pgfplots : style Pop par défaut
\pgfplotsset{
  every axis/.append style={axis line style={popink,line width=1pt},tick style={popink},
    grid style={popline,line width=0.5pt},label style={font=\small},tick label style={font=\footnotesize}},
  popcurve/.style={poporange,line width=1.8pt,no markers,smooth},
}
% Même style disponible dans un \draw TikZ simple (hors pgfplots)
\tikzset{popcurve/.style={poporange,line width=1.8pt}}

% ============================================================
% Pill / squiggle
% ============================================================
\newcommand{\poppill}[3]{%
  \tikz[baseline=-0.55ex]\node[draw=popink,line width=1.2pt,rounded corners=2.6pt,%
    fill=#2,inner xsep=6.5pt,inner ysep=3pt]{%
    \popxboldls\fontsize{7}{7}\selectfont\color{#3}\MakeUppercase{#1}};%
}
\newcommand{\popsquiggle}[1]{%
  \tikz[baseline=-0.4ex]{
    \draw[#1,line width=2.6pt,line cap=round]
      plot [smooth] coordinates {(0,0.09) (0.34,0.01) (0.68,0.21) (1.02,0.01) (1.36,0.10)};
  }%
}
% Mot-clé surligné (marqueur orange sous le texte)
\newcommand{\cle}[1]{{\popxbold\color{popdark}#1}}

% ============================================================
% En-tête / pied
% ============================================================
\pagestyle{fancy}
\fancyhf{}
\renewcommand{\headrulewidth}{0pt}
\renewcommand{\footrulewidth}{0pt}
\lhead{\raisebox{-0.15ex}{\poplogo\fontsize{13}{13}\selectfont\color{popink}OnBuch\color{poporange}+}}
\rhead{\poppill{\DOCMATIERE\ \textbullet\ \DOCNIVEAU\ \textbullet\ Leçon \DOCLECON}{white}{popink}}
\lfoot{\popxbold\fontsize{7.6}{7.6}\selectfont\color{popmuted}\MakeUppercase{Cours signé OnBuch+}\ \textbullet\ {\popsemi\DOCTITRE}}
\rfoot{\tikz[baseline=-0.6ex]\node[rounded corners=8.5pt,fill=popink,minimum height=17pt,minimum width=17pt,inner xsep=5pt,inner sep=0]{\popxbold\fontsize{8}{8}\selectfont\color{popcream}\thepage\,/\,\pageref*{LastPage}};}

% ============================================================
% Titres
% ============================================================
\newcommand{\chaptitre}[2]{%
  \begin{tcolorbox}[enhanced,colback=poporange,colframe=popink,boxrule=2.6pt,arc=11pt,
      drop shadow=popink,left=15pt,right=15pt,top=12pt,bottom=11pt,before skip=3pt,after skip=18pt]
    \poppill{#2}{popink}{popcream}\par\vspace{9pt}
    {\raggedright\hyphenpenalty=10000\exhyphenpenalty=10000\popdisplay\fontsize{22}{24}\selectfont\color{popcream}\MakeUppercase{#1}\par}
  \end{tcolorbox}
}
% Section numérotée : pastille orange avec le numéro + titre Archivo
\newcounter{popsec}
\newcommand{\coursec}[1]{%
  \stepcounter{popsec}\setcounter{popsub}{0}%
  \par\vspace{16pt}\noindent
  \begin{minipage}{\linewidth}
  \tikz[baseline=-0.7ex]\node[draw=popink,line width=1.6pt,fill=poporange,rounded corners=4pt,minimum size=22pt,inner sep=0]{\popdisplay\fontsize{12}{12}\selectfont\color{popcream}\thepopsec};%
  \hspace{8pt}{\popdisplay\fontsize{15}{17}\selectfont\color{popink}\MakeUppercase{#1}}\par
  \vspace{4pt}\noindent{\color{poporange}\rule{36pt}{3.6pt}}
  \end{minipage}\par\nopagebreak\vspace{8pt}
}
\newcounter{popsub}
\newcommand{\courssub}[1]{%
  \stepcounter{popsub}%
  \par\vspace{9pt}\noindent{\popxbold\fontsize{11.5}{13}\selectfont\color{popink}{\color{poporange}\thepopsec.\thepopsub}\hspace{6pt}#1}\par\nopagebreak\vspace{4pt}
}

% ============================================================
% Boîtes pédagogiques — même recette : fond blanc, bordure épaisse colorée,
% ombre dure encre, badge plein. Titre optionnel [#1].
% ============================================================
\tcbset{popbase/.style={breakable,enhanced,colback=white,boxrule=2.3pt,arc=9pt,
  shadow={3pt}{-3pt}{0pt}{popink},left=14pt,right=14pt,top=11pt,bottom=11pt,before skip=11pt,after skip=15pt,
  fontupper=\popsemi\fontsize{10.6}{14.5}\selectfont\color{popink},
  fontlower=\popsemi\fontsize{10.6}{14.5}\selectfont\color{popink}}}
% Coche dessinée (le glyphe ✓ n'existe pas dans les polices du document)
\newcommand{\popcheck}[1]{\tikz[baseline=-0.5ex]\draw[#1,line width=1.6pt,line cap=round,line join=round](-0.13,0.0)--(-0.04,-0.09)--(0.13,0.1);}
\newcommand{\popboxhead}[3]{\poppill{#1}{#2}{white}\ifx\relax#3\relax\else\hspace{6pt}{\popxbold\fontsize{10.6}{12}\selectfont\color{popink}#3}\fi\par\vspace{7pt}}

\newtcolorbox{aretenir}[1][]{popbase,colframe=popgreen,before upper={\popboxhead{À retenir}{popgreen}{#1}}}
\newtcolorbox{exemplebox}[1][]{popbase,colframe=poporange,before upper={\popboxhead{Exemple}{poporange}{#1}}}
\newtcolorbox{definition}[1][]{popbase,colframe=popblue,before upper={\popboxhead{Définition}{popblue}{#1}}}
\newtcolorbox{propriete}[1][]{popbase,colframe=poppurple,before upper={\popboxhead{Propriété}{poppurple}{#1}}}
\newtcolorbox{methode}[1][]{popbase,colframe=popgold,before upper={\popboxhead{Méthode}{popgold}{#1}}}
\newtcolorbox{attention}[1][]{popbase,colframe=poppink,before upper={\popboxhead{Attention}{poppink}{#1}}}
\newtcolorbox{experience}[1][]{popbase,colframe=popblue,colback=popblueL,before upper={\popboxhead{Expérience}{popblue}{#1}}}
% « Le savais-tu ? » : carte violette pleine (contexte, culture, Cameroun)
\newtcolorbox{savaistu}[1][]{popbase,colback=poppurple,colframe=popink,
  fontupper=\popsemi\fontsize{10.4}{14.2}\selectfont\color{white},
  before upper={\poppill{Le savais-tu\,?}{white}{poppurple}\ifx\relax#1\relax\else\hspace{6pt}{\popxbold\color{white}#1}\fi\par\vspace{7pt}}}
% Exercice résolu : énoncé en haut, \tcblower puis la solution sur fond crème
\newcounter{popexo}\newcounter{popexr}
\newtcolorbox{exoresolu}[1][]{popbase,colframe=poporange,colbacklower=popcream,
  segmentation style={popink,line width=1.2pt,dashed},
  before upper={\stepcounter{popexr}\popboxhead{Exercice résolu \thepopexr}{poporange}{#1}},
  before lower={\poppill{Solution}{popink}{popcream}\par\vspace{6pt}}}
% Exercice (non résolu en place) — la correction est dans la partie Corrigés
\newtcolorbox{exercice}[1][]{popbase,colframe=popink,
  before upper={\stepcounter{popexo}\popboxhead{Exercice \thepopexo}{popink}{#1}}}
% Corrigé
\newtcolorbox{corrige}[1][]{popbase,colframe=popgreen,colback=popgreenL,
  before upper={\popboxhead{Corrigé}{popgreen}{#1}}}
% Objectifs / prérequis / bilan
\newtcolorbox{objectifs}{popbase,colframe=popink,colback=poporangeL,
  before upper={\popboxhead{Objectifs de la leçon}{popink}{}}}
\newtcolorbox{prerequis}{popbase,colframe=popink,colback=popblueL,
  before upper={\popboxhead{Prérequis}{popblue}{}}}
\newtcolorbox{bilan}{popbase,colframe=popink,colback=white,boxrule=2.8pt,
  before upper={\popboxhead{Fiche bilan}{poporange}{L'essentiel à connaître pour l'examen}}}
% Figure encadrée avec légende
\newtcolorbox{popfigure}[1][]{enhanced,breakable=false,colback=white,colframe=popline,boxrule=1.4pt,arc=8pt,
  left=10pt,right=10pt,top=10pt,bottom=8pt,before skip=10pt,after skip=12pt,halign=center,
  lower separated=false,halign lower=center,
  fontlower=\popsemi\fontsize{9}{11.5}\selectfont\color{popmuted},
  #1}
\newcommand{\legende}[1]{\par\vspace{6pt}{\popsemi\fontsize{9}{11.5}\selectfont\color{popmuted}#1}}

% ============================================================
% Signature OnBuch+
% ============================================================
\newcommand{\poplogoinline}[1]{{\poplogo\fontsize{#1}{#1}\selectfont\color{popink}OnBuch\color{poporange}+}}
\newcommand{\signatureonbuch}{%
  \begin{tcolorbox}[enhanced,colback=popink,colframe=popink,boxrule=2.2pt,arc=10pt,
      drop shadow=poporange,left=14pt,right=14pt,top=10pt,bottom=10pt,before skip=0pt,after skip=0pt]
    \tikz[baseline=-0.7ex]\node[circle,fill=popgreen,draw=popcream,line width=1.3pt,minimum size=22pt,inner sep=0]{\popcheck{white}};%
    \hspace{9pt}\begin{minipage}[c]{0.62\linewidth}
      {\popxbold\fontsize{12}{13}\selectfont\color{popcream}Signé\ \textbullet\ {\poplogo OnBuch\color{poporange}+}}\par\vspace{2pt}
      {\raggedright\popsemi\fontsize{8}{10.5}\selectfont\color{popline}\MakeUppercase{Équipe pédagogique OnBuch+}\\\MakeUppercase{Conforme au programme officiel MINESEC}\par}
    \end{minipage}\hfill
    {\poplogo\fontsize{22}{22}\selectfont\color{popcream}OnBuch\color{poporange}+}
  \end{tcolorbox}%
}

% ============================================================
% Page de garde
% #1 titre · #2 matière · #3 niveau · #4 module · #5 n° leçon · #6 durée ·
% #7 description / objectifs (texte ou itemize)
% ============================================================
\newcommand{\pagegarde}[7]{%
  \thispagestyle{empty}
  \newgeometry{a4paper,margin=1.7cm,top=1.6cm,bottom=1.4cm}%
  \begin{tikzpicture}[remember picture,overlay]
    \fill[poporange!18] ([xshift=-3.2cm,yshift=-3.6cm]current page.north east) circle (4.6cm);
    \fill[poppurple!14] ([xshift=2.4cm,yshift=7.6cm]current page.south west) circle (3.8cm);
  \end{tikzpicture}%
  {\poplogo\fontsize{30}{30}\selectfont\color{popink}OnBuch\color{poporange}+}\hfill
  \raisebox{0.6ex}{\poppill{Cours signé \textbullet\ Premium}{popink}{popcream}}\par
  \vspace{1pt}\popsquiggle{poppurple}\par
  \vspace{8pt}
  \begin{tcolorbox}[enhanced,colback=poporange,colframe=popink,boxrule=2.8pt,arc=13pt,
      drop shadow=popink,left=18pt,right=18pt,top=12pt,bottom=13pt,after skip=10pt]
    \poppill{#2 \textbullet\ #3}{popink}{popcream}\hspace{5pt}\poppill{#4}{white}{popink}\par\vspace{8pt}
    {\popdisplay\fontsize{14}{14}\selectfont\color{popink}LEÇON #5}\par\vspace{4pt}
    {\raggedright\hyphenpenalty=10000\exhyphenpenalty=10000\popdisplay\fontsize{25}{27}\selectfont\color{popcream}\MakeUppercase{#1}\par}
    \vspace{8pt}
    {\popxbold\fontsize{9.5}{9.5}\selectfont\color{popink}\MakeUppercase{Durée conseillée : #6}}
  \end{tcolorbox}
  \begin{tcolorbox}[enhanced,colback=white,colframe=popink,boxrule=2.2pt,arc=10pt,
      drop shadow=popink,left=16pt,right=16pt,top=10pt,bottom=10pt,after skip=8pt]
    \poppill{Dans cette leçon}{poppurple}{white}\par\vspace{5pt}
    \popsemi\fontsize{9.3}{12.6}\selectfont\color{popink}#7
  \end{tcolorbox}
  \noindent\begin{minipage}[t]{0.487\linewidth}
    \begin{tcolorbox}[enhanced,colback=popgreen,colframe=popink,boxrule=2.2pt,arc=10pt,
        drop shadow=popink,left=11pt,right=11pt,top=10pt,bottom=10pt,height=3.1cm,valign=center]
      \tcbox[colback=white,colframe=popink,boxrule=1.3pt,arc=5pt,boxsep=0pt,nobeforeafter,
        left=3pt,right=3pt,top=3pt,bottom=3pt,valign=center]{\qrcode[height=1.9cm]{https://play.google.com/store/apps/details?id=com.luvvix.onbuch&hl=fr}}%
      \hspace{8pt}\begin{minipage}[c]{0.5\linewidth}
        {\popdisplay\fontsize{11}{12.5}\selectfont\color{white}\MakeUppercase{Télécharge OnBuch}}\par\vspace{4pt}
        {\raggedright\popsemi\fontsize{8.4}{11}\selectfont\color{white}Gratuit \textbullet\ Google Play\\Tuteur IA, annales, concours, résultats\par}
      \end{minipage}
    \end{tcolorbox}
  \end{minipage}\hfill
  \begin{minipage}[t]{0.487\linewidth}
    \begin{tcolorbox}[enhanced,colback=poppurple,colframe=popink,boxrule=2.2pt,arc=10pt,
        drop shadow=popink,left=11pt,right=11pt,top=10pt,bottom=10pt,height=3.1cm,valign=center]
      \tikz[baseline=-0.7ex]\node[circle,fill=white,draw=popink,line width=1.3pt,minimum size=1.6cm,inner sep=0]{\popdisplay\fontsize{24}{24}\selectfont\color{poppurple}+};%
      \hspace{8pt}\begin{minipage}[c]{0.6\linewidth}
        {\popdisplay\fontsize{11}{12.5}\selectfont\color{white}\MakeUppercase{Passe à OnBuch+}}\par\vspace{4pt}
        {\raggedright\popsemi\fontsize{8.4}{11}\selectfont\color{white}Tous les cours signés, Tuteur IA illimité, fiches PDF — onglet OnBuch+ de l'app\par}
      \end{minipage}
    \end{tcolorbox}
  \end{minipage}
  \vspace{8pt}
  \popcontenu
  \vfill
  \signatureonbuch
  \restoregeometry
  \newpage
}

% Rangée « Ce cours contient » (page de garde)
\newcommand{\poptile}[3]{% #1 couleur #2 gros texte #3 légende
  \begin{tcolorbox}[enhanced,colback=#1,colframe=popink,boxrule=2pt,arc=9pt,shadow={2.5pt}{-2.5pt}{0pt}{popink},
      left=6pt,right=6pt,top=9pt,bottom=9pt,halign=center,valign=center,height=1.75cm,width=\linewidth,nobeforeafter]
    {\popdisplay\fontsize{12.5}{13}\selectfont\color{white}\MakeUppercase{#2}}\par\vspace{3pt}
    {\centering\popsemi\fontsize{7.8}{9.5}\selectfont\color{white}#3\par}
  \end{tcolorbox}}
\newcommand{\popcontenu}{%
  \par\noindent{\popxboldls\fontsize{7.5}{7.5}\selectfont\color{popmuted}CE COURS SIGNÉ CONTIENT}\par\vspace{7pt}
  \noindent\begin{minipage}[t]{0.235\linewidth}\poptile{popblue}{Cours}{complet, illustré, pas à pas}\end{minipage}\hfill
  \begin{minipage}[t]{0.235\linewidth}\poptile{poporange}{Méthodes}{et exercices résolus}\end{minipage}\hfill
  \begin{minipage}[t]{0.235\linewidth}\poptile{poppink}{Exercices}{type examen + corrigés}\end{minipage}\hfill
  \begin{minipage}[t]{0.235\linewidth}\poptile{popgreen}{Fiche bilan}{l'essentiel à retenir}\end{minipage}\par}

% Sommaire
\newtcolorbox{sommaire}{enhanced,colback=white,colframe=popink,boxrule=2.2pt,arc=10pt,
  drop shadow=popink,left=18pt,right=18pt,top=13pt,bottom=13pt,before skip=6pt,after skip=20pt,
  before upper={\poppill{Sommaire}{poppurple}{white}\par\vspace{9pt}\popsemi\fontsize{10.6}{14.5}\selectfont\color{popink}}}

% Signature de fin de cours — poussée tout en bas de la dernière page
\newcommand{\findecours}{%
  \par\vfill\nopagebreak
  \begin{center}\popsquiggle{poporange}\end{center}\vspace{4pt}
  \signatureonbuch
}
\AtBeginDocument{\def\llbracket{\mathopen{[\mkern-3.5mu[}}\def\rrbracket{\mathclose{]\mkern-3.5mu]}}} % intervalles d'entiers (glyphe ⟦ absent de la police)
% Glyphes absents de Fira Math : redéfinis à partir de glyphes disponibles
\AtBeginDocument{%
  \def\setminus{\mathbin{\backslash}}\let\smallsetminus\setminus
  \def\vdots{\mathord{\vbox{\baselineskip3.2pt\lineskiplimit0pt\kern4pt\hbox{.}\hbox{.}\hbox{.}}}}%
  \def\ddots{\mathinner{\mkern1mu\raise6.5pt\hbox{.}\mkern2mu\raise3.6pt\hbox{.}\mkern2mu\raise0.7pt\hbox{.}\mkern1mu}}%
  \def\triangle{\mathord{\scalebox{0.9}{$\symup{\Delta}$}}}%
  \def\nabla{\mathord{\raisebox{\depth}{\scalebox{1}[-1]{$\symup{\Delta}$}}}}%
  \def\checkmark{\popcheck{popdark}}%
  \def\star{\mathbin{\ast}}\def\bigstar{\mathord{\ast}}%
  \def\diamond{\mathbin{\scalebox{0.6}{$\Diamond$}}}%
  \def\bigcup{\mathop{\mathchoice{\raisebox{-0.3em}{\scalebox{1.7}{$\cup$}}}{\scalebox{1.25}{$\cup$}}{\cup}{\cup}}\displaylimits}%
  \def\bigcap{\mathop{\mathchoice{\raisebox{-0.3em}{\scalebox{1.7}{$\cap$}}}{\scalebox{1.25}{$\cap$}}{\cap}{\cap}}\displaylimits}%
  \def\lesseqqgtr{\mathrel{\lessgtr}}\def\bigoplus{\mathop{\oplus}\displaylimits}%
}
\providecommand{\ssi}{si et seulement si} % abréviation fréquente
\AtBeginDocument{\providecommand{\wideparen}{\overparen}} % arc de cercle

%% FIGFIX-BEGIN v4 — correctifs globaux des figures (docs/onbuch-generation/tools/figfix)
\usepackage{adjustbox}\usepackage{etoolbox}
% toute figure TikZ/pgfplots plus large que la ligne est réduite à la largeur disponible
\BeforeBeginEnvironment{tikzpicture}{\begin{adjustbox}{max width=\linewidth}}
\AfterEndEnvironment{tikzpicture}{\end{adjustbox}}
% légendes pgfplots lisibles par-dessus les courbes
\pgfplotsset{every axis/.append style={legend style={fill=white,fill opacity=0.92,text opacity=1,draw=black!15,inner sep=3pt}}}
% étiquettes d'arêtes (syntaxe edge["..."]) avec fond blanc
\tikzset{every edge quotes/.append style={fill=white,inner sep=1.2pt}}
%% FIGFIX-END
\begin{document}

\pagegarde{Séries statistiques regroupées en classes de même amplitude}{Mathématiques}{1re Tech}{Mathématiques – classe de Première technique}{N05}{3 séances (fiche de progression harmonisée)}{Cette leçon étend l'étude des séries statistiques au cas où les données sont regroupées en classes de même amplitude. L'élève apprend à construire les tableaux d'effectifs et de fréquences (cumulés ou non), à calculer les caractéristiques de position et de dispersion, et à représenter la série par un histogramme.
\vspace{4pt}
\begin{multicols}{2}\small
\begin{itemize}
\item À la fin de cette leçon, je sais regrouper une série statistique en classes de même amplitude et construire le tableau des effectifs.
\item À la fin de cette leçon, je sais calculer les effectifs cumulés croissants et décroissants d'une série regroupée en classes.
\item À la fin de cette leçon, je sais calculer les fréquences et les fréquences cumulées et les exprimer en pourcentages.
\item À la fin de cette leçon, je sais déterminer la classe modale, la moyenne (par les centres de classes) et la classe médiane d'une série regroupée.
\item À la fin de cette leçon, je sais déterminer la médiane et les quartiles par interpolation linéaire ou par lecture graphique.
\item À la fin de cette leçon, je sais calculer la variance et l'écart-type d'une série regroupée en classes.
\end{itemize}\end{multicols}}

\begin{sommaire}
\begin{multicols}{2}
\begin{enumerate}
\item Regroupement en classes et effectifs
\item Effectifs cumulés et fréquences
\item Caractéristiques de position
\item Caractéristiques de dispersion
\item Histogramme et polygone
\item Synthèse et applications à la vie courante
\item Activité d'intégration
\item Méthodes et exercices résolus
\item Exercices
\item Corrigés des exercices
\item Fiche bilan
\end{enumerate}
\end{multicols}
\end{sommaire}

\begin{objectifs}
\begin{itemize}
\item À la fin de cette leçon, je sais regrouper une série statistique en classes de même amplitude et construire le tableau des effectifs.
\item À la fin de cette leçon, je sais calculer les effectifs cumulés croissants et décroissants d'une série regroupée en classes.
\item À la fin de cette leçon, je sais calculer les fréquences et les fréquences cumulées et les exprimer en pourcentages.
\item À la fin de cette leçon, je sais déterminer la classe modale, la moyenne (par les centres de classes) et la classe médiane d'une série regroupée.
\item À la fin de cette leçon, je sais déterminer la médiane et les quartiles par interpolation linéaire ou par lecture graphique.
\item À la fin de cette leçon, je sais calculer la variance et l'écart-type d'une série regroupée en classes.
\item À la fin de cette leçon, je sais construire et interpréter un histogramme d'une série regroupée en classes de même amplitude.
\item À la fin de cette leçon, je sais rédiger et justifier une démarche statistique complète appliquée à une situation concrète de la vie courante.
\end{itemize}
\end{objectifs}

\begin{prerequis}
\begin{itemize}
\item Séries statistiques à caractère quantitatif discret (effectifs, fréquences, moyenne) vues en classe de Seconde
\item Effectifs cumulés et fréquences cumulées d'une série discrète
\item Calculs de pourcentages et proportionnalité
\item Lecture et construction de diagrammes en bâtons
\item Notion de variance et d'écart-type pour une série discrète
\end{itemize}
\end{prerequis}

\newpage

\coursec{Regroupement en classes et effectifs}

\textbf{Situation d'accroche.} La SONARA de Limbé, la raffinerie nationale, veut analyser les salaires mensuels de ses $120$ employés, exprimés en milliers de FCFA. Le directeur du personnel dispose d'une longue liste de valeurs : $132$ ; $187$ ; $215$ ; $148$ ; $263$ ; $171$ ; \dots{} Les valeurs sont si nombreuses et si dispersées qu'un tableau classique « une ligne par valeur » serait illisible : presque chaque salaire est différent ! Il décide alors de \emph{regrouper} les salaires en \cle{classes} de même amplitude : $[100\,;\,150[$, $[150\,;\,200[$, $[200\,;\,250[$, etc. Comment construire ces classes ? Comment compter les employés de chaque classe sans erreur ? Et comment, ensuite, calculer une moyenne ou tracer un histogramme à partir de données regroupées ? C'est tout l'objet de cette leçon, dont nous posons ici les fondations : le regroupement en classes et le calcul des effectifs.

\courssub{Rappels de vocabulaire statistique}

Avant de regrouper des données, remettons en place le vocabulaire de base, déjà rencontré dans les classes précédentes.

\begin{definition}[Population, individu, caractère]
En statistique :
\begin{itemize}
\item la \cle{population} est l'ensemble sur lequel porte l'étude (exemple : les $120$ employés de la SONARA) ;
\item un \cle{individu} est un élément de cette population (exemple : un employé) ;
\item le \cle{caractère} est la propriété que l'on observe sur chaque individu (exemple : le salaire mensuel).
\end{itemize}
Le caractère est dit \cle{quantitatif} lorsqu'il s'exprime par un nombre. Il est :
\begin{itemize}
\item \cle{discret} s'il ne prend que des valeurs isolées (nombre d'enfants, nombre de pannes d'une machine par mois) ;
\item \cle{continu} s'il peut prendre n'importe quelle valeur dans un intervalle (taille, masse, salaire, durée, température).
\end{itemize}
\end{definition}

\begin{exemplebox}[Deux situations d'atelier]
\begin{itemize}
\item Un chef d'atelier compte le nombre de pièces défectueuses produites chaque jour : le caractère est quantitatif \textbf{discret} (valeurs $0, 1, 2, 3, \dots$).
\item Un contrôleur qualité mesure le diamètre, en millimètres, d'arbres de transmission : le caractère est quantitatif \textbf{continu} ($24{,}98$ ; $25{,}03$ ; $25{,}00$ ; \dots).
\end{itemize}
\end{exemplebox}

\courssub{Pourquoi regrouper en classes ?}

Imaginons que l'on relève les tailles de $50$ élèves d'un lycée technique de Douala. On obtient $50$ valeurs presque toutes différentes : $162{,}4$ ; $158{,}1$ ; $170{,}2$ ; \dots{} Un tableau « une ligne par valeur » comporterait près de $50$ lignes, chacune d'effectif $1$ ou $2$ : aucune lecture globale n'est possible.

\begin{aretenir}[Principe du regroupement en classes]
Lorsque le caractère est \textbf{continu}, ou \textbf{discret avec un très grand nombre de valeurs}, on regroupe les valeurs dans des intervalles appelés \cle{classes}. On perd un peu de précision (on ne connaît plus chaque valeur exacte), mais on gagne une vision claire de la \emph{répartition} de la population.
\end{aretenir}

C'est exactement la démarche du directeur de la SONARA : plutôt que $120$ salaires isolés, il préfère $5$ ou $6$ classes bien choisies, dont il pourra compter les effectifs puis tirer des conclusions.

\courssub{Classe, amplitude, centre, effectif}

\begin{definition}[Classe]
Une \cle{classe} est un intervalle $[a\,;\,b[$ dans lequel on regroupe les valeurs du caractère.
\begin{itemize}
\item $a$ est la \cle{borne inférieure} de la classe ;
\item $b$ est la \cle{borne supérieure} de la classe ;
\item \textbf{convention d'appartenance} : une valeur $x$ appartient à la classe $[a\,;\,b[$ si et seulement si $a \leq x < b$. La borne inférieure est \textbf{incluse}, la borne supérieure est \textbf{exclue}.
\end{itemize}
\end{definition}

\begin{attention}[Le piège de la borne : ne jamais compter deux fois]
L'erreur la plus fréquente au Probatoire est de compter une valeur située exactement sur une borne dans \emph{deux} classes à la fois. Avec la convention $[a\,;\,b[$, ce risque disparaît : la valeur $x = 150$ appartient à $[150\,;\,200[$ et \textbf{pas} à $[100\,;\,150[$, car $150$ est exclu de la première classe et inclus dans la seconde. Retenez : \textbf{borne de gauche incluse, borne de droite exclue}. Pour la \emph{dernière} classe, on ferme parfois l'intervalle des deux côtés, par exemple $[15\,;\,20]$, afin d'inclure la plus grande valeur de la série.
\end{attention}

\begin{definition}[Amplitude et centre d'une classe]
Soit une classe $[a\,;\,b[$.
\begin{itemize}
\item Son \cle{amplitude} est le nombre $b - a$.
\item Son \cle{centre} (ou milieu) est le nombre
\[ c_i = \frac{a + b}{2}. \]
\end{itemize}
Des classes sont dites de \cle{même amplitude} lorsque toutes les amplitudes $b - a$ sont égales.
\end{definition}

Le centre de classe jouera un rôle capital dans la suite de la leçon : une fois les données regroupées, on ne connaît plus les valeurs exactes, et l'on fera \emph{comme si} tous les individus d'une classe avaient pour valeur le centre de la classe. C'est cette convention qui permettra de calculer une moyenne approchée.

\begin{exemplebox}[Classes de la SONARA]
Le directeur du personnel choisit les classes $[100\,;\,150[$, $[150\,;\,200[$, $[200\,;\,250[$, $[250\,;\,300[$ (en milliers de FCFA).
\begin{itemize}
\item Amplitude de chaque classe : $150 - 100 = 50$. Les classes sont de même amplitude.
\item Centre de la première classe : $c_1 = \dfrac{100 + 150}{2} = 125$.
\item Un employé gagnant exactement $200$ milliers de FCFA est rangé dans $[200\,;\,250[$, pas dans $[150\,;\,200[$.
\end{itemize}
\end{exemplebox}

\begin{popfigure}
\begin{center}
\begin{tikzpicture}[scale=0.9]
% Axe
\draw[->, thick, popdark] (-0.5,0) -- (11,0);
% Graduations
\foreach \x/\lab in {0/0, 2.5/5, 5/10, 7.5/15, 10/20}{
  \draw[thick, popdark] (\x,0.12) -- (\x,-0.12) node[below=2pt, popdark] {\lab};
}
% Classe [10 ; 15[
\draw[line width=2.5pt, popblue] (5,0.55) -- (7.5,0.55);
\fill[popblue] (5,0.55) circle (2.6pt);
\draw[fill=white, draw=popblue, line width=1.4pt] (7.5,0.55) circle (2.6pt);
\node[above=4pt, popblue] at (5,0.55) {\footnotesize borne incluse};
\node[above=4pt, popblue] at (7.5,0.55) {\footnotesize borne exclue};
% Amplitude
\draw[<->, poporange, thick] (5,-0.85) -- (7.5,-0.85) node[midway, below, poporange] {\footnotesize amplitude $= 15 - 10 = 5$};
% Centre
\fill[poppurple] (6.25,0) circle (2.4pt);
\node[above=10pt, poppurple] at (6.25,0.05) {\footnotesize centre $c = \dfrac{10+15}{2} = 12{,}5$};
\draw[dashed, poppurple] (6.25,0) -- (6.25,0.55);
\end{tikzpicture}
\end{center}
\legende{La classe $[10\,;\,15[$ sur une droite graduée : point plein pour la borne incluse, point creux pour la borne exclue, amplitude $5$ et centre $12{,}5$.}
\end{popfigure}

\begin{definition}[Effectif d'une classe, effectif total]
L'\cle{effectif} d'une classe, noté $n_i$, est le nombre d'individus dont la valeur du caractère appartient à cette classe. L'\cle{effectif total}, noté $N$, est la somme des effectifs de toutes les classes :
\[ N = n_1 + n_2 + \dots + n_k = \sum_{i=1}^{k} n_i, \]
où $k$ est le nombre de classes.
\end{definition}

\begin{attention}[Vérification obligatoire]
Après tout dépouillement, vérifiez que la somme des effectifs $n_i$ est égale à l'effectif total $N$ de la population. Si la somme est inférieure, des individus ont été oubliés ; si elle est supérieure, des valeurs ont été comptées deux fois (souvent à cause d'une mauvaise gestion des bornes). Cette vérification, rapide, vous évite de fausser tous les calculs suivants (fréquences, moyenne, histogramme).
\end{attention}

\courssub{Méthode de dépouillement : du tableau brut au tableau de classes}

\begin{methode}[Regrouper une série brute en classes de même amplitude]
\begin{enumerate}
\item \textbf{Repérer les extrêmes} : noter la plus petite valeur $x_{\min}$ et la plus grande valeur $x_{\max}$ de la série, puis calculer l'étendue $e = x_{\max} - x_{\min}$.
\item \textbf{Choisir le nombre de classes} $k$ (souvent entre $4$ et $8$, ou donné par l'énoncé).
\item \textbf{Calculer l'amplitude commune} : $a \approx \dfrac{e}{k}$, arrondie à une valeur simple (par excès) pour que toutes les données soient couvertes.
\item \textbf{Écrire les classes} $[a_1\,;\,a_2[$, $[a_2\,;\,a_3[$, \dots{} en respectant la convention « borne inférieure incluse, borne supérieure exclue » ; fermer éventuellement la dernière classe.
\item \textbf{Pointer} : parcourir la liste brute \emph{une seule fois}, et pour chaque valeur, tracer un trait de pointage dans la ligne de sa classe (traits groupés par paquets de $5$ pour faciliter le comptage).
\item \textbf{Compter} les traits pour obtenir les effectifs $n_i$, puis \textbf{vérifier} que $\sum n_i = N$.
\end{enumerate}
\end{methode}

\begin{exoresolu}[Dépouillement des notes de 40 élèves]
\textbf{Énoncé.} Voici les notes sur $20$ obtenues par $40$ élèves de Première technique à un devoir de mathématiques :
\begin{center}
\small
8 ; 12 ; 5 ; 17 ; 10 ; 14 ; 3 ; 9 ; 11 ; 15 ; 7 ; 13 ; 10 ; 18 ; 6 ; 12 ; 9 ; 14 ; 11 ; 8 ;\\
16 ; 4 ; 12 ; 10 ; 13 ; 7 ; 15 ; 9 ; 11 ; 19 ; 5 ; 12 ; 14 ; 8 ; 10 ; 16 ; 13 ; 6 ; 11 ; 9.
\end{center}
Regrouper ces notes en classes $[0\,;\,5[$, $[5\,;\,10[$, $[10\,;\,15[$, $[15\,;\,20]$, dresser le tableau de dépouillement avec pointage, puis donner les effectifs et vérifier l'effectif total.
\tcblower
\textbf{Solution détaillée.}
\begin{enumerate}
\item Les notes vont de $x_{\min} = 3$ à $x_{\max} = 19$ : les quatre classes proposées, d'amplitude commune $5$, couvrent bien toutes les valeurs.
\item On parcourt la liste et on pointe chaque note dans sa classe. Par exemple, la note $10$ va dans $[10\,;\,15[$ (borne inférieure incluse), la note $15$ va dans $[15\,;\,20]$ (car $15$ est exclu de $[10\,;\,15[$).
\end{enumerate}
\begin{center}
\begin{tabularx}{\linewidth}{|c|X|c|c|}
\hline
\rowcolor{popblueL} \textbf{Classe} & \textbf{Pointage} & \textbf{Effectif } $n_i$ & \textbf{Centre } $c_i$ \\
\hline
$[0\,;\,5[$ & |||| & $4$ & $2{,}5$ \\
\hline
$[5\,;\,10[$ & |||| |||| \;|| & $12$ & $7{,}5$ \\
\hline
$[10\,;\,15[$ & |||| |||| |||| \;| & $16$ & $12{,}5$ \\
\hline
$[15\,;\,20]$ & |||| \;||| & $8$ & $17{,}5$ \\
\hline
\rowcolor{popcream} \textbf{Total} & & $\mathbf{40}$ & \\
\hline
\end{tabularx}
\end{center}
Vérification : $n_1 + n_2 + n_3 + n_4 = 4 + 12 + 16 + 8 = 40 = N$. Le dépouillement est correct. On lit immédiatement que la classe $[10\,;\,15[$ est la plus chargée : c'est la \emph{classe modale}, que nous étudierons dans la section sur les caractéristiques de position.
\end{exoresolu}

\begin{exemplebox}[Tailles de 50 élèves d'un lycée de Douala]
On mesure les tailles, en centimètres, de $50$ élèves d'un lycée technique de Douala. La plus petite taille est $152$ cm, la plus grande $191$ cm. L'étendue vaut $191 - 152 = 39$ cm. On choisit $k = 4$ classes d'amplitude $10$ cm : $[150\,;\,160[$, $[160\,;\,170[$, $[170\,;\,180[$, $[180\,;\,190]$. Après pointage, on obtient :
\begin{center}
\begin{tabularx}{\linewidth}{|c|X|c|c|}
\hline
\rowcolor{popgreenL} \textbf{Classe (cm)} & \textbf{Pointage} & \textbf{Effectif } $n_i$ & \textbf{Centre } $c_i$ \\
\hline
$[150\,;\,160[$ & |||| \;|| & $7$ & $155$ \\
\hline
$[160\,;\,170[$ & |||| |||| |||| & $15$ & $165$ \\
\hline
$[170\,;\,180[$ & |||| |||| |||| \;|| & $17$ & $175$ \\
\hline
$[180\,;\,190]$ & |||| |||| \;| & $11$ & $185$ \\
\hline
\rowcolor{popcream} \textbf{Total} & & $\mathbf{50}$ & \\
\hline
\end{tabularx}
\end{center}
Vérification : $7 + 15 + 17 + 11 = 50 = N$. Remarquez qu'un élève mesurant exactement $170$ cm est rangé dans $[170\,;\,180[$, conformément à la convention.
\end{exemplebox}

\begin{savaistu}[La règle de Sturges : combien de classes choisir ?]
Comment les statisticiens professionnels décident-ils du nombre de classes ? Une règle célèbre, proposée en 1926 par le statisticien allemand Herbert Sturges, donne une estimation du nombre de classes en fonction de l'effectif total $N$ :
\[ k \approx 1 + 3{,}3 \log N, \]
où $\log$ désigne le logarithme décimal. Par exemple, pour les $120$ employés de la SONARA : $k \approx 1 + 3{,}3 \times \log 120 \approx 1 + 3{,}3 \times 2{,}08 \approx 7{,}9$, soit environ $8$ classes. Trop peu de classes écrasent l'information ; trop de classes rendent le tableau aussi illisible que la liste brute. Cette règle n'est pas exigible au Probatoire, mais elle explique pourquoi les énoncés proposent presque toujours entre $4$ et $8$ classes.
\end{savaistu}

\begin{exercice}[À vous de jouer]
Un garage de Yaoundé relève le kilométrage de $30$ véhicules venus en révision. Les valeurs, en milliers de km, vont de $12$ à $108$.
\begin{enumerate}
\item Calculer l'étendue de la série.
\item Le garagiste souhaite $4$ classes de même amplitude. Proposer une amplitude convenable et écrire les quatre classes.
\item Indiquer dans quelle classe on range un véhicule affichant exactement $60$ milliers de km, en justifiant par la convention d'appartenance.
\end{enumerate}
\end{exercice}

\begin{corrige}[Exercice 1]
\begin{enumerate}
\item Étendue : $e = 108 - 12 = 96$ milliers de km.
\item Amplitude : $\dfrac{96}{4} = 24$. On peut arrondir à $25$ pour simplifier : classes $[10\,;\,35[$, $[35\,;\,60[$, $[60\,;\,85[$, $[85\,;\,110]$, qui couvrent bien $[12\,;\,108]$.
\item La valeur $60$ est la borne supérieure de $[35\,;\,60[$ (exclue) et la borne inférieure de $[60\,;\,85[$ (incluse) : le véhicule est rangé dans $[60\,;\,85[$.
\end{enumerate}
\end{corrige}

\begin{aretenir}[L'essentiel de la section]
\begin{itemize}
\item On regroupe en \cle{classes} lorsque le caractère est continu ou prend trop de valeurs.
\item Une classe $[a\,;\,b[$ contient les valeurs $x$ telles que $a \leq x < b$ : borne inférieure incluse, borne supérieure exclue.
\item Amplitude : $b - a$ ; centre : $c_i = \dfrac{a+b}{2}$ ; effectif $n_i$ : nombre d'individus de la classe.
\item Le dépouillement se fait par \cle{pointage}, et l'on vérifie toujours que $\sum n_i = N$.
\end{itemize}
Dans la section suivante, nous enrichirons ce tableau avec les effectifs cumulés et les fréquences, outils indispensables pour calculer la médiane et comparer des séries entre elles.
\end{aretenir}

\coursec{Effectifs cumulés et fréquences}

Après avoir constitué les classes de même amplitude et compté les effectifs de chaque classe (section précédente), nous disposons d'une première vision de la répartition des salaires à la SONARA. Mais pour répondre à des questions comme « Combien d'employés gagnent \emph{moins de} 250\,000 FCFA ? » ou « Quel \emph{pourcentage} du personnel perçoit un salaire \emph{supérieur ou égal à} 300\,000 FCFA ? », il faut organiser l'information de manière cumulative et relative. C'est l'objet de cette section.

\courssub{Effectifs cumulés croissants et décroissants}

L'effectif simple $n_i$ d'une classe ne dit rien sur les classes voisines. Pour situer une valeur par rapport à l'ensemble de la population, on cumule les effectifs dans les deux sens.

\begin{definition}[Effectif cumulé croissant (ECC)]
L'\textbf{effectif cumulé croissant} d'une classe est le nombre total d'individus dont la valeur de la variable est \textbf{strictement inférieure} à la borne supérieure de cette classe.
On le note souvent $N_{< x_{\text{sup}}}$.
\end{definition}

\begin{definition}[Effectif cumulé décroissant (ECD)]
L'\textbf{effectif cumulé décroissant} d'une classe est le nombre total d'individus dont la valeur de la variable est \textbf{supérieure ou égale} à la borne inférieure de cette classe.
On le note souvent $N_{\ge x_{\text{inf}}}$.
\end{definition}

\begin{attention}[Sens des inégalités : piège classique]
\emph{Croissant} $\rightarrow$ on accumule \textbf{depuis le bas} (valeurs petites) $\rightarrow$ inégalité \textbf{stricte} ($<$) sur la borne \textbf{supérieure}.
\emph{Décroissant} $\rightarrow$ on accumule \textbf{depuis le haut} (valeurs grandes) $\rightarrow$ inégalité \textbf{large} ($\ge$) sur la borne \textbf{inférieure}.
Ne les confondez pas : à l'examen, la rédaction précise « strictement inférieur » ou « supérieur ou égal » fait partie des points notés.
\end{attention}

Reprenons l'exemple des salaires mensuels (en milliers de FCFA) de $N=200$ employés de la SONARA, regroupés en 5 classes d'amplitude $h=50$.

\begin{exemplebox}[Construction des ECC et ECD — Salaires SONARA]
\begin{center}
\begin{tabular}{|c|c|c|c|c|}\hline
\textbf{Classes (kFCFA)} & \textbf{Centres} & \textbf{Effectifs $n_i$} & \textbf{ECC} & \textbf{ECD} \\ \hline
$[100 ; 150[$ & 125 & 20 & 20 & 200 \\ \hline
$[150 ; 200[$ & 175 & 40 & 60 & 180 \\ \hline
$[200 ; 250[$ & 225 & 70 & 130 & 140 \\ \hline
$[250 ; 300[$ & 275 & 50 & 180 & 70 \\ \hline
$[300 ; 350[$ & 325 & 20 & 200 & 20 \\ \hline
\textbf{Total} & & \textbf{200} & & \\ \hline
\end{tabular}
\end{center}
\emph{Calculs :}
\begin{itemize}
\item ECC ligne 1 : $20$ (valeurs $< 150$).
\item ECC ligne 2 : $20+40=60$ (valeurs $< 200$).
\item ECC ligne 3 : $60+70=130$ (valeurs $< 250$).
\item ECC ligne 4 : $130+50=180$ (valeurs $< 300$).
\item ECC ligne 5 : $180+20=200=N$ (valeurs $< 350$, soit toute la population).
\item ECD ligne 1 : $200$ (valeurs $\ge 100$, toute la population).
\item ECD ligne 2 : $200-20=180$ (valeurs $\ge 150$).
\item ECD ligne 3 : $180-40=140$ (valeurs $\ge 200$).
\item ECD ligne 4 : $140-70=70$ (valeurs $\ge 250$).
\item ECD ligne 5 : $70-50=20$ (valeurs $\ge 300$).
\end{itemize}
\end{exemplebox}

\begin{propriete}[Propriétés fondamentales des effectifs cumulés]
Pour une série groupée en $k$ classes d'effectif total $N$ :
\begin{enumerate}
\item Le dernier ECC (celui de la dernière classe) vaut \textbf{$N$}.
\item Le premier ECD (celui de la première classe) vaut \textbf{$N$}.
\item Pour toute classe $i$, $\text{ECC}_i + \text{ECD}_{i+1} = N$ (avec la convention $\text{ECD}_{k+1}=0$).
\item Les ECC forment une suite \textbf{croissante} ; les ECD forment une suite \textbf{décroissante}.
\end{enumerate}
\end{propriete}

\courssub{Fréquences et fréquences cumulées}

Les effectifs bruts dépendent de la taille $N$ de l'échantillon. Pour comparer deux séries de tailles différentes (ex. : salaires SONARA vs salaires SNH) ou pour exprimer une proportion, on utilise les \textbf{fréquences}.

\begin{definition}[Fréquence d'une classe]
La \textbf{fréquence} $f_i$ de la classe $i$ est le quotient de son effectif $n_i$ par l'effectif total $N$ :
\[
f_i = \frac{n_i}{N}
\]
La fréquence est un nombre compris entre $0$ et $1$. Elle s'exprime souvent en \textbf{pourcentage} : $f_i \times 100\,\%$.
\end{definition}

\begin{propriete}[Somme des fréquences]
La somme des fréquences de toutes les classes vaut \textbf{1} (soit $100\,\%$).
\[
\sum_{i=1}^{k} f_i = \frac{\sum n_i}{N} = \frac{N}{N} = 1
\]
\end{propriete}

\begin{experience}[Démonstration guidée : pourquoi la somme vaut 1 ?]
\textbf{Consigne :} Complétez le raisonnement.
\begin{enumerate}
\item Par définition, $f_i = \dfrac{n_i}{N}$.
\item La somme des fréquences est $\sum f_i = \sum \dfrac{n_i}{N}$.
\item Factoriser par $\dfrac{1}{N}$ : $\sum f_i = \dfrac{1}{N} \sum n_i$.
\item Que vaut $\sum n_i$ par définition de l'effectif total ?
\item Conclure sur la valeur de $\sum f_i$.
\end{enumerate}
\tcblower
\textbf{Solution :}
$\sum n_i = N$, donc $\sum f_i = \dfrac{1}{N} \times N = 1$. CQFD.
Cette propriété est un \textbf{test de vérification} immédiat de vos calculs.
\end{experience}

\begin{definition}[Fréquences cumulées croissantes (FCC) et décroissantes (FCD)]
On cumule les fréquences exactement comme les effectifs :
\begin{itemize}
\item $\mathrm{FCC}_i$ = somme des fréquences des classes de rang $\le i$ = $\dfrac{\text{ECC}_i}{N}$.
\item $\mathrm{FCD}_i$ = somme des fréquences des classes de rang $\ge i$ = $\dfrac{\text{ECD}_i}{N}$.
\end{itemize}
La dernière FCC vaut $1$ (ou $100\,\%$) ; la première FCD vaut $1$ (ou $100\,\%$).
\end{definition}

\begin{methode}[Passage fréquence $\leftrightarrow$ pourcentage]
\begin{enumerate}
\item \textbf{Fréquence $\to$ Pourcentage :} multiplier par $100$ (ex. : $0,35 \to 35\,\%$).
\item \textbf{Pourcentage $\to$ Fréquence :} diviser par $100$ (ex. : $25\,\% \to 0,25$).
\item \textbf{Effectif $\to$ Fréquence :} diviser par $N$ (ex. : $70/200 = 0,35$).
\item \textbf{Fréquence $\to$ Effectif :} multiplier par $N$ (ex. : $0,10 \times 200 = 20$).
\end{enumerate}
\end{methode}

\courssub{Tableau complet et contrôles de cohérence}

Un tableau statistique sérieux présente toutes les colonnes utiles. Voici le tableau complet pour l'exemple SONARA, avec les fréquences en pourcentage (arrondies au dixième) et les FCC en pourcentage.

\begin{popfigure}
\begin{center}
\begin{tabular}{|c|c|c|c|c|c|}\hline
\rowcolor{popblue}
\textbf{Classes (kFCFA)} & \textbf{Centres} & \textbf{Effectifs $n_i$} & \textbf{ECC} & \textbf{Fréq. (\%)} & \textbf{FCC (\%)} \\ \hline
$[100 ; 150[$ & 125 & 20 & 20 & 10,0 & 10,0 \\ \hline
$[150 ; 200[$ & 175 & 40 & 60 & 20,0 & 30,0 \\ \hline
\rowcolor{popblueL!30}
$[200 ; 250[$ & 225 & 70 & 130 & 35,0 & 65,0 \\ \hline
$[250 ; 300[$ & 275 & 50 & 180 & 25,0 & 90,0 \\ \hline
$[300 ; 350[$ & 325 & 20 & 200 & 10,0 & 100,0 \\ \hline
\rowcolor{popgoldL}
\textbf{Total} & \textbf{--} & \textbf{200} & \textbf{--} & \textbf{100,0} & \textbf{--} \\ \hline
\end{tabular}
\end{center}
\legende{Tableau statistique complet : salaires SONARA (200 employés). La classe modale $[200;250[$ est surlignée.}
\end{popfigure}

\begin{methode}[Vérifier la cohérence d'un tableau statistique (Check-list)]
Avant de rendre une copie ou un rapport technique, contrôlez systématiquement :
\begin{enumerate}
\item \textbf{Somme des effectifs :} $\sum n_i = N$ (dernier ECC = $N$).
\item \textbf{Somme des fréquences :} $\sum f_i = 1$ (ou $\sum f_i\,\% = 100\,\%$).
\item \textbf{Dernière FCC :} vaut $1$ (ou $100\,\%$).
\item \textbf{Ordre des ECC :} strictement croissants (sauf classes vides).
\item \textbf{Ordre des ECD :} strictement décroissants.
\item \textbf{Cohérence ECC/ECD :} $\text{ECC}_i + \text{ECD}_{i+1} = N$.
\end{enumerate}
Si un seul point échoue, il y a une erreur de saisie ou de calcul.
\end{methode}

\begin{attention}[Le piège des arrondis : $99,9\,\%$ ou $100,1\,\%$]
En arrondissant les fréquences en pourcentage à une décimale (ex. : $33,3\,\%$ pour $1/3$), la somme peut ne pas faire exactement $100\,\%$.
\begin{itemize}
\item \textbf{Règle :} On ajuste \textbf{la dernière fréquence} (ou celle de la classe modale) pour forcer la somme à $100,0\,\%$.
\item \textbf{Exemple :} Fréquences exactes : $33,333...\,\%$, $33,333...\,\%$, $33,333...\,\%$. Arrondis : $33,3\,\%$, $33,3\,\%$, $33,3\,\%$ $\to$ Somme $99,9\,\%$.
\item \textbf{Correction :} On écrit $33,3\,\%$, $33,3\,\%$, $\mathbf{33,4\,\%}$.
\item \textbf{Ne JAMAIS} écrire « Total $99,9\,\%$ » : cela prouve que vous n'avez pas vérifié.
\end{itemize}
\end{attention}

\courssub{Lecture d'informations sur les lignes cumulées}

Les lignes ECC et FCC permettent de répondre directement aux questions de seuils sans refaire de calculs.

\begin{exoresolu}[Lecture directe sur le tableau SONARA]
\textbf{Énoncé :} À partir du tableau complet ci-dessus, répondez aux questions suivantes :
\begin{enumerate}
\item Combien d'employés perçoivent un salaire \textbf{strictement inférieur} à 250\,000 FCFA ?
\item Quel \textbf{pourcentage} d'employés perçoit un salaire \textbf{supérieur ou égal} à 300\,000 FCFA ?
\item On veut accorder une prime aux 30\,\% des employés les \textbf{mieux payés}. Quel est le salaire minimal pour en bénéficier (seuil approximatif) ?
\end{enumerate}

\textbf{Solution détaillée :}
\begin{enumerate}
\item « Strictement inférieur à 250 » $\rightarrow$ borne supérieure de la classe $[200;250[$.
On lit l'\textbf{ECC} de cette classe : \textbf{130 employés}.
(Verification : $20+40+70=130$).
\item « Supérieur ou égal à 300 » $\rightarrow$ borne inférieure de la classe $[300;350[$.
On lit l'\textbf{ECD} de cette classe : $20$ employés.
En pourcentage : $\dfrac{20}{200} \times 100 = \mathbf{10\,\%}$.
(Ou lecture directe : FCD de la classe $[300;350[$ = $100 - 90 = 10\,\%$).
\item « 30\,\% les mieux payés » $\rightarrow$ on cherche le seuil \textbf{en partant du haut}.
30\,\% de 200 = 60 employés.
On regarde les ECD : 
\begin{itemize}
\item Classe $[300;350[$ : ECD = 20 (insuffisant).
\item Classe $[250;300[$ : ECD = 70 (dépasse 60).
\end{itemize}
Le seuil se situe \textbf{dans la classe $[250;300[$}.
Comme on n'a pas la répartition exacte à l'intérieur de la classe, on donne une estimation : le seuil est \textbf{proche de 250\,000 FCFA} (borne inférieure de la classe).
\emph{Note :} En Terminale, on utilisera l'interpolation linéaire pour affiner (médiane, quartiles).
\end{enumerate}
\end{exoresolu}

\begin{savaistu}[Statistiques officielles au Cameroun]
L'INS (Institut National de la Statistique) publie chaque année l'\emph{Enquête Camerounaise Auprès des Ménages (ECAM)}. Les revenus y sont présentés par déciles (10\,\% les plus pauvres, 10\,\% les plus riches) grâce aux fréquences cumulées. Savoir lire une FCC permet de comprendre où se situe un ménage par rapport à la population nationale.
\end{savaistu}

\begin{exercice}[Application : Contrôle qualité atelier]
Dans un atelier de mécanique, on mesure le diamètre (en mm) de 50 axes. Résultats groupés :
\begin{center}
\begin{tabular}{|c|c|c|}\hline
Classes (mm) & Effectifs & ECC \\ \hline
$[19,90 ; 19,95[$ & 5 & 5 \\ \hline
$[19,95 ; 20,00[$ & 15 & 20 \\ \hline
$[20,00 ; 20,05[$ & 20 & 40 \\ \hline
$[20,05 ; 20,10[$ & 8 & 48 \\ \hline
$[20,10 ; 20,15[$ & 2 & 50 \\ \hline
\end{tabular}
\end{center}
\begin{enumerate}
\item Compléter le tableau avec les colonnes : Centres, Fréquences (\%), FCC (\%), ECD.
\item Vérifier la cohérence (somme effectifs, somme fréquences, dernier ECC).
\item Combien d'axes ont un diamètre $\ge 20,00$ mm ?
\item Quel pourcentage d'axes a un diamètre $< 19,95$ mm ?
\item La norme exige un diamètre dans $[20,00 ; 20,05[$. Quel pourcentage est conforme ?
\end{enumerate}
\end{exercice}

\begin{corrige}[Exercice : Contrôle qualité atelier]
\begin{enumerate}
\item Tableau complet :
\begin{center}
\begin{tabular}{|c|c|c|c|c|c|c|}\hline
Classes & Centres & $n_i$ & ECC & Fréquences (\%) & FCC (\%) & ECD \\ \hline
$[19,90 ; 19,95[$ & 19,925 & 5 & 5 & 10,0 & 10,0 & 50 \\ \hline
$[19,95 ; 20,00[$ & 19,975 & 15 & 20 & 30,0 & 40,0 & 45 \\ \hline
$[20,00 ; 20,05[$ & 20,025 & 20 & 40 & 40,0 & 80,0 & 30 \\ \hline
$[20,05 ; 20,10[$ & 20,075 & 8 & 48 & 16,0 & 96,0 & 10 \\ \hline
$[20,10 ; 20,15[$ & 20,125 & 2 & 50 & 4,0 & 100,0 & 2 \\ \hline
Total & & 50 & & 100,0 & & \\ \hline
\end{tabular}
\end{center}
\item Contrôles : $\sum n_i = 50 = N$ (OK). Dernier ECC = 50 (OK). $\sum f_i\,\% = 100\,\%$ (OK). ECC croissants, ECD décroissants (OK).
\item Diamètre $\ge 20,00$ mm $\rightarrow$ ECD de la classe $[20,00 ; 20,05[$ = \textbf{30 axes}.
\item Diamètre $< 19,95$ mm $\rightarrow$ ECC de la classe $[19,90 ; 19,95[$ = 5 axes $\rightarrow$ \textbf{10\,\%}.
\item Conformes (classe $[20,00 ; 20,05[$) : 20 axes $\rightarrow$ \textbf{40\,\%}.
\end{enumerate}
\end{corrige}

\coursec{Caractéristiques de position}

Dans la section précédente, nous avons appris à organiser une série statistique regroupée en classes : effectifs, effectifs cumulés croissants (ECC), fréquences et fréquences cumulées. Ces tableaux décrivent la série, mais ils restent longs à lire. Pour \emph{résumer} une série par quelques nombres significatifs, on utilise des \cle{caractéristiques de position} : elles indiquent autour de quelles valeurs se situent les données. Nous étudions ici la \cle{classe modale}, la \cle{moyenne}, la \cle{médiane} et les \cle{quartiles}, dans le cas d'une série regroupée en classes de même amplitude.

Pour toute cette section, nous utilisons la situation suivante, qui nous servira de fil conducteur.

\begin{exemplebox}[Situation fil conducteur : les salaires à la SONARA]
Une enquête porte sur les salaires mensuels (en milliers de francs CFA) de $N = 200$ employés d'une entreprise de la zone industrielle de Limbe. Les résultats sont regroupés en classes de même amplitude $50$ :

\begin{center}
\begin{tabularx}{\linewidth}{|l|*{5}{>{\centering\arraybackslash}X|}}
\hline
\rowcolor{popblueL} Classe (salaire) & $[100\,;\,150[$ & $[150\,;\,200[$ & $[200\,;\,250[$ & $[250\,;\,300[$ & $[300\,;\,350[$ \\
\hline Effectif $n_i$ & $30$ & $50$ & $70$ & $35$ & $15$ \\
\hline ECC & $30$ & $80$ & $150$ & $185$ & $200$ \\
\hline
\end{tabularx}
\end{center}

On vérifie : $30+50+70+35+15 = 200 = N$.
\end{exemplebox}

\courssub{La classe modale}

\textbf{Intuition.} Si l'on demande : « Quelle tranche de salaire est la plus fréquente dans l'entreprise ? », il suffit de repérer la classe qui contient le plus d'employés. C'est l'idée de la classe modale.

\begin{definition}[Classe modale]
Dans une série statistique regroupée en classes, la \cle{classe modale} est la classe qui a le plus grand effectif (ou, de façon équivalente, la plus grande fréquence).
\end{definition}

\begin{attention}[Classes de même amplitude]
Comparer les effectifs pour trouver la classe modale n'est valable que parce que les classes ont \textbf{toutes la même amplitude}. Si les amplitudes étaient différentes, il faudrait comparer les \emph{densités} (effectif divisé par l'amplitude). Dans cette leçon, toutes les classes ont la même amplitude, donc la comparaison directe des effectifs est correcte.
\end{attention}

Pour la série SONARA, le plus grand effectif est $70$, atteint par la classe $[200\,;\,250[$. La classe modale est donc $[200\,;\,250[$ : c'est dans cette tranche (entre \num{200000} et \num{250000} FCFA) que se situe le plus grand nombre de salariés.

\courssub{Moyenne d'une série regroupée en classes}

\textbf{Intuition.} Pour une série non regroupée, la moyenne est $\bar{x} = \dfrac{1}{N}\sum n_i x_i$. Mais après regroupement en classes, nous ne connaissons plus les valeurs exactes : nous savons seulement que $30$ salaires sont \emph{quelque part} entre $100$ et $250$... On fait alors une hypothèse naturelle.

\begin{definition}[Hypothèse des centres de classes et moyenne]
Pour une série regroupée en classes, on fait l'\cle{hypothèse de répartition uniforme} : on suppose que toutes les valeurs d'une classe sont concentrées au \cle{centre de la classe} $c_i = \dfrac{\text{borne inférieure} + \text{borne supérieure}}{2}$.

La \cle{moyenne} de la série regroupée est alors :
\[ \bar{x} = \frac{1}{N}\sum_{i} n_i\, c_i = \frac{n_1 c_1 + n_2 c_2 + \cdots + n_k c_k}{N} \]
où $n_i$ est l'effectif de la classe numéro $i$, $c_i$ son centre et $N$ l'effectif total.
\end{definition}

\begin{methode}[Calculer la moyenne d'une série regroupée]
\begin{enumerate}
\item Ajouter au tableau une colonne (ou une ligne) « centre $c_i$ » : moyenne des deux bornes de chaque classe.
\item Ajouter une colonne « produit $n_i \times c_i$ ».
\item Additionner tous les produits $n_i c_i$.
\item Diviser cette somme par l'effectif total $N$.
\end{enumerate}
\end{methode}

Appliquons cette méthode à la série SONARA. Le tableau de calcul complet est le suivant :

\begin{center}
\begin{tabularx}{\linewidth}{|l|*{4}{>{\centering\arraybackslash}X|}}
\hline
\rowcolor{popblueL} Classe & Centre $c_i$ & Effectif $n_i$ & $n_i \times c_i$ & ECC \\
\hline $[100\,;\,150[$ & $125$ & $30$ & $\num{3750}$ & $30$ \\
\hline $[150\,;\,200[$ & $175$ & $50$ & $\num{8750}$ & $80$ \\
\hline $[200\,;\,250[$ & $225$ & $70$ & $\num{15750}$ & $150$ \\
\hline $[250\,;\,300[$ & $275$ & $35$ & $\num{9625}$ & $185$ \\
\hline $[300\,;\,350[$ & $325$ & $15$ & $\num{4875}$ & $200$ \\
\hline \textbf{Total} & & $N = 200$ & $\sum n_i c_i = \num{42750}$ & \\
\hline
\end{tabularx}
\end{center}

La moyenne est donc :
\[ \bar{x} = \frac{\num{42750}}{200} = \num{213,75}. \]

\textbf{Interprétation.} Le salaire mensuel moyen des employés enquêtés est d'environ $\num{213750}$ FCFA (puisque les classes sont exprimées en milliers de francs).

\begin{savaistu}[Une moyenne approchée]
La moyenne calculée avec les centres de classes est une \emph{valeur approchée} de la vraie moyenne, car nous avons remplacé chaque salaire réel par le centre de sa classe. Elle est d'autant plus \textbf{précise} que les classes sont \textbf{resserrées} (amplitude petite) : avec des classes de $10$ milliers de francs au lieu de $50$, l'erreur commise serait bien plus faible. C'est pourquoi les instituts de statistique, comme l'INS à Yaoundé, choisissent soigneusement le découpage en classes lors des enquêtes sur les revenus des ménages camerounais.
\end{savaistu}

\courssub{Classe médiane et médiane par interpolation linéaire}

\textbf{Intuition.} La médiane est la valeur qui partage la série en deux groupes de même effectif : $50\,\%$ des salariés gagnent moins que la médiane, $50\,\%$ gagnent plus. Pour une série regroupée, on ne connaît pas les valeurs exactes : on procède en deux temps, d'abord repérer la \emph{classe} médiane, puis préciser la valeur à l'intérieur de cette classe.

\begin{definition}[Classe médiane]
La \cle{classe médiane} est la \textbf{première} classe dont l'effectif cumulé croissant (ECC) atteint ou dépasse $\dfrac{N}{2}$.
\end{definition}

Pour la série SONARA : $\dfrac{N}{2} = \dfrac{200}{2} = 100$. Les ECC sont $30$, $80$, $150$, $185$, $200$. Le premier ECC qui atteint ou dépasse $100$ est $150$, qui correspond à la classe $[200\,;\,250[$. La classe médiane est donc $[200\,;\,250[$.

\begin{attention}[Classe modale $\neq$ classe médiane]
Deux erreurs sont très fréquentes au Probatoire :
\begin{itemize}
\item \textbf{Confondre classe modale et classe médiane.} La classe modale se repère dans la colonne des \emph{effectifs} (le plus grand), la classe médiane dans la colonne des \emph{ECC} (le premier qui atteint $N/2$). Ici, les deux coïncident ($[200\,;\,250[$), mais ce n'est pas toujours le cas !
\item \textbf{Prendre le centre de la classe médiane comme médiane.} Dire « la médiane est $225$ » sans calcul est faux : le centre n'est la médiane que si l'effectif cumulé avant la classe vaut exactement $N/2$. Il faut appliquer l'\emph{interpolation linéaire}.
\end{itemize}
\end{attention}

\begin{propriete}[Médiane par interpolation linéaire (admise)]
Soit $[a\,;\,b[$ la classe médiane, d'effectif $n_{\text{med}}$, et soit $E$ l'effectif cumulé croissant de la classe \textbf{précédente} (c'est-à-dire l'ECC juste avant d'entrer dans la classe médiane). Sous l'hypothèse de \textbf{répartition uniforme} des valeurs dans la classe médiane, la médiane est :
\[ Me = a + (b-a)\times \frac{\dfrac{N}{2} - E}{n_{\text{med}}}. \]
Cette formule est admise ; sa justification intuitive est donnée dans l'activité ci-dessous.
\end{propriete}

\begin{experience}[Pourquoi cette formule ? Justification intuitive]
Imaginons les $n_{\text{med}}$ valeurs de la classe médiane $[a\,;\,b[$ \emph{régulièrement espacées} entre $a$ et $b$ (hypothèse de répartition uniforme).
\begin{enumerate}
\item Avant la classe médiane, $E$ valeurs ont déjà été comptées. Il nous manque donc $\dfrac{N}{2} - E$ valeurs pour atteindre la moitié de l'effectif.
\item La fraction de la classe médiane qu'il faut « parcourir » est donc $\dfrac{\frac{N}{2} - E}{n_{\text{med}}}$.
\item La classe ayant une amplitude $b - a$, la distance à ajouter à la borne $a$ est $(b-a)\times \dfrac{\frac{N}{2} - E}{n_{\text{med}}}$.
\item D'où $Me = a + (b-a)\times \dfrac{\frac{N}{2} - E}{n_{\text{med}}}$ : c'est une \emph{interpolation linéaire}, c'est-à-dire une répartition proportionnelle le long du segment $[a\,;\,b]$.
\end{enumerate}
\end{experience}

\begin{popfigure}
\begin{center}
\begin{tikzpicture}[scale=1]
% Axe horizontal : valeurs de la classe médiane
\draw[->, thick, popdark] (0,0) -- (10.5,0) node[below] {\small Valeurs};
% Axe vertical : ECC
\draw[->, thick, popdark] (0,0) -- (0,4.6) node[above] {\small ECC};
% Graduations x : a=200 -> x=1, b=250 -> x=9
\draw[thick, popdark] (1,0.08) -- (1,-0.08) node[below] {$a=200$};
\draw[thick, popdark] (9,0.08) -- (9,-0.08) node[below] {$b=250$};
% Niveaux ECC : E=80 -> y=1.2 ; N/2=100 -> y=2.1 ; E+n=150 -> y=3.6
\draw[dashed, popmuted] (0,1.2) node[left, popdark] {\small $E=80$} -- (1,1.2);
\draw[dashed, popmuted] (0,2.1) node[left, popdark] {\small $\frac{N}{2}=100$} -- (4.6,2.1);
\draw[dashed, popmuted] (0,3.6) node[left, popdark] {\small $E+n_{\text{med}}=150$} -- (9,3.6);
% Segment d'interpolation
\draw[very thick, popblue] (1,1.2) -- (9,3.6);
% Point médiane
\fill[poporange] (4.6,2.1) circle (2.5pt);
\draw[thick, poporange] (4.6,2.1) -- (4.6,0) node[below, poporange] {$Me$};
% Accolades / annotations
\draw[<->, popgreen, thick] (1,0.55) -- (9,0.55) node[midway, above, popgreen] {\small amplitude $b-a=50$};
\node[popblue, align=center] at (6.3,3.9) {\small répartition uniforme :\\ \small l'ECC croît en ligne droite};
\end{tikzpicture}
\end{center}
\legende{Interpolation linéaire dans la classe médiane : on suppose que l'effectif cumulé croît régulièrement de $E$ à $E+n_{\text{med}}$ le long de la classe $[a\,;\,b[$ ; la médiane est l'abscisse du point d'ordonnée $N/2$.}
\end{popfigure}

\begin{methode}[Déterminer la médiane d'une série regroupée]
\begin{enumerate}
\item Calculer $\dfrac{N}{2}$.
\item Repérer la classe médiane dans la colonne des ECC : c'est la première classe dont l'ECC atteint ou dépasse $\dfrac{N}{2}$.
\item Noter $a$, $b$, $n_{\text{med}}$ et l'ECC $E$ de la classe précédente, puis appliquer la formule d'interpolation linéaire :
\[ Me = a + (b-a)\times \frac{\dfrac{N}{2} - E}{n_{\text{med}}}. \]
\end{enumerate}
\end{methode}

\courssub{Quartiles $Q_1$ et $Q_3$ par interpolation linéaire}

\textbf{Intuition.} Les quartiles affinent le découpage de la série : $Q_1$ laisse $25\,\%$ des valeurs en dessous, $Q_3$ en laisse $75\,\%$. Avec la médiane, ils partagent la série en quatre groupes de même effectif.

\begin{definition}[Quartiles d'une série regroupée]
Le \cle{premier quartile} $Q_1$ est la valeur telle qu'au moins $25\,\%$ des valeurs lui sont inférieures ; le \cle{troisième quartile} $Q_3$ est la valeur telle qu'au moins $75\,\%$ des valeurs lui sont inférieures.

On les détermine comme la médiane, par interpolation linéaire, en remplaçant la position $\dfrac{N}{2}$ par $\dfrac{N}{4}$ pour $Q_1$ et par $\dfrac{3N}{4}$ pour $Q_3$ :
\[ Q_1 = a_1 + (b_1 - a_1)\times\frac{\dfrac{N}{4} - E_1}{n_1}, \qquad Q_3 = a_3 + (b_3 - a_3)\times\frac{\dfrac{3N}{4} - E_3}{n_3}, \]
où $[a_1\,;\,b_1[$ est la première classe dont l'ECC atteint ou dépasse $\dfrac{N}{4}$, $E_1$ l'ECC de la classe précédente et $n_1$ son effectif (notations analogues pour $Q_3$).
\end{definition}

\begin{aretenir}[Les trois positions à retenir]
\begin{center}
\begin{tabularx}{0.9\linewidth}{|l|*{3}{>{\centering\arraybackslash}X|}}
\hline
\rowcolor{popblueL} Caractéristique & $Q_1$ & $Me$ & $Q_3$ \\
\hline Position dans l'effectif & $\dfrac{N}{4}$ & $\dfrac{N}{2}$ & $\dfrac{3N}{4}$ \\
\hline Proportion en dessous & $25\,\%$ & $50\,\%$ & $75\,\%$ \\
\hline
\end{tabularx}
\end{center}
La démarche est toujours la même : position $\to$ classe dans les ECC $\to$ interpolation linéaire.
\end{aretenir}

\courssub{Exemple résolu complet : la série des salaires SONARA}

\begin{exoresolu}[Moyenne, médiane et quartiles des salaires SONARA]
On reprend la série des salaires mensuels (en milliers de FCFA) de $N=200$ employés :
\begin{center}
\begin{tabular}{|l|c|c|c|c|c|}
\hline
Classe & $[100;150[$ & $[150;200[$ & $[200;250[$ & $[250;300[$ & $[300;350[$ \\
\hline Effectif $n_i$ & $30$ & $50$ & $70$ & $35$ & $15$ \\
\hline ECC & $30$ & $80$ & $150$ & $185$ & $200$ \\
\hline
\end{tabular}
\end{center}
\begin{enumerate}
\item Déterminer la classe modale.
\item Calculer le salaire moyen $\bar{x}$.
\item Déterminer la médiane $Me$ par interpolation linéaire.
\item Déterminer $Q_1$ et $Q_3$, puis interpréter les résultats.
\end{enumerate}
\tcblower
\textbf{1. Classe modale.} Le plus grand effectif est $70$ : la classe modale est $[200\,;\,250[$.

\textbf{2. Moyenne.} Avec les centres $125$, $175$, $225$, $275$, $325$ :
\begin{align*}
\bar{x} &= \frac{30\times 125 + 50\times 175 + 70\times 225 + 35\times 275 + 15\times 325}{200} \\
&= \frac{\num{3750} + \num{8750} + \num{15750} + \num{9625} + \num{4875}}{200} = \frac{\num{42750}}{200} = \num{213,75}.
\end{align*}
Le salaire moyen est d'environ $\num{213750}$ FCFA.

\textbf{3. Médiane.} $\dfrac{N}{2} = 100$. Le premier ECC $\geq 100$ est $150$ : la classe médiane est $[200\,;\,250[$, avec $a=200$, $b=250$, $n_{\text{med}} = 70$ et ECC précédent $E = 80$.
\begin{align*}
Me &= 200 + (250-200)\times\frac{100 - 80}{70} = 200 + 50\times\frac{20}{70} \\
&= 200 + \frac{\num{1000}}{70} \approx 200 + \num{14,29} \approx \num{214,3}.
\end{align*}
La médiane est d'environ $\num{214300}$ FCFA : la moitié des employés gagne moins que cette somme, l'autre moitié davantage.

\textbf{4. Quartiles.} Pour $Q_1$ : $\dfrac{N}{4} = 50$. Le premier ECC $\geq 50$ est $80$ : la classe de $Q_1$ est $[150\,;\,200[$, avec $E_1 = 30$ et $n_1 = 50$.
\[ Q_1 = 150 + 50\times\frac{50-30}{50} = 150 + 20 = 170. \]
Pour $Q_3$ : $\dfrac{3N}{4} = 150$. Le premier ECC $\geq 150$ est $150$ : la classe de $Q_3$ est $[200\,;\,250[$, avec $E_3 = 80$ et $n_3 = 70$.
\[ Q_3 = 200 + 50\times\frac{150-80}{70} = 200 + 50\times\frac{70}{70} = 250. \]

\textbf{Interprétation concrète.} Dans cette entreprise de Limbe :
\begin{itemize}
\item $25\,\%$ des employés gagnent moins de $\num{170000}$ FCFA ($Q_1 = 170$) ;
\item la moitié gagne moins d'environ $\num{214300}$ FCFA ($Me \approx \num{214,3}$) ;
\item $75\,\%$ gagnent moins de $\num{250000}$ FCFA ($Q_3 = 250$), autrement dit le quart le mieux payé gagne plus de $\num{250000}$ FCFA.
\end{itemize}
On remarque que la moyenne ($\num{213,75}$) et la médiane ($\num{214,3}$) sont proches : la répartition des salaires est à peu près équilibrée autour de la valeur centrale.
\end{exoresolu}

\begin{attention}[Pièges de rédaction au Probatoire]
\begin{itemize}
\item \textbf{Ne pas oublier l'ECC précédent.} Dans la formule, $E$ est l'effectif cumulé de la classe \emph{avant} la classe médiane, pas l'ECC de la classe médiane elle-même. Ici, $E = 80$ et non $150$.
\item \textbf{Bien choisir la « première » classe.} Pour $Q_3$, l'ECC de la classe $[200\,;\,250[$ vaut exactement $150 = \frac{3N}{4}$ : la règle « atteint ou dépasse » désigne bien cette classe.
\item \textbf{Donner l'unité.} Les classes sont en milliers de FCFA : conclure « $Me = \num{214,3}$ » sans préciser l'unité fait perdre des points. Rédiger : $Me \approx \num{214300}$ FCFA.
\item \textbf{Vérifier la cohérence.} On doit toujours avoir $Q_1 \leq Me \leq Q_3$. Ici : $170 \leq \num{214,3} \leq 250$ : le résultat est cohérent.
\end{itemize}
\end{attention}

\begin{exercice}[À vous de jouer : les tailles des pièces en atelier]
Un atelier de mécanique à Douala contrôle le diamètre (en mm) de $120$ pièces usinées :
\begin{center}
\begin{tabular}{|l|c|c|c|c|}
\hline
Diamètre (mm) & $[49,5;49,7[$ & $[49,7;49,9[$ & $[49,9;50,1[$ & $[50,1;50,3[$ \\
\hline Effectif & $18$ & $30$ & $48$ & $24$ \\
\hline
\end{tabular}
\end{center}
\begin{enumerate}
\item Construire la ligne des ECC et déterminer la classe modale.
\item Calculer le diamètre moyen des pièces.
\item Déterminer la médiane, $Q_1$ et $Q_3$ par interpolation linéaire.
\item Le cahier des charges exige un diamètre compris entre $Q_1$ et $Q_3$ pour au moins la moitié de la production : est-ce le cas ici ?
\end{enumerate}
\end{exercice}

\begin{corrige}[Exercice : les tailles des pièces en atelier]
\textbf{1.} ECC : $18$, $48$, $96$, $120$. Classe modale : $[49,9\,;\,50,1[$ (effectif $48$).

\textbf{2.} Centres : $\num{49,6}$, $\num{49,8}$, $50$, $\num{50,2}$.
\[ \bar{x} = \frac{18\times \num{49,6} + 30\times \num{49,8} + 48\times 50 + 24\times \num{50,2}}{120} = \frac{\num{5991,6}}{120} = \num{49,93}\ \text{mm}. \]

\textbf{3.} $\frac{N}{2} = 60$ : classe médiane $[49,9\,;\,50,1[$, $E = 48$, $n_{\text{med}} = 48$.
\[ Me = \num{49,9} + \num{0,2}\times\frac{60-48}{48} = \num{49,9} + \num{0,05} = \num{49,95}\ \text{mm}. \]
$\frac{N}{4} = 30$ : classe de $Q_1$ : $[49,7\,;\,49,9[$, $E_1 = 18$, $n_1 = 30$.
\[ Q_1 = \num{49,7} + \num{0,2}\times\frac{30-18}{30} = \num{49,7} + \num{0,08} = \num{49,78}\ \text{mm}. \]
$\frac{3N}{4} = 90$ : classe de $Q_3$ : $[49,9\,;\,50,1[$, $E_3 = 48$, $n_3 = 48$.
\[ Q_3 = \num{49,9} + \num{0,2}\times\frac{90-48}{48} = \num{49,9} + \num{0,175} = \num{50,075}\ \text{mm}. \]

\textbf{4.} Par construction des quartiles, environ $50\,\%$ des pièces ont un diamètre compris entre $Q_1$ et $Q_3$, c'est-à-dire entre $\num{49,78}$ et $\num{50,075}$ mm : la condition du cahier des charges est satisfaite (de justesse).
\end{corrige}

\begin{aretenir}[Ce qu'il faut retenir de cette section]
\begin{itemize}
\item La \cle{classe modale} : plus grand effectif ; la \cle{classe médiane} : premier ECC $\geq \frac{N}{2}$. Ce sont deux notions différentes.
\item La moyenne d'une série regroupée se calcule avec les \cle{centres de classes} : $\bar{x} = \dfrac{1}{N}\sum n_i c_i$.
\item Médiane et quartiles se calculent par \cle{interpolation linéaire} sous l'hypothèse de répartition uniforme :
\[ Me = a + (b-a)\times\frac{\frac{N}{2} - E}{n_{\text{med}}}, \]
et de même pour $Q_1$ (position $\frac{N}{4}$) et $Q_3$ (position $\frac{3N}{4}$).
\item Toujours vérifier $Q_1 \leq Me \leq Q_3$ et conclure avec l'unité et une phrase d'interprétation.
\end{itemize}
\end{aretenir}

\coursec{Caractéristiques de dispersion}

Après avoir localisé la série grâce aux caractéristiques de position (moyenne, médiane, quartiles), il est indispensable de mesurer \emph{comment les valeurs s'éparpillent} autour de ce centre. Deux séries de salaires peuvent avoir la même moyenne de \SI{200000}{FCFA} mais des réalités sociales très différentes : l'une où tous les salaires sont proches de \SI{200000}{FCFA}, l'autre où quelques cadres touchent \SI{1000000}{FCFA} tandis que la majorité perçoit le SMIC. Les caractéristiques de dispersion rendent compte de cette hétérogénéité.

\courssub{Étendue et intervalle interquartile}

\textbf{Étendue}\par
Pour une série regroupée en classes, l'étendue totale n'est pas connue exactement (on ne connaît pas les valeurs extrêmes réelles). On utilise une \textbf{valeur approchée} :
\begin{definition}[Étendue approchée d'une série regroupée]
L'étendue approchée est la différence entre la \textbf{borne supérieure de la dernière classe} et la \textbf{borne inférieure de la première classe} :
\[
E \approx b_{\text{max}} - a_{\text{min}}
\]
où $[a_{\text{min}}; b_{\text{max}}[$ est l'intervalle couvrant toutes les classes.
\end{definition}

\begin{attention}[Piège : classes ouvertes]
Si la première classe est de la forme $[a_1; b_1[$ sans borne inférieure connue (ex. « moins de 20 ») ou la dernière classe est $[a_k; +\infty[$, l'étendue \textbf{ne peut pas être calculée}. On se rabat alors sur l'intervalle interquartile.
\end{attention}

\textbf{Quartiles et intervalle interquartile}\par
Les quartiles $Q_1$ et $Q_3$ coupent la série en quatre parties d'effectifs égaux. Pour une série regroupée, on les détermine \textbf{par interpolation linéaire} dans la classe où se trouve le rang recherché (même méthode que pour la médiane vue en section 3).

\begin{propriete}[Intervalle interquartile et écart interquartile]
L'intervalle interquartile est $[Q_1 ; Q_3]$. Il contient les \textbf{50\% centraux} de la population (les valeurs « typiques »).
L'écart interquartile (EIQ) est :
\[
\text{EIQ} = Q_3 - Q_1
\]
C'est un indicateur de dispersion \textbf{robuste} : il n'est pas sensible aux valeurs extrêmes (contrairement à l'étendue).
\end{propriete}

\begin{aretenir}[Choix de l'indicateur]
\begin{itemize}
    \item \cle{Étendue} : simple, mais très sensible aux extrêmes, imprécise pour séries regroupées.
    \item \cle{Écart interquartile (EIQ)} : robuste, utilisable même avec classes ouvertes, privilégié pour comparer des dispersions « centrales ».
    \item \cle{Variance / Écart-type} : utilisent \textbf{toutes} les valeurs (via les centres de classes), base de l'inférence statistique, sensibles aux extrêmes.
\end{itemize}
\end{aretenir}

\courssub{Variance et écart-type d'une série regroupée}

\textbf{Intuition : mesurer l'écart à la moyenne}\par
La moyenne $\bar{x}$ est le centre de gravité. Pour quantifier la dispersion, on regarde la distance de chaque centre de classe $c_i$ à $\bar{x}$. Comme les écarts peuvent être négatifs, on les \textbf{met au carré}. On pondère par les effectifs $n_i$ et on fait la moyenne.

\begin{definition}[Variance et écart-type d'une série regroupée en classes de même amplitude]
Soit une série de $N$ individus regroupée en $k$ classes de même amplitude. On note $c_i$ le centre de la classe $i$ et $n_i$ son effectif ($\sum n_i = N$). La moyenne est $\bar{x} = \frac{1}{N}\sum n_i c_i$.

La \textbf{variance} $V$ (ou $\sigma^2$) est la moyenne des carrés des écarts aux centres de classes :
\[
V = \frac{1}{N}\sum_{i=1}^{k} n_i (c_i - \bar{x})^2
\]

L'\textbf{écart-type} $\sigma$ est la racine carrée de la variance :
\[
\sigma = \sqrt{V}
\]
Il s'exprime dans la \textbf{même unité} que la variable étudiée (FCFA, mm, kg, etc.).
\end{definition}

\begin{savaistu}[Unité de la variance]
Si $X$ est en mètres, $V$ est en $\text{m}^2$ et $\sigma$ en $\text{m}$. L'écart-type est donc directement interprétable : « en moyenne, les valeurs s'écartent de $\sigma$ unités de la moyenne ».
\end{savaistu}

\textbf{Théorème de König-Huygens : une formule de calcul simplifiée}\par
Calculer $(c_i - \bar{x})^2$ pour chaque classe entraîne des arrondis successifs sur $\bar{x}$ et des manipulations fastidieuses. Le théorème suivant permet de calculer $V$ directement à partir des sommes $\sum n_i c_i$ et $\sum n_i c_i^2$.

\begin{experience}[Démonstration guidée du théorème de König-Huygens]
\emph{Objectif : transformer $\sum n_i (c_i - \bar{x})^2$ en expression sans $\bar{x}$ à l'intérieur de la somme.}

\begin{enumerate}
    \item \textbf{Développer le carré} : $(c_i - \bar{x})^2 = c_i^2 - 2c_i\bar{x} + \bar{x}^2$.
    \item \textbf{Multiplier par $n_i$ et sommer} :
    \[
    \sum n_i (c_i - \bar{x})^2 = \sum n_i c_i^2 - 2\bar{x}\sum n_i c_i + \bar{x}^2\sum n_i
    \]
    \item \textbf{Utiliser les définitions} : $\sum n_i c_i = N\bar{x}$ et $\sum n_i = N$.
    \item \textbf{Substituer} :
    \[
    \sum n_i (c_i - \bar{x})^2 = \sum n_i c_i^2 - 2\bar{x}(N\bar{x}) + \bar{x}^2 N = \sum n_i c_i^2 - N\bar{x}^2
    \]
    \item \textbf{Diviser par $N$} pour obtenir la variance :
    \[
    V = \frac{1}{N}\sum n_i c_i^2 - \bar{x}^2
    \]
\end{enumerate}
\end{experience}

\begin{propriete}[Formule de König-Huygens (exigible au Probatoire)]
Pour une série statistique regroupée en classes de centres $c_i$ et d'effectifs $n_i$ (total $N$), de moyenne $\bar{x}$ :
\[
\boxed{V = \frac{1}{N}\sum_{i=1}^{k} n_i c_i^2 - \bar{x}^2}
\]
Cette formule est \textbf{numériquement plus stable} (moins d'erreurs d'arrondi) et plus rapide à l'examen.
\end{propriete}

\begin{methode}[Calcul de l'écart-type par König-Huygens (étapes)]
\begin{enumerate}
    \item Construire le tableau avec les colonnes : $c_i$ (centres), $n_i$, $n_i c_i$, $c_i^2$, $n_i c_i^2$.
    \item Calculer les totaux : $N = \sum n_i$, $S_1 = \sum n_i c_i$, $S_2 = \sum n_i c_i^2$.
    \item Calculer la moyenne : $\bar{x} = S_1 / N$.
    \item Appliquer König-Huygens : $V = \frac{S_2}{N} - \bar{x}^2$.
    \item En déduire l'écart-type : $\sigma = \sqrt{V}$.
    \item \textbf{Interpréter} : « L'écart-type est de \dots ; les valeurs se situent majoritairement dans l'intervalle $[\bar{x}-\sigma ; \bar{x}+\sigma]$. »
\end{enumerate}
\end{methode}

\begin{popfigure}
\centering
\begin{tabularx}{\linewidth}{|c|>{\centering\arraybackslash}X|>{\centering\arraybackslash}X|>{\centering\arraybackslash}X|>{\centering\arraybackslash}X|>{\centering\arraybackslash}X|}
\hline
\rowcolor{popblueL}
\textbf{Classe} & \textbf{Centre $c_i$} & \textbf{Eff. $n_i$} & \textbf{$n_i c_i$} & \textbf{$c_i^2$} & \textbf{$n_i c_i^2$} \\ \hline
$[150; 200[$ & 175 & 12 & 2100 & 30625 & 367500 \\ \hline
$[200; 250[$ & 225 & 28 & 6300 & 50625 & 1417500 \\ \hline
$[250; 300[$ & 275 & 35 & 9625 & 75625 & 2646875 \\ \hline
$[300; 350[$ & 325 & 18 & 5850 & 105625 & 1901250 \\ \hline
$[350; 400[$ & 375 & 7 & 2625 & 140625 & 984375 \\ \hline
\rowcolor{popblueL}
\textbf{Total} & & \textbf{100} & \textbf{26500} & & \textbf{7317500} \\ \hline
\end{tabularx}
\legende{Tableau de calcul de la variance (exemple salaires en milliers de FCFA).}
\end{popfigure}

\begin{exoresolu}[Calcul de $V$ et $\sigma$ pour une série de salaires]
\emph{Énoncé :} Dans une entreprise de 100 employés, les salaires mensuels (en milliers de FCFA) sont regroupés ainsi :
\begin{center}
\begin{tabular}{|c|c|c|c|c|c|}\hline
Classe (kFCFA) & $[150;200[$ & $[200;250[$ & $[250;300[$ & $[300;350[$ & $[350;400[$ \\ \hline
Effectif & 12 & 28 & 35 & 18 & 7 \\ \hline
\end{tabular}
\end{center}
Calculer la variance et l'écart-type. Interpréter.

\tcblower
\emph{Solution détaillée :}
On utilise le tableau ci-dessus (colonnes $c_i, n_i c_i, c_i^2, n_i c_i^2$).
\begin{align*}
N &= 100 \\
S_1 = \sum n_i c_i &= 26500 \quad \Rightarrow \quad \bar{x} = \frac{26500}{100} = 265 \text{ kFCFA} \\
S_2 = \sum n_i c_i^2 &= 7317500
\end{align*}
Variance (König-Huygens) :
\[
V = \frac{S_2}{N} - \bar{x}^2 = \frac{7317500}{100} - 265^2 = 73175 - 70225 = 2950 \text{ (kFCFA)}^2
\]
Écart-type :
\[
\sigma = \sqrt{2950} \approx 54,3 \text{ kFCFA}
\]
\emph{Interprétation :} Le salaire moyen est de \SI{265000}{FCFA}. Les salaires « typiques » s'écartent en moyenne d'environ \SI{54300}{FCFA} de cette moyenne. La majorité des salaires se situe dans l'intervalle $[210,7 ; 319,3]$ kFCFA.
\end{exoresolu}

\begin{attention}[Erreur fréquente : bornes au lieu des centres]
\emph{Ne jamais utiliser les bornes des classes ($a_i$ ou $b_i$) ni l'amplitude $h$ dans les formules de $\bar{x}$, $V$, $\sigma$.}
Seuls les \textbf{centres de classes $c_i$} représentent les valeurs de la variable dans la classe. Utiliser la borne inférieure sous-estime systématiquement la moyenne et fausse complètement la variance.
\end{attention}

\courssub{Comparaison de deux séries à l'aide de l'écart-type}

L'écart-type permet de comparer la \textbf{dispersion} de deux séries, même si leurs moyennes sont différentes. Si les moyennes sont très différentes, on utilise le \textbf{coefficient de variation} (CV) :
\[
\text{CV} = \frac{\sigma}{\bar{x}} \times 100 \quad (\text{en \%})
\]
Un CV faible ($< 15\%$) indique une série \textbf{homogène} ; un CV élevé indique une série \textbf{hétérogène}.

\begin{popfigure}
\centering
\begin{tikzpicture}
\begin{axis}[
    width=\linewidth, height=7cm,
    axis lines=middle,
    xlabel={Salaire (kFCFA)}, ylabel={Densité},
    xmin=100, xmax=500,
    ymin=0, ymax=0.012,
    xtick={200,265,300,400}, xticklabels={200,265,300,400},
    ytick=\empty,
    clip=false,
    domain=100:500, samples=200,
    legend style={at={(0.98,0.98)}, anchor=north east, fill=white, draw=popmuted}
]
% Entreprise A : moyenne 265, sigma 54 (étroite)
\addplot[popblue, thick, domain=100:500] {1/(54.3*sqrt(max(0,2*pi)))*exp(-(x-265)^2/(2*54.3^2))};
\addlegendentry{Entreprise A : $\bar{x}=265$, $\sigma\approx 54$ (homogène)}
% Entreprise B : moyenne 265, sigma 120 (large)
\addplot[poporange, thick, dashed, domain=100:500] {1/(120*sqrt(max(0,2*pi)))*exp(-(x-265)^2/(2*120^2))};
\addlegendentry{Entreprise B : $\bar{x}=265$, $\sigma=120$ (hétérogène)}
% Lignes moyennes
\draw[popblue, dashed] (axis cs:265,0) -- (axis cs:265,0.0073) node[above, font=\tiny, popblue] {$\bar{x}$};
\draw[poporange, dashed] (axis cs:265,0) -- (axis cs:265,0.0033) node[above, font=\tiny, poporange] {$\bar{x}$};
\end{axis}
\end{tikzpicture}
\legende{Deux distributions de même moyenne (265 kFCFA) mais d'écarts-types différents. L'écart-type $\sigma$ mesure la « largeur » de la cloche.}
\end{popfigure}

\begin{exemplebox}[Comparaison de deux entreprises : même moyenne, dispersions différentes]
L'entreprise A (ci-dessus) a $\bar{x}_A = 265$ kFCFA, $\sigma_A \approx 54$ kFCFA $\Rightarrow$ CV$_A \approx 20,4\%$.
L'entreprise B a la \textbf{même moyenne} $\bar{x}_B = 265$ kFCFA mais $\sigma_B = 120$ kFCFA $\Rightarrow$ CV$_B \approx 45,3\%$.

\emph{Conclusion :} Bien que le salaire moyen soit identique, l'entreprise B est \textbf{bien plus inégalitaire} : les salaires sont très dispersés (écarts importants entre bas et hauts salaires), alors que l'entreprise A pratique une grille de salaires plus resserrée. L'écart-type (ou le CV) révèle ce que la moyenne cache.
\end{exemplebox}

\begin{exoresolu}[Comparaison avec moyennes différentes : coefficient de variation]
\emph{Énoncé :} Deux machines produisent des pièces de diamètre cible 50 mm.
\begin{itemize}
    \item Machine 1 : $\bar{x}_1 = 50,02$ mm, $\sigma_1 = 0,05$ mm.
    \item Machine 2 : $\bar{x}_2 = 49,98$ mm, $\sigma_2 = 0,08$ mm.
\end{itemize}
Quelle machine est la plus \textbf{régulière} (précise) ?

\tcblower
\emph{Solution :}
On compare les coefficients de variation (ou directement les $\sigma$ car les moyennes sont proches).
CV$_1 = \frac{0,05}{50,02}\times 100 \approx 0,10\%$
CV$_2 = \frac{0,08}{49,98}\times 100 \approx 0,16\%$
$\sigma_1 < \sigma_2$ et CV$_1 <$ CV$_2$.
\emph{Réponse :} La \textbf{Machine 1} est plus régulière (dispersion plus faible autour de sa moyenne). Elle produit des pièces de qualité plus constante.
\end{exoresolu}

\courssub{Synthèse : quel indicateur de dispersion choisir ?}

\begin{center}
\begin{tabularx}{\linewidth}{|l|X|X|}\hline
\rowcolor{popblueL}
\textbf{Indicateur} & \textbf{Formule (série regroupée)} & \textbf{Usage principal} \\ \hline
Étendue $E$ & $b_{\text{max}} - a_{\text{min}}$ & Vue rapide de l'amplitude totale. Sensible aux extrêmes. \\ \hline
Écart interquartile EIQ & $Q_3 - Q_1$ & Dispersion des 50\% centraux. Robuste. Classes ouvertes OK. \\ \hline
Variance $V$ & $\frac{1}{N}\sum n_i c_i^2 - \bar{x}^2$ & Calculs théoriques, base tests statistiques. Unité$^2$. \\ \hline
Écart-type $\sigma$ & $\sqrt{V}$ & \textbf{Indicateur roi.} Même unité que $X$. Interprétable ($\bar{x}\pm\sigma$). \\ \hline
Coeff. de variation CV & $\frac{\sigma}{\bar{x}}\times 100$ & \textbf{Comparaison} de séries d'unités ou moyennes différentes. \\ \hline
\end{tabularx}
\end{center}

\begin{aretenir}[Règle d'or pour le Probatoire]
\begin{enumerate}
    \item Toujours calculer les \textbf{centres de classes $c_i$} en premier.
    \item Utiliser \textbf{exclusivement la formule de König-Huygens} pour la variance : $V = \frac{\sum n_i c_i^2}{N} - \bar{x}^2$.
    \item Donner l'écart-type $\sigma$ avec \textbf{l'unité} de la variable.
    \item Pour comparer : même unité + moyennes proches $\to$ comparer $\sigma$ ; sinon $\to$ comparer CV.
\end{enumerate}
\end{aretenir}

\coursec{Histogramme et polygone}

Dans la section précédente, nous avons mesuré la \cle{dispersion} d'une série regroupée en classes à l'aide de nombres : étendue, écart interquartile, variance, écart-type. Ces indicateurs sont précis, mais ils restent abstraits. Un chef d'atelier, un comptable ou un responsable des ressources humaines préfère souvent \emph{voir} la répartition des données d'un seul coup d'œil. C'est exactement le rôle des représentations graphiques que nous étudions maintenant : l'\cle{histogramme}, qui donne une image immédiate de la répartition, et le \cle{polygone des effectifs cumulés}, qui permet de lire graphiquement la médiane. Nous illustrerons toute la section avec la série des salaires mensuels (en milliers de FCFA) des 40 employés d'une PME de Douala, regroupés en classes d'amplitude 50.

\begin{center}
\begin{tabularx}{\linewidth}{|l|X|X|X|X|X|}\hline
\rowcolor{popblueL}
Classes (milliers FCFA) & $[100\,;150[$ & $[150\,;200[$ & $[200\,;250[$ & $[250\,;300[$ & $[300\,;350[$ \\ \hline
Effectifs $n_i$ & 6 & 10 & 14 & 7 & 3 \\ \hline
ECC & 6 & 16 & 30 & 37 & 40 \\ \hline
\end{tabularx}
\end{center}

\courssub{L'histogramme d'une série regroupée en classes}

\textbf{Intuition.} Dans un diagramme en bâtons (série non regroupée), chaque valeur du caractère est représentée par un trait vertical séparé des autres. Mais quand les données sont regroupées en \emph{classes}, chaque classe est un \emph{intervalle} : il est donc naturel de représenter chaque classe par un \emph{rectangle} dont la base recouvre exactement l'intervalle, et dont la hauteur indique combien d'individus tombent dans cette classe. Les rectangles se touchent, car les classes se succèdent sans trou.

\begin{definition}[Histogramme d'une série regroupée en classes de même amplitude]
L'\cle{histogramme} d'une série statistique regroupée en classes de \textbf{même amplitude} est un diagramme formé de \textbf{rectangles juxtaposés} (sans espace entre eux) tels que :
\begin{itemize}
\item la \cle{base} de chaque rectangle est l'intervalle de la classe, porté sur l'axe des abscisses ;
\item la \cle{hauteur} de chaque rectangle est l'effectif $n_i$ de la classe (ou, au choix, sa fréquence $f_i$).
\end{itemize}
Comme toutes les classes ont la même amplitude, les hauteurs sont directement comparables : une classe deux fois plus chargée donne un rectangle deux fois plus haut. L'aire de chaque rectangle est alors proportionnelle à l'effectif de la classe.
\end{definition}

\begin{methode}[Construire un histogramme]
\begin{enumerate}
\item \textbf{Choisir l'échelle des abscisses} : l'axe horizontal est gradué avec les \emph{bornes des classes} (ici 100, 150, 200, 250, 300, 350). Prévoir une graduation régulière, par exemple 1 cm pour 50 milliers de FCFA.
\item \textbf{Choisir l'échelle des ordonnées} : l'axe vertical porte les effectifs, de 0 jusqu'au plus grand effectif (ici 14). Par exemple 0,5 cm par employé.
\item \textbf{Tracer les rectangles} : pour chaque classe, tracer un rectangle dont la base va de la borne inférieure à la borne supérieure, et dont la hauteur est l'effectif. Les rectangles doivent être \textbf{juxtaposés}.
\item \textbf{Soigner la présentation} : hachurer ou colorier les rectangles, donner un \emph{titre} au graphique et nommer les axes (grandeur et unité).
\end{enumerate}
\end{methode}

\begin{attention}[Rectangles juxtaposés, pas de bâtons séparés]
L'erreur la plus fréquente au Probatoire est de laisser des \textbf{espaces entre les rectangles}, comme dans un diagramme en bâtons. C'est faux : les classes se succèdent sans interruption ($[100\,;150[$ puis $[150\,;200[$, etc.), donc les rectangles doivent se toucher. Un histogramme avec des espaces entre les barres est sanctionné. Autre piège : oublier d'indiquer les bornes des classes sur l'axe des abscisses, ce qui rend le graphique illisible.
\end{attention}

\begin{exemplebox}[Histogramme de la série des salaires]
Avec les données de la PME de Douala, on obtient l'histogramme suivant. La classe $[200\,;250[$, qui contient 14 employés, donne le rectangle le plus haut : c'est la \cle{classe modale}.
\end{exemplebox}

\begin{popfigure}
\begin{tikzpicture}
\begin{axis}[
    width=\linewidth, height=7cm,
    ybar interval,
    ymin=0, ymax=16,
    xmin=90, xmax=360,
    xtick={100,150,200,250,300,350},
    xlabel={Salaire (milliers de FCFA)},
    ylabel={Effectif},
    ytick={0,2,4,6,8,10,12,14},
    axis lines=left,
    every axis plot/.append style={fill=popblue!60, draw=popdark},
]
\addplot+[ybar interval] coordinates {(100,6) (150,10) (200,14) (250,7) (300,3) (350,0)};
\end{axis}
\end{tikzpicture}
\legende{Histogramme des salaires mensuels des 40 employés (classes d'amplitude 50 milliers de FCFA).}
\end{popfigure}

\textbf{Lecture d'informations sur un histogramme.} Un histogramme se lit comme une carte :
\begin{itemize}
\item la \cle{classe modale} est visible immédiatement : c'est le rectangle le plus haut (ici $[200\,;250[$) ;
\item la \emph{symétrie} ou l'\emph{asymétrie} de la répartition : ici, la distribution s'étire vers les salaires élevés (étalement à droite), ce qui est typique des distributions de revenus ;
\item l'\emph{étalement} : les salaires vont de 100 à 350 milliers de FCFA, mais l'essentiel des employés (30 sur 40, soit 75 \%) se concentre entre 150 et 300 milliers.
\end{itemize}

\courssub{Polygone des effectifs}

\begin{definition}[Polygone des effectifs]
Le \cle{polygone des effectifs} est la ligne brisée obtenue en joignant par des segments les points de coordonnées $(c_i\,;\,n_i)$, où $c_i$ est le \textbf{centre} de la classe numéro $i$ et $n_i$ son effectif. Par convention, on prolonge la ligne jusqu'à l'axe des abscisses en ajoutant une classe fictive d'effectif nul de chaque côté.
\end{definition}

Pour la série des salaires, les centres de classes sont :
\[ c_1 = \frac{100+150}{2} = 125, \quad c_2 = 175, \quad c_3 = 225, \quad c_4 = 275, \quad c_5 = 325. \]
Le polygone joint donc les points $(125\,;6)$, $(175\,;10)$, $(225\,;14)$, $(275\,;7)$, $(325\,;3)$. Il donne la même information que l'histogramme, mais sous forme de courbe brisée : il est pratique pour \emph{comparer deux séries} sur le même graphique (par exemple les salaires de deux ateliers).

\courssub{Polygone des effectifs cumulés croissants et médiane}

\textbf{Intuition.} L'effectif cumulé croissant répond à la question : « combien d'employés gagnent \emph{moins} qu'un certain montant ? » En traçant cette quantité en fonction du montant, on obtient une courbe qui ne fait que monter. La médiane étant la valeur qui partage l'effectif en deux moitiés égales, il suffit de chercher sur cette courbe le point dont l'ordonnée vaut la moitié de l'effectif total.

\begin{definition}[Polygone des ECC]
Le \cle{polygone des effectifs cumulés croissants} (ECC) est la ligne brisée joignant les points $(b_i\,;\,N_i)$, où $b_i$ est la \textbf{borne supérieure} de la classe numéro $i$ et $N_i$ l'effectif cumulé croissant de cette classe. On fait partir la courbe du point $(b_0\,;\,0)$, où $b_0$ est la borne inférieure de la première classe.
\end{definition}

Pour notre série, le polygone des ECC joint les points :
\[ (100\,;0),\ (150\,;6),\ (200\,;16),\ (250\,;30),\ (300\,;37),\ (350\,;40). \]

\begin{propriete}[Lecture graphique de la médiane — admise]
Soit une série regroupée en classes, d'effectif total $N$. La \cle{médiane} $Me$ est l'\textbf{abscisse du point du polygone des ECC dont l'ordonnée est égale à $\dfrac{N}{2}$}.\\
Construction : on trace la droite horizontale d'équation $y = \dfrac{N}{2}$ ; elle coupe le polygone en un point $M$ ; on projette $M$ sur l'axe des abscisses : la valeur lue est $Me$. Cette lecture graphique est \textbf{cohérente avec le calcul par interpolation linéaire}, car le polygone est précisément formé de segments de droite : chercher le point d'ordonnée $N/2$ sur un segment, c'est exactement faire une interpolation linéaire sur ce segment.
\end{propriete}

\begin{popfigure}
\begin{tikzpicture}
\begin{axis}[
    width=\linewidth, height=8cm,
    xmin=90, xmax=360, ymin=0, ymax=44,
    xtick={100,150,200,214.3,250,300,350},
    x tick label style={rotate=30, anchor=east, font=\small},
    ytick={0,6,16,20,30,37,40},
    xlabel={Salaire (milliers de FCFA)},
    ylabel={Effectif cumul\'e croissant},
    axis lines=left, grid=both, grid style={popline!40},
]
\addplot[popcurve, mark=*, mark options={fill=poporange}] coordinates {(100,0) (150,6) (200,16) (250,30) (300,37) (350,40)};
\addplot[poporange, dashed, thick] coordinates {(90,20) (214.3,20)};
\addplot[poporange, dashed, thick] coordinates {(214.3,0) (214.3,20)};
\node[popdark, anchor=south west] at (axis cs:100,20) {$\dfrac{N}{2}=20$};
\node[popdark, anchor=north] at (axis cs:214.3,0) {$Me$};
\end{axis}
\end{tikzpicture}
\legende{Polygone des ECC et construction graphique de la médiane : on lit $Me \approx 214$ milliers de FCFA.}
\end{popfigure}

\begin{exoresolu}[Histogramme, polygone des ECC et médiane de la série des salaires]
Énoncé. Les salaires mensuels (en milliers de FCFA) des 40 employés d'une PME sont répartis ainsi : $[100\,;150[$ : 6 ; $[150\,;200[$ : 10 ; $[200\,;250[$ : 14 ; $[250\,;300[$ : 7 ; $[300\,;350[$ : 3.
\begin{enumerate}
\item Construire l'histogramme de la série.
\item Construire le polygone des ECC et en déduire graphiquement la médiane.
\item Retrouver la médiane par interpolation linéaire et comparer.
\end{enumerate}
\tcblower
Solution détaillée.
\begin{enumerate}
\item \textbf{Histogramme.} Axe des abscisses : bornes 100, 150, 200, 250, 300, 350 (1 cm pour 50 milliers). Axe des ordonnées : effectifs de 0 à 14 (0,5 cm par employé). On trace cinq rectangles \emph{juxtaposés} de hauteurs 6, 10, 14, 7, 3 (voir la figure 1). Le rectangle le plus haut correspond à la classe modale $[200\,;250[$.
\item \textbf{Médiane graphique.} L'effectif total est $N = 40$, donc $\dfrac{N}{2} = 20$. On trace la droite horizontale $y = 20$ : elle coupe le polygone des ECC sur le segment joignant $(200\,;16)$ à $(250\,;30)$. La projection du point d'intersection sur l'axe des abscisses donne $Me \approx 214$ milliers de FCFA (voir la figure 2).
\item \textbf{Interpolation linéaire.} La médiane appartient à la classe médiane $[200\,;250[$, car $16 < 20 \leq 30$. La formule d'interpolation donne :
\begin{align*}
Me &= 200 + 50 \times \frac{20 - 16}{30 - 16} \\
   &= 200 + 50 \times \frac{4}{14} \\
   &= 200 + \frac{200}{14} \approx 200 + 14{,}3 \approx 214{,}3.
\end{align*}
\textbf{Comparaison.} La lecture graphique ($Me \approx 214$) et le calcul ($Me \approx 214{,}3$) donnent le même résultat, aux erreurs de lecture près. C'est normal : la lecture graphique \emph{est} une interpolation linéaire, réalisée avec une règle au lieu d'une formule.
\end{enumerate}
Conclusion : la moitié des employés gagne moins d'environ 214 300 FCFA par mois, et l'autre moitié davantage.
\end{exoresolu}

\begin{attention}[Borne supérieure ou centre de classe ?]
Ne confonds pas les deux polygones :
\begin{itemize}
\item \textbf{polygone des effectifs} : points $(\text{centre de la classe}\,;\,n_i)$ ;
\item \textbf{polygone des ECC} : points $(\text{borne supérieure de la classe}\,;\,N_i)$.
\end{itemize}
Placer les effectifs cumulés au centre des classes est une erreur classique qui décale toute la courbe et fausse la lecture de la médiane. Retiens : « cumulé $\Rightarrow$ borne supérieure », car l'ECC répond à « moins de $b_i$ ».
\end{attention}

\begin{methode}[Lire la médiane sur un polygone des ECC]
\begin{enumerate}
\item Calculer $N$ puis $\dfrac{N}{2}$.
\item Repérer $\dfrac{N}{2}$ sur l'axe des ordonnées et tracer la droite horizontale correspondante.
\item Repérer le point d'intersection avec le polygone (il se trouve sur le segment de la classe médiane).
\item Projeter ce point verticalement sur l'axe des abscisses et lire la valeur de $Me$ avec l'échelle.
\item Vérifier la cohérence : $Me$ doit appartenir à la classe médiane trouvée dans le tableau.
\end{enumerate}
\end{methode}

\begin{savaistu}[Et si les classes n'ont pas la même amplitude ?]
Dans tout ce cours, les classes ont la même amplitude, donc les \emph{hauteurs} des rectangles sont proportionnelles aux effectifs. Mais dans la vie réelle (recensements de l'INS Cameroun, bulletins météo), on rencontre parfois des classes d'amplitudes différentes, par exemple $[0\,;10[$ puis $[10\,;50[$. Dans ce cas — ce sera utile en Terminale — ce sont les \textbf{aires} des rectangles qui doivent être proportionnelles aux effectifs, et non leurs hauteurs : on divise chaque effectif par l'amplitude de la classe pour obtenir la hauteur (appelée \emph{densité}). Un rectangle deux fois plus large doit être deux fois plus bas pour représenter le même effectif. Au programme de Première technique, on se limite aux classes de même amplitude : les hauteurs suffisent.
\end{savaistu}

\begin{exercice}[Atelier de mécanique à Yaoundé]
Les durées de vie (en centaines d'heures) de 30 roulements testés dans un atelier sont regroupées ainsi : $[10\,;20[$ : 4 ; $[20\,;30[$ : 9 ; $[30\,;40[$ : 11 ; $[40\,;50[$ : 5 ; $[50\,;60[$ : 1.
\begin{enumerate}
\item Construire l'histogramme de cette série sur papier millimétré (échelles au choix, à justifier).
\item Dresser le tableau des ECC, construire le polygone des ECC et lire graphiquement la médiane.
\item Calculer la médiane par interpolation linéaire et comparer avec la lecture graphique.
\end{enumerate}
\end{exercice}

\begin{corrige}[Exercice : atelier de mécanique]
\begin{enumerate}
\item Abscisses : bornes 10 à 60 (par exemple 2 cm pour 10 centaines d'heures) ; ordonnées : effectifs de 0 à 11 (1 cm pour 2 roulements). Cinq rectangles juxtaposés de hauteurs 4, 9, 11, 5, 1. Classe modale : $[30\,;40[$.
\item ECC : 4 ; 13 ; 24 ; 29 ; 30. Points du polygone : $(10\,;0)$, $(20\,;4)$, $(30\,;13)$, $(40\,;24)$, $(50\,;29)$, $(60\,;30)$. On a $\dfrac{N}{2} = 15$ : la droite $y = 15$ coupe le segment joignant $(30\,;13)$ à $(40\,;24)$ ; on lit $Me \approx 32$ centaines d'heures.
\item Interpolation : $Me = 30 + 10 \times \dfrac{15 - 13}{24 - 13} = 30 + \dfrac{20}{11} \approx 31{,}8$ centaines d'heures, soit environ 3 180 heures. Lecture graphique et calcul concordent.
\end{enumerate}
\end{corrige}

\begin{aretenir}[L'essentiel de la section]
\begin{itemize}
\item L'\cle{histogramme} d'une série en classes de même amplitude est formé de rectangles \textbf{juxtaposés} : base = classe, hauteur = effectif (ou fréquence).
\item Le \cle{polygone des effectifs} joint les points (centre ; effectif) ; le \cle{polygone des ECC} joint les points (borne supérieure ; effectif cumulé).
\item La \cle{médiane} se lit graphiquement sur le polygone des ECC : c'est l'abscisse du point d'ordonnée $\dfrac{N}{2}$. Cette lecture équivaut au calcul par interpolation linéaire.
\item Piège à éviter : des espaces entre les rectangles de l'histogramme, ou des ECC placés aux centres des classes.
\end{itemize}
\end{aretenir}

\coursec{Synthèse et applications à la vie courante}

Après avoir maîtrisé la construction d'histogrammes et de polygones des fréquences, nous disposons désormais de tous les outils pour mener une étude statistique complète. Cette section synthétise la démarche globale, insiste sur la rédaction rigoureuse d'une conclusion — compétence clé au Probatoire — et montre comment comparer deux populations réelles. Nous illustrerons le tout par une étude détaillée de notes de deux classes de Première technique du lycée de Bafoussam, et nous verrons comment l'Institut National de la Statistique (INS) utilise exactement ces méthodes pour les enquêtes nationales (ECAM).

\courssub{La démarche complète d'étude d'une série statistique regroupée}

Face à un jeu de données brutes (notes, salaires, rendements, diamètres de pièces...), la démarche suit toujours les mêmes étapes logiques. La maîtriser permet de ne rien oublier et de structurer sa copie d'examen.

\begin{methode}[Démarche type pour l'étude d'une série regroupée en classes de même amplitude]
\begin{enumerate}
    \item \textbf{Organiser les données :} Construire le tableau de répartition avec les classes (borne inférieure, borne supérieure, milieu), les effectifs $n_i$, les effectifs cumulés croissants $N_i$, les fréquences $f_i = \frac{n_i}{N}$ (souvent en \%) et les fréquences cumulées $F_i$.
    \item \textbf{Calculer les caractéristiques de position :} Moyenne $\bar{x}$ (centre de gravité), Médiane $Me$ (valeur centrale), Quartiles $Q_1, Q_3$ (encadrement de 50\% des données), Mode (classe modale).
    \item \textbf{Calculer les caractéristiques de dispersion :} Étendue $E$, Écart interquartile $EIQ = Q_3 - Q_1$, Variance $s^2$ (ou $\sigma^2$), Écart-type $s$ (ou $\sigma$), Coefficient de variation $CV = \frac{s}{\bar{x}} \times 100$ (pour comparer des séries d'unités ou d'échelles différentes).
    \item \textbf{Représenter graphiquement :} Tracer l'histogramme (aires proportionnelles aux effectifs car amplitude constante) et/ou le polygone des fréquences. Ajouter si possible la médiane, les quartiles, la moyenne $\pm$ écart-type.
    \item \textbf{Interpréter et conclure :} Rédiger une réponse en langage courant, en citant \emph{explicitement} les valeurs calculées, en comparant position et dispersion, et en répondant à la question posée.
\end{enumerate}
\end{methode}

\begin{popfigure}
\begin{tikzpicture}[node distance=1.2cm and 1.5cm, >=Stealth,
    box/.style={draw=popblue, thick, rounded corners, fill=popblueL, text width=3.2cm, align=center, minimum height=1.2cm},
    arrow/.style={->, thick, popdark}
]
\node[box, align=center] (brut){Données brutes\\ (liste de valeurs)};
\node[box, right=of brut, align=center] (tab){Tableau de répartition\\ Classes, $n_i$, $N_i$, $f_i$, $F_i$};
\node[box, below=of tab, align=center] (ind){Indicateurs numériques\\ $\bar{x}$, $Me$, $Q_1$, $Q_3$, $s$, $CV$};
\node[box, right=of ind, align=center] (graph){Représentation graphique\\ Histogramme, Polygone};
\node[box, below=of graph, align=center] (conc){Interprétation \& Conclusion\\ Réponse argumentée};
\draw[arrow] (brut) -- (tab);
\draw[arrow] (tab) -- (ind);
\draw[arrow] (tab) -- (graph);
\draw[arrow] (ind) -- (conc);
\draw[arrow] (graph) -- (conc);
\end{tikzpicture}
\legende{Organigramme de la démarche statistique complète : des données brutes à la conclusion argumentée.}
\end{popfigure}

\courssub{Rédiger et justifier une conclusion statistique}

Au Probatoire, il ne suffit pas de donner des chiffres ; il faut \textbf{parler le langage des chiffres}. Une conclusion correcte cite les indicateurs, les compare le cas échéant, et répond à la problématique en français courant.

\begin{methode}[Rédiger une conclusion argumentée (modèle pour l'examen)]
\begin{enumerate}
    \item \textbf{Annoncer le résultat principal :} « La moyenne des notes est de $\bar{x} = 11,3$ sur 20. »
    \item \textbf{Préciser la dispersion :} « L'écart-type $s \approx 3,3$ montre un écart moyen modéré autour de la moyenne. » ou « Le coefficient de variation $CV = 29\%$ indique une hétérogénéité notable. »
    \item \textbf{Situer la médiane / quartiles :} « La médiane $Me = 11,3$ ; 50\% des élèves ont entre $Q_1 = 9,0$ et $Q_3 = 13,7$. »
    \item \textbf{Comparer (si deux séries) :} « La classe A a une moyenne supérieure ($11,3$ vs $9,7$) \textbf{et} une dispersion relative plus faible ($CV=29\%$ vs $CV=34\%$), elle est donc globalement meilleure et plus homogène. »
    \item \textbf{Conclure en langage courant :} « On peut dire que le niveau global est moyen, avec une minorité d'élèves en grande difficulté (notes $< 8$). »
\end{enumerate}
\end{methode}

\begin{attention}[Piège majeur : conclure sur la seule moyenne]
Deux séries peuvent avoir \textbf{exactement la même moyenne} mais des réalités totalement différentes.
\begin{itemize}
    \item Série 1 : Notes $10, 10, 10, 10, 10$ $\rightarrow$ $\bar{x}=10$, $s=0$ (homogène, tout le monde à la moyenne).
    \item Série 2 : Notes $0, 0, 10, 20, 20$ $\rightarrow$ $\bar{x}=10$, $s \approx 8,9$ (très dispersée, la moitié a échoué, l'autre a excellé).
\end{itemize}
\textbf{Jamais de conclusion sans citer la dispersion (écart-type, $CV$, $EIQ$ ou étendue).} C'est l'erreur n°1 qui fait perdre le plus de points au Probatoire.
\end{attention}

\courssub{Comparaison de deux populations : étude complète et chiffrée}

Comparer deux groupes (deux classes, deux machines, deux champs, deux traitements médicaux) est l'application reine des statistiques descriptives. On compare \textbf{position} (moyenne, médiane) \textbf{ET} \textbf{dispersion} (écart-type, $CV$, $EIQ$).

\begin{exemplebox}[Étude comparée : Notes de deux classes de Première Tech au Lycée de Bafoussam]
On a relevé les notes sur 20 au devoir commun de mathématiques de deux classes de Première technique (effectif total $N=60$ par classe). Les notes sont regroupées en classes d'amplitude $2$ : $[4;6[$, $[6;8[$, $[8;10[$, $[10;12[$, $[12;14[$, $[14;16[$, $[16;18[$, $[18;20]$.

\begin{center}
\begin{tabularx}{\linewidth}{|c|X|X|X|X|}\hline
\rowcolor{popblueL}
\textbf{Classe (notes)} & \textbf{Milieu $x_i$} & \textbf{Effectifs Classe A ($n_i^A$)} & \textbf{Effectifs Classe B ($n_i^B$)} & \textbf{Eff. cumulés A / B} \\ \hline
$[4;6[$ & 5 & 4 & 8 & 4 / 8 \\
$[6;8[$ & 7 & 6 & 12 & 10 / 20 \\
$[8;10[$ & 9 & 10 & 15 & 20 / 35 \\
$[10;12[$ & 11 & 15 & 10 & 35 / 45 \\
$[12;14[$ & 13 & 12 & 8 & 47 / 53 \\
$[14;16[$ & 15 & 8 & 5 & 55 / 58 \\
$[16;18[$ & 17 & 4 & 1 & 59 / 59 \\
$[18;20]$ & 19 & 1 & 1 & 60 / 60 \\ \hline
\textbf{Total} & & \textbf{60} & \textbf{60} & \\ \hline
\end{tabularx}
\end{center}

\textbf{1. Calculs des caractéristiques (détail pour la Classe A, résultats pour les deux) :}

\emph{Classe A :}
\begin{align*}
\sum n_i x_i &= 4(5) + 6(7) + 10(9) + 15(11) + 12(13) + 8(15) + 4(17) + 1(19) = 680 \\
\bar{x}_A &= \frac{680}{60} \approx \textbf{11,33} \\
\sum n_i x_i^2 &= 4(25) + 6(49) + 10(81) + 15(121) + 12(169) + 8(225) + 4(289) + 1(361) = 8364 \\
s_A^2 &= \frac{8364}{60} - 11,33^2 \approx 139,4 - 128,4 = \textbf{11,0} \quad \Rightarrow \quad s_A \approx \textbf{3,31} \\
CV_A &= \frac{3,31}{11,33} \times 100 \approx \textbf{29,2\%}
\end{align*}

\emph{Classe B :}
\begin{align*}
\sum n_i x_i &= 8(5) + 12(7) + 15(9) + 10(11) + 8(13) + 5(15) + 1(17) + 1(19) = 584 \\
\bar{x}_B &= \frac{584}{60} \approx \textbf{9,73} \\
\sum n_i x_i^2 &= 8(25) + 12(49) + 15(81) + 10(121) + 8(169) + 5(225) + 1(289) + 1(361) = 6340 \\
s_B^2 &= \frac{6340}{60} - 9,73^2 \approx 105,67 - 94,67 = \textbf{11,0} \quad \Rightarrow \quad s_B \approx \textbf{3,31} \\
CV_B &= \frac{3,31}{9,73} \times 100 \approx \textbf{34,0\%}
\end{align*}

\textbf{2. Médianes et quartiles (interpolation linéaire, $h=2$) :}
\begin{itemize}
    \item \textbf{Classe A :} $N=60$. Médiane (30\textsuperscript{e}-31\textsuperscript{e}) : classe $[10;12[$, $N_{\text{ant}}=20$, $n=15$.
    $Me_A \approx 10 + \frac{30-20}{15} \times 2 = \textbf{11,33}$.
    $Q_1^A$ (15\textsuperscript{e}) : classe $[8;10[$, $N_{\text{ant}}=10$, $n=10$.
    $Q_1^A \approx 8 + \frac{15-10}{10} \times 2 = \textbf{9,0}$.
    $Q_3^A$ (45\textsuperscript{e}) : classe $[12;14[$, $N_{\text{ant}}=35$, $n=12$.
    $Q_3^A \approx 12 + \frac{45-35}{12} \times 2 = \textbf{13,67}$.
    $EIQ_A \approx 4,67$.
    \item \textbf{Classe B :} Médiane (30\textsuperscript{e}-31\textsuperscript{e}) : classe $[8;10[$, $N_{\text{ant}}=20$, $n=15$.
    $Me_B \approx 8 + \frac{30-20}{15} \times 2 = \textbf{9,33}$.
    $Q_1^B$ (15\textsuperscript{e}) : classe $[6;8[$, $N_{\text{ant}}=8$, $n=12$.
    $Q_1^B \approx 6 + \frac{15-8}{12} \times 2 = \textbf{7,17}$.
    $Q_3^B$ (45\textsuperscript{e}) : classe $[10;12[$, $N_{\text{ant}}=35$, $n=10$.
    $Q_3^B \approx 10 + \frac{45-35}{10} \times 2 = \textbf{12,0}$.
    $EIQ_B \approx 4,83$.
\end{itemize}

\textbf{3. Histogrammes comparés (même échelle) :}
\begin{popfigure}
\begin{center}
\begin{tabular}{cc}
\begin{tikzpicture}
\begin{axis}[
    width=0.45\linewidth, height=6cm,
    ybar=0pt, bar width=6pt,
    xlabel={Notes}, ylabel={Effectifs},
    xtick={5,7,9,11,13,15,17,19},
    xticklabels={$[4;6[$, $[6;8[$, $[8;10[$, $[10;12[$, $[12;14[$, $[14;16[$, $[16;18[$, $[18;20]$},
    x tick label style={rotate=45, anchor=east, font=\tiny},
    ymin=0, ymax=18,
    title={Classe A ($\bar{x}\approx11,3$ ; $s\approx3,3$)},
    title style={font=\bfseries\small},
    legend style={at={(0.5,-0.3)}, anchor=north, font=\tiny}
]
\addplot[fill=popblue!70, draw=popblue] coordinates {(5,4) (7,6) (9,10) (11,15) (13,12) (15,8) (17,4) (19,1)};
\addplot[poporange, very thick, dashed] coordinates {(11.33,0) (11.33,16)};
\legend{Effectifs, $\bar{x}_A$}
\end{axis}
\end{tikzpicture}
&
\begin{tikzpicture}
\begin{axis}[
    width=0.45\linewidth, height=6cm,
    ybar=0pt, bar width=6pt,
    xlabel={Notes}, ylabel={Effectifs},
    xtick={5,7,9,11,13,15,17,19},
    xticklabels={$[4;6[$, $[6;8[$, $[8;10[$, $[10;12[$, $[12;14[$, $[14;16[$, $[16;18[$, $[18;20]$},
    x tick label style={rotate=45, anchor=east, font=\tiny},
    ymin=0, ymax=18,
    title={Classe B ($\bar{x}\approx9,7$ ; $s\approx3,3$)},
    title style={font=\bfseries\small},
    legend style={at={(0.5,-0.3)}, anchor=north, font=\tiny}
]
\addplot[fill=popgreen!70, draw=popgreen] coordinates {(5,8) (7,12) (9,15) (11,10) (13,8) (15,5) (17,1) (19,1)};
\addplot[poporange, very thick, dashed] coordinates {(9.73,0) (9.73,16)};
\legend{Effectifs, $\bar{x}_B$}
\end{axis}
\end{tikzpicture}
\end{tabular}
\end{center}
\legende{Histogrammes des notes des classes A et B (même échelle). La classe A est centrée plus haut (moyenne 11,3 vs 9,7) avec une dispersion absolue similaire mais relative plus faible.}
\end{popfigure}

\textbf{4. Conclusion argumentée (modèle de rédaction pour le Probatoire) :}
\begin{quote}
\emph{« La classe A obtient des résultats globalement supérieurs à ceux de la classe B. En effet, la moyenne de la classe A ($\bar{x}_A \approx 11,3$) est supérieure de 1,6 point à celle de la classe B ($\bar{x}_B \approx 9,7$). Les deux classes ont un écart-type absolu quasi identique ($s \approx 3,3$), mais le coefficient de variation de la classe A ($CV_A \approx 29\%$) est inférieur à celui de la classe B ($CV_B \approx 34\%$), ce qui prouve une meilleure homogénéité relative pour la classe A. La médiane confirme ce constat ($Me_A \approx 11,3$ vs $Me_B \approx 9,3$). On note aussi que la classe B compte deux fois plus d'élèves en difficulté (20 élèves $< 8$ contre 10 pour la classe A). En conclusion, le niveau de la classe A est meilleur et plus homogène. »}
\end{quote}
\end{exemplebox}

\courssub{Applications concrètes à la vie courante et professionnelle}

Les séries regroupées en classes ne sont pas un exercice scolaire abstrait : c'est \emph{le} format standard de publication des données massives.

\begin{savaistu}[L'INS et les enquêtes ECAM : la statistique au service de la décision publique]
L'\textbf{Institut National de la Statistique (INS)} du Cameroun réalise régulièrement l'\textbf{Enquête Camerounaise auprès des Ménages (ECAM)}. Les résultats (niveau de vie, pauvreté, emploi, éducation, santé) sont publiés sous forme de \textbf{tableaux de séries statistiques regroupées en classes} (tranches d'âge, tranches de dépenses par tête, quintiles de niveau de vie).
\begin{itemize}
    \item Exemple : « Répartition de la population selon la dépense de consommation par adulte équivalent » $\rightarrow$ classes de dépenses, effectifs (ménages), fréquences.
    \item Les indicateurs calculés (médiane des dépenses, écart interquartile, taux de pauvreté = proportion sous un seuil) orientent les politiques publiques (allocation budgétaire, ciblage des zones prioritaires).
\end{itemize}
Maîtriser cette leçon, c'est être capable de lire et critiquer les rapports officiels de son pays.
\end{savaistu}

\begin{aretenir}[Principaux domaines d'application en technique]
\begin{itemize}
    \item \textbf{Bulletins de notes / Orientation :} Moyenne de classe, écart-type pour juger l'homogénéité, quartiles pour les seuils de mention (Assez bien, Bien, Très bien).
    \item \textbf{Gestion de stock / Approvisionnement :} Distribution des demandes journalières (classes de quantités) $\rightarrow$ calcul du stock de sécurité (lié à l'écart-type), point de commande.
    \item \textbf{Contrôle qualité en entreprise (usinage, agro-alimentaire, BTP) :} Diamètres de pièces, teneur en humidité, résistance à la compression. \textbf{Règle des $3\sigma$ :} Si le processus est centré et stable, $99,7\%$ des valeurs sont dans $[\mu - 3\sigma ; \mu + 3\sigma]$. Tout point hors limites $\rightarrow$ alerte.
    \item \textbf{Agriculture / Élevage :} Rendements à l'hectare (cacao, maïs), poids des porcs à l'abattage. Comparaison variétés / aliments via moyenne et $CV$.
    \item \textbf{Électricité / Maintenance :} Durée de vie de lampes/transformateurs (classes de temps), consommation horaire (classes de kWh) pour dimensionner les câbles et transformateurs.
\end{itemize}
\end{aretenir}

\courssub{Limites du regroupement en classes : perte d'information et valeurs approchées}

Regrouper en classes simplifie le traitement de gros volumes, mais ce n'est pas neutre.

\begin{attention}[Ce que le regroupement masque — Points de vigilance]
\begin{enumerate}
    \item \textbf{Perte de l'information individuelle :} On ne connaît plus les valeurs exactes, seulement les milieux $x_i$. \textbf{Tous les calculs (moyenne, variance, médiane, quartiles) deviennent des \emph{valeurs approchées}.} La précision dépend de l'amplitude des classes (plus $h$ est petit, meilleure est l'approximation).
    \item \textbf{Hypothèse de répartition uniforme dans la classe :} Pour la médiane et les quartiles (formule d'interpolation linéaire), on suppose que les valeurs sont régulièrement réparties à l'intérieur de la classe. C'est souvent faux (ex. : notes souvent entières, arrondies).
    \item \textbf{Choix des bornes (effet de seuil) :} Déplacer une borne de $0,5$ peut changer la classe modale, la médiane, et même l'allure de l'histogramme. \textbf{Toujours vérifier la robustesse :} si on change légèrement l'amplitude ou l'origine, la conclusion change-t-elle ?
    \item \textbf{Classes ouvertes (première et dernière) :} Si la première classe est $]-\infty; a[$ ou la dernière $[b; +\infty[$, on \textbf{ne peut pas calculer la moyenne ni l'écart-type} (milieu inconnu). On ne peut donner que la médiane et les quartiles (si l'effectif total est connu).
    \item \textbf{Amplitudes inégales :} Si les classes n'ont pas la même amplitude, l'histogramme doit avoir des \textbf{aires} proportionnelles aux effectifs (hauteur = densité = $n_i / h_i$), pas les hauteurs. La formule de la moyenne reste $\bar{x} = \frac{\sum n_i x_i}{N}$ avec $x_i$ = milieu, mais la variance approximée est moins fiable.
\end{enumerate}
\end{attention}

\begin{propriete}[Critère de choix du nombre de classes (Règle de Sturges — indicatif)]
Pour un effectif total $N$, un nombre de classes $k$ souvent utilisé est :
\[ k \approx 1 + 3,3 \log_{10} N \quad \text{ou} \quad k \approx \sqrt{N} \]
L'amplitude $h$ se déduit : $h = \frac{\text{Étendue}}{k}$, arrondie à une valeur « commode » (1, 2, 5, 10, 20...).
\textit{Note : Au Probatoire, les classes sont généralement données.}
\end{propriete}

\courssub{Exercice résumé : Synthèse complète sur un cas d'entreprise}

\begin{exoresolu}[Contrôle qualité d'une chaîne de conditionnement — Jus d'ananas « Tropic'Sun »]
Une entreprise agro-alimentaire de Douala conditionne des bouteilles de jus d'ananas étiquetées « 1 Litre ». Un échantillon de $N=80$ bouteilles est prélevé. Le volume réel (en mL) est mesuré. On regroupe en classes d'amplitude $h=5$ mL.

\begin{center}
\begin{tabularx}{\linewidth}{|c|X|X|X|X|}\hline
\rowcolor{popblueL}
\textbf{Classe (mL)} & \textbf{Milieu $x_i$} & \textbf{Effectif $n_i$} & \textbf{Eff. cumulé $N_i$} & \textbf{Fréquence $f_i$ (\%)} \\ \hline
$[985; 990[$ & 987,5 & 4 & 4 & 5,0 \\
$[990; 995[$ & 992,5 & 8 & 12 & 10,0 \\
$[995; 1000[$ & 997,5 & 18 & 30 & 22,5 \\
$[1000; 1005[$ & 1002,5 & 25 & 55 & 31,25 \\
$[1005; 1010[$ & 1007,5 & 15 & 70 & 18,75 \\
$[1010; 1015[$ & 1012,5 & 7 & 77 & 8,75 \\
$[1015; 1020[$ & 1017,5 & 3 & 80 & 3,75 \\ \hline
\textbf{Total} & & \textbf{80} & & \textbf{100} \\ \hline
\end{tabularx}
\end{center}

\textbf{Questions :}
\begin{enumerate}
    \item Calculer la moyenne $\bar{x}$ et l'écart-type $s$ (valeurs approchées).
    \item Déterminer la médiane $Me$ et les quartiles $Q_1, Q_3$.
    \item L'entreprise garantit « 1 L $\pm$ 10 mL » (soit $[990; 1010]$). Quel pourcentage de bouteilles est hors spécification ?
    \item Rédiger une conclusion technique pour le responsable qualité.
\end{enumerate}

\tcblower

\textbf{Solution détaillée}

\textbf{1. Moyenne et écart-type}
\begin{align*}
\sum n_i x_i &= 4(987,5) + 8(992,5) + 18(997,5) + 25(1002,5) + 15(1007,5) + 7(1012,5) + 3(1017,5) \\
&= 3950 + 7940 + 17955 + 25062,5 + 15112,5 + 7087,5 + 3052,5 = \textbf{80160} \\
\bar{x} &= \frac{80160}{80} = \textbf{1002,0 mL}
\end{align*}
\begin{align*}
\sum n_i x_i^2 &= 4(987,5^2) + 8(992,5^2) + 18(997,5^2) + 25(1002,5^2) + 15(1007,5^2) + 7(1012,5^2) + 3(1017,5^2) \\
&= 3\,900\,625 + 7\,880\,450 + 17\,910\,112,5 + 25\,125\,156,25 + 15\,225\,843,75 + 7\,176\,093,75 + 3\,105\,918,75 \\
&= \textbf{80\,324\,200} \\
s^2 &= \frac{80\,324\,200}{80} - 1002,0^2 = 1\,004\,052,5 - 1\,004\,004 = \textbf{48,5} \\
s &\approx \sqrt{48,5} \approx \textbf{6,96 mL} \quad (\text{arrondi à } \textbf{7,0 mL}) \\
CV &= \frac{6,96}{1002} \times 100 \approx \textbf{0,69\%}
\end{align*}
\textit{Note : L'écart-type d'environ 7 mL indique une bonne précision de la machine, mais la moyenne est décalée de +2 mL par rapport à la cible 1000 mL.}

\textbf{2. Médiane et quartiles}
$N=80$. Médiane = moyenne des 40\textsuperscript{e} et 41\textsuperscript{e} valeurs. $N_{3}=30 < 40 \le N_{4}=55 \Rightarrow$ classe médiane $[1000; 1005[$.
\[
Me \approx 1000 + \frac{40 - 30}{25} \times 5 = 1000 + 2 = \textbf{1002,0 mL}
\]
$Q_1$ (20\textsuperscript{e} et 21\textsuperscript{e}) : $N_{2}=12 < 20 \le N_{3}=30 \Rightarrow$ classe $[995; 1000[$.
\[
Q_1 \approx 995 + \frac{20 - 12}{18} \times 5 = 995 + 2,22 \approx \textbf{997,2 mL}
\]
$Q_3$ (60\textsuperscript{e} et 61\textsuperscript{e}) : $N_{4}=55 < 60 \le N_{5}=70 \Rightarrow$ classe $[1005; 1010[$.
\[
Q_3 \approx 1005 + \frac{60 - 55}{15} \times 5 = 1005 + 1,67 \approx \textbf{1006,7 mL}
\]
$EIQ = Q_3 - Q_1 \approx 9,5$ mL. 50\% des bouteilles contiennent entre 997,2 mL et 1006,7 mL.

\textbf{3. Pourcentage hors spécification $[990; 1010]$}
Hors specs = classes $[985;990[$ (4 bouteilles) + $[1010;1015[$ (7) + $[1015;1020[$ (3) = 14 bouteilles.
\[
\frac{14}{80} \times 100 = \textbf{17,5\%}
\]
\textbf{Alerte :} 17,5\% de non-conformité est inacceptable pour une norme $\pm 10$ mL (on vise généralement $< 0,3\%$ soit $3\sigma$).

\textbf{4. Conclusion technique (modèle)}
\begin{quote}
\emph{« Monsieur le Responsable Qualité, l'analyse de l'échantillon de 80 bouteilles montre que le volume moyen est de 1002,0 mL (légèrement au-dessus de l'étiquetage 1000 mL) avec une dispersion modérée ($s \approx 7,0$ mL, $CV \approx 0,7\%$). Le processus n'est pas centré : la médiane (1002 mL) confirme un décalage de +2 mL. Surtout, \textbf{17,5\% des bouteilles sont hors tolérance $[990;1010]$} (principalement en sur-remplissage $>1010$ mL). Bien que l'écart-type soit correct, le décalage de la moyenne génère un taux de rebut énorme (coût matière perdu, risque légal). \textbf{Action corrective urgente :} recentrer la remplisseuse sur 998--999 mL pour ramener le taux de non-conformité sous 0,3\% (règle des $3\sigma$ : $3 \times 7 \approx 21$ mL, intervalle $[979; 1021]$ centré sur 1000 couvre la spec). »}
\end{quote}
\end{exoresolu}

\begin{exercice}[Entraînement : Comparaison de deux fournisseurs de câbles électriques]
Deux fournisseurs (Alpha et Beta) proposent des câbles cuivre de section nominale 2,5 mm². La résistance électrique au mètre (en m$\Omega$/m) est mesurée sur 50 échantillons par fournisseur. Séries regroupées en classes d'amplitude 0,2 m$\Omega$/m.

\begin{center}
\begin{tabularx}{\linewidth}{|c|X|X|}\hline
\rowcolor{popblueL}
\textbf{Classe (m$\Omega$/m)} & \textbf{Effectifs Alpha} & \textbf{Effectifs Beta} \\ \hline
$[7,0; 7,2[$ & 5 & 2 \\
$[7,2; 7,4[$ & 10 & 8 \\
$[7,4; 7,6[$ & 20 & 15 \\
$[7,6; 7,8[$ & 10 & 15 \\
$[7,8; 8,0[$ & 5 & 10 \\ \hline
\textbf{Total} & \textbf{50} & \textbf{50} \\ \hline
\end{tabularx}
\end{center}
\begin{enumerate}
    \item Calculer $\bar{x}$, $s$, $CV$ pour chaque fournisseur.
    \item Déterminer $Me$, $Q_1$, $Q_3$ pour chaque fournisseur.
    \item Tracer les deux histogrammes sur le même graphique (ou côte à côte).
    \item \textbf{Choisir le fournisseur} en justifiant par une conclusion argumentée (spécification client : résistance $< 7,8$ m$\Omega$/m).
\end{enumerate}
\end{exercice}

\begin{corrige}[Exercice : Comparaison fournisseurs]
\textbf{Calculs détaillés (milieux : 7,1 ; 7,3 ; 7,5 ; 7,7 ; 7,9)}

\textbf{Fournisseur Alpha ($N=50$) :}
\begin{align*}
\sum n_i x_i &= 5(7,1) + 10(7,3) + 20(7,5) + 10(7,7) + 5(7,9) = 375,0 \\
\bar{x}_\alpha &= \frac{375}{50} = \textbf{7,50 m$\Omega$/m} \\
\sum n_i x_i^2 &= 5(50,41) + 10(53,29) + 20(56,25) + 10(59,29) + 5(62,41) = 2814,9 \\
s_\alpha^2 &= \frac{2814,9}{50} - 7,50^2 = 56,298 - 56,25 = 0,048 \quad \Rightarrow \quad s_\alpha \approx \textbf{0,22 m$\Omega$/m} \\
CV_\alpha &= \frac{0,22}{7,50} \times 100 \approx \textbf{2,9\%}
\end{align*}
Médiane (25\textsuperscript{e}-26\textsuperscript{e}) : classe $[7,4;7,6[$, $N_{\text{ant}}=15$, $n=20$.
$Me_\alpha \approx 7,4 + \frac{25-15}{20} \times 0,2 = \textbf{7,50}$.
$Q_1$ (12,5\textsuperscript{e}) : classe $[7,2;7,4[$, $N_{\text{ant}}=5$, $n=10$.
$Q_1^\alpha \approx 7,2 + \frac{12,5-5}{10} \times 0,2 = \textbf{7,35}$.
$Q_3$ (37,5\textsuperscript{e}) : classe $[7,6;7,8[$, $N_{\text{ant}}=35$, $n=10$.
$Q_3^\alpha \approx 7,6 + \frac{37,5-35}{10} \times 0,2 = \textbf{7,65}$.

\textbf{Fournisseur Beta ($N=50$) :}
\begin{align*}
\sum n_i x_i &= 2(7,1) + 8(7,3) + 15(7,5) + 15(7,7) + 10(7,9) = 379,6 \\
\bar{x}_\beta &= \frac{379,6}{50} = \textbf{7,59 m$\Omega$/m} \\
\sum n_i x_i^2 &= 2(50,41) + 8(53,29) + 15(56,25) + 15(59,29) + 10(62,41) = 2884,34 \\
s_\beta^2 &= \frac{2884,34}{50} - 7,59^2 \approx 57,687 - 57,608 = 0,079 \quad \Rightarrow \quad s_\beta \approx \textbf{0,28 m$\Omega$/m} \\
CV_\beta &= \frac{0,28}{7,59} \times 100 \approx \textbf{3,7\%}
\end{align*}
Médiane (25\textsuperscript{e}-26\textsuperscript{e}) : classe $[7,4;7,6[$, $N_{\text{ant}}=10$, $n=15$.
$Me_\beta \approx 7,4 + \frac{25-10}{15} \times 0,2 = \textbf{7,60}$.
$Q_1$ (12,5\textsuperscript{e}) : classe $[7,2;7,4[$, $N_{\text{ant}}=2$, $n=8$.
$Q_1^\beta \approx 7,2 + \frac{12,5-2}{8} \times 0,2 = \textbf{7,46}$.
$Q_3$ (37,5\textsuperscript{e}) : classe $[7,6;7,8[$, $N_{\text{ant}}=25$, $n=15$.
$Q_3^\beta \approx 7,6 + \frac{37,5-25}{15} \times 0,2 = \textbf{7,77}$.

\textbf{Histogrammes :} (À tracer : Alpha centré sur 7,5, symétrique ; Beta centré sur 7,6, étalé vers la droite).

\textbf{Choix et conclusion argumentée :}
\begin{quote}
\emph{« Le fournisseur \textbf{Alpha} est à privilégier. Sa résistance moyenne est plus faible (7,50 vs 7,59 m$\Omega$/m), ce qui indique une meilleure conductivité. Sa dispersion est également plus faible ($s_\alpha \approx 0,22$ vs $s_\beta \approx 0,28$ ; $CV_\alpha \approx 2,9\%$ vs $CV_\beta \approx 3,7\%$), gage de régularité. Concernant la spécification client ($R < 7,8$ m$\Omega$/m) : chez Alpha, le 3\textsuperscript{e} quartile $Q_3^\alpha = 7,65$ est bien en deçà du seuil, et seulement 5 câbles (10\%) se trouvent dans la classe $[7,8;8,0[$ (potentiellement hors specs). Chez Beta, $Q_3^\beta = 7,77$ frôle le seuil et 10 câbles (20\%) sont dans la classe $[7,8;8,0[$. Le risque de non-conformité est deux fois plus élevé chez Beta. Alpha offre donc la meilleure performance moyenne, la meilleure homogénéité et le meilleur respect des specs. »}
\end{quote}
\end{corrige}

\newpage

\coursec{Activité d'intégration}

\coursec{Leçon N05 : Séries statistiques regroupées en classes de même amplitude}

\courssub{Objectifs de l'activité}

À l'issue de cette activité d'intégration, l'élève sera capable de :
\begin{itemize}
\item regrouper une série statistique brute en \cle{classes de même amplitude} et construire le tableau complet (effectifs, effectifs cumulés croissants, fréquences, fréquences cumulées croissantes) ;
\item calculer les \cle{caractéristiques de position} (moyenne, classe modale, médiane par interpolation linéaire) et de \cle{dispersion} (variance, écart-type) d'une série regroupée ;
\item tracer un \cle{histogramme} et le \cle{polygone des effectifs cumulés croissants}, puis déterminer graphiquement la médiane ;
\item \cle{rédiger et justifier} une conclusion argumentée à destination d'un partenaire économique (une coopérative), en mobilisant le vocabulaire statistique correct.
\end{itemize}

\courssub{Situation}

La coopérative \emph{« Les Maraîchers Réunis de Yaoundé »} souhaite acheter du gombo en gros au marché de Mokolo pour approvisionner des restaurants de la ville. Afin de fixer un \textbf{prix d'achat groupé} équitable, elle a missionné des élèves de Première technique pour mener une enquête. Pendant une matinée, les élèves ont relevé le prix de vente au kilogramme (en FCFA) du gombo auprès de \textbf{60 vendeuses} du marché. Voici les données brutes :

\begin{center}
\begin{tabularx}{\linewidth}{|X|X|X|X|X|X|}
\hline
\rowcolor{popblueL} 250 & 380 & 420 & 510 & 640 & 330 \\
\hline 460 & 290 & 550 & 410 & 370 & 480 \\
\hline 620 & 440 & 350 & 500 & 430 & 260 \\
\hline 390 & 470 & 580 & 320 & 450 & 410 \\
\hline 530 & 400 & 360 & 490 & 280 & 440 \\
\hline 610 & 420 & 340 & 560 & 460 & 300 \\
\hline 430 & 520 & 450 & 380 & 650 & 470 \\
\hline 410 & 350 & 590 & 440 & 400 & 270 \\
\hline 480 & 430 & 540 & 390 & 460 & 420 \\
\hline 360 & 500 & 450 & 680 & 410 & 440 \\
\hline
\end{tabularx}
\end{center}

La coopérative pose deux questions aux élèves :
\begin{enumerate}
\item Quel \textbf{prix moyen} au kilogramme peut-on proposer comme prix d'achat groupé ?
\item Le marché est-il \textbf{homogène} (prix peu dispersés, faible écart-type) ou \textbf{très contrasté} (forte dispersion) ? Qu'en conclure pour la négociation ?
\end{enumerate}

\courssub{Consignes}

Travail en groupes de 4 élèves. Chaque groupe rédigera un rapport.

\textbf{Partie A : organisation des données.}
\begin{enumerate}
\item Vérifier que toutes les données sont comprises entre 200 et 700 FCFA. Regrouper les données en 5 classes d'amplitude 100 FCFA : $[200\,;300[$, $[300\,;400[$, $[400\,;500[$, $[500\,;600[$, $[600\,;700]$.
\item Construire le tableau statistique complet : effectifs $n_i$, effectifs cumulés croissants (ECC), fréquences $f_i$ en pourcentage, fréquences cumulées croissantes (FCC).
\end{enumerate}

\textbf{Partie B : indicateurs.}
\begin{enumerate}
\setcounter{enumi}{2}
\item Calculer le centre $x_i$ de chaque classe, puis la moyenne $\bar{x}$ de la série.
\item Déterminer la classe modale et interpréter.
\item Calculer la médiane $M_e$ par interpolation linéaire.
\item Calculer la variance $V$ puis l'écart-type $\sigma$. Interpréter la dispersion.
\end{enumerate}

\textbf{Partie C : représentations graphiques.}
\begin{enumerate}
\setcounter{enumi}{6}
\item Tracer l'histogramme des effectifs sur papier millimétré.
\item Tracer le polygone des ECC et vérifier graphiquement la valeur de la médiane.
\end{enumerate}

\textbf{Partie D : rapport à la coopérative.}
\begin{enumerate}
\setcounter{enumi}{8}
\item Rédiger un court rapport (10 à 15 lignes) répondant aux deux questions de la coopérative, en justifiant chaque affirmation par un indicateur calculé.
\end{enumerate}

\courssub{Ressources à mobiliser}

\begin{aretenir}[Formules utiles — série regroupée en classes]
\begin{itemize}
\item Centre d'une classe : $x_i=\dfrac{\text{borne inf}+\text{borne sup}}{2}$.
\item Fréquence : $f_i=\dfrac{n_i}{N}$ ; en pourcentage : $f_i\times 100$.
\item Moyenne : $\bar{x}=\dfrac{1}{N}\displaystyle\sum n_i x_i$.
\item Variance : $V=\dfrac{1}{N}\displaystyle\sum n_i x_i^2-\bar{x}^2$ ; écart-type : $\sigma=\sqrt{V}$.
\item Médiane par interpolation linéaire : si $M_e\in[a\,;b[$, alors
\[ M_e=a+(b-a)\times\dfrac{\dfrac{N}{2}-\text{ECC avant la classe}}{n_{\text{classe médiane}}}. \]
\item La médiane se lit graphiquement : abscisse du point du polygone des ECC d'ordonnée $N/2$.
\end{itemize}
\end{aretenir}

\courssub{Démarche et corrigé commenté}

\begin{exoresolu}[Enquête au marché de Mokolo — corrigé complet]
Regrouper les 60 prix en 5 classes, calculer les indicateurs, tracer les graphiques et rédiger le rapport.
\tcblower
\textbf{Étape 1 : tableau statistique complet.} En dépouillant les 60 données, on obtient les effectifs suivants :

\begin{center}
\begin{tabular}{|c|c|c|c|c|}
\hline
\rowcolor{popblueL} Classe (FCFA) & $n_i$ & ECC & $f_i$ (en \%) & FCC (en \%) \\
\hline $[200\,;300[$ & 5 & 5 & 8,3 & 8,3 \\
\hline $[300\,;400[$ & 13 & 18 & 21,7 & 30,0 \\
\hline $[400\,;500[$ & 27 & 45 & 45,0 & 75,0 \\
\hline $[500\,;600[$ & 10 & 55 & 16,7 & 91,7 \\
\hline $[600\,;700]$ & 5 & 60 & 8,3 & 100 \\
\hline Total & 60 & & 100 & \\
\hline
\end{tabular}
\end{center}

Vérification : $5+13+27+10+5=60$ ; la dernière FCC vaut bien $100\,\%$. La classe modale est $[400\,;500[$ (effectif maximal 27).

\textbf{Étape 2 : moyenne.} Les centres sont $x_i=250\,;350\,;450\,;550\,;650$.
\[ \bar{x}=\frac{5\times250+13\times350+27\times450+10\times550+5\times650}{60}=\frac{26700}{60}=445\ \text{FCFA}. \]

\textbf{Étape 3 : classe modale.} L'effectif maximal est $27$ : la classe modale est $[400\,;500[$. C'est dans cette tranche de prix que se situent le plus de vendeuses (45\,\% d'entre elles).

\textbf{Étape 4 : médiane par interpolation linéaire.} $\dfrac{N}{2}=30$. D'après les ECC, $18<30\leq 45$ : la classe médiane est $[400\,;500[$.
\[ M_e=400+100\times\frac{30-18}{27}=400+100\times\frac{12}{27}=400+\frac{400}{9}\approx 444{,}4\ \text{FCFA}. \]
Au moins la moitié des vendeuses vendent le kilogramme à moins de $444$ FCFA environ.

\textbf{Étape 5 : variance et écart-type.}
\begin{align*}
\sum n_i x_i^2 &= 5\times250^2+13\times350^2+27\times450^2+10\times550^2+5\times650^2 \\
&= 5\times62500 + 13\times122500 + 27\times202500 + 10\times302500 + 5\times422500 \\
&= 312500 + 1592500 + 5467500 + 3025000 + 2112500 = 12\,510\,000,\\
V &= \frac{12\,510\,000}{60}-(445)^2 = 208\,500 - 198\,025 = 10\,475,\\
\sigma &= \sqrt{V}\approx 102{,}3\ \text{FCFA}.
\end{align*}
L'écart-type ($\approx 102$ FCFA) représente environ $23\,\%$ de la moyenne : la dispersion est notable, le marché est contrasté.

\textbf{Étape 6 : histogramme et polygone des ECC.} Les classes ayant la même amplitude, la hauteur de chaque rectangle est proportionnelle à l'effectif.

\begin{popfigure}
\begin{tikzpicture}[x=0.018cm,y=0.22cm]
\draw[->] (150,0) -- (730,0) node[right]{\footnotesize prix (FCFA)};
\draw[->] (150,0) -- (150,30) node[above]{\footnotesize effectif};
\foreach \x/\h in {200/5,300/13,400/27,500/10,600/5}{
  \fill[popblue!45] (\x,0) rectangle (\x+100,\h);
  \draw[popdark] (\x,0) rectangle (\x+100,\h);}
\foreach \x in {200,300,400,500,600,700} \draw (\x,0) node[below]{\footnotesize \x};
\foreach \y in {5,10,15,20,25} \draw (150,\y) node[left]{\footnotesize \y};
\end{tikzpicture}
\legende{Histogramme des effectifs (prix du kg de gombo au marché de Mokolo).}
\end{popfigure}

\begin{popfigure}
\begin{tikzpicture}[x=0.016cm,y=0.075cm]
\draw[->] (150,0) -- (740,0) node[right]{\footnotesize prix (FCFA)};
\draw[->] (150,0) -- (150,66) node[above]{\footnotesize ECC};
\draw[poporange,very thick] plot coordinates {(200,0) (300,5) (400,18) (500,45) (600,55) (700,60)};
\foreach \p in {(200,0),(300,5),(400,18),(500,45),(600,55),(700,60)} \fill[popdark] \p circle (1.6);
\draw[dashed,popgreen,thick] (150,30) -- (444.4,30) -- (444.4,0);
\draw (444.4,0) node[below]{\footnotesize $M_e\approx444$};
\draw (150,30) node[left]{\footnotesize 30};
\foreach \x in {200,300,400,500,600,700} \draw (\x,0) node[below]{\footnotesize \x};
\end{tikzpicture}
\legende{Polygone des ECC : la médiane se lit à l'ordonnée $N/2=30$.}
\end{popfigure}

La lecture graphique ($M_e\approx 444$ FCFA) confirme le calcul par interpolation.

\textbf{Étape 7 : rapport à la coopérative (modèle de rédaction).}

\emph{« Notre enquête porte sur 60 vendeuses du marché de Mokolo. Le prix moyen du kilogramme de gombo est de \textbf{445 FCFA}, et la moitié des vendeuses le vendent à moins de \textbf{444 FCFA} (médiane). La tranche de prix la plus fréquente est $[400\,;500[$ FCFA (45\,\% des vendeuses). Nous proposons donc un prix d'achat groupé de l'ordre de \textbf{440 à 450 FCFA} le kilogramme. Cependant, l'écart-type est élevé (environ 102 FCFA, soit 23\,\% de la moyenne) : le marché est \textbf{contrasté}, avec des prix allant de 250 à 680 FCFA. La coopérative devra donc négocier vendeuse par vendeuse ou prévoir une marge, car un prix unique risque d'être refusé par les vendeuses les plus chères. »}
\end{exoresolu}

\begin{attention}[Erreurs fréquentes à éviter]
\begin{itemize}
\item Ne pas confondre ECC et fréquences cumulées : l'un est un nombre de vendeuses, l'autre un pourcentage.
\item Pour la médiane, utiliser $\dfrac{N}{2}=30$ et non $\dfrac{N+1}{2}$ dans l'interpolation sur classes.
\item Ne pas oublier de prendre les \textbf{centres} des classes pour la moyenne et la variance.
\item Dans la formule de la variance, c'est $\dfrac{\sum n_i x_i^2}{N}-\bar{x}^2$, et non $\left(\dfrac{\sum n_i x_i}{N}\right)^2$.
\item Vérifier que la somme des effectifs retrouvés est bien égale à l'effectif total $N=60$.
\end{itemize}
\end{attention}

\courssub{Bilan de l'activité}

Cette activité a permis de mobiliser, sur une situation authentique de la vie économique camerounaise, l'ensemble des savoir-faire de la leçon :
\begin{itemize}
\item le \cle{regroupement en classes} de même amplitude et la construction rigoureuse du tableau complet (effectifs, ECC, fréquences, FCC) ;
\item le calcul et surtout l'\cle{interprétation} des indicateurs : la moyenne ($445$ FCFA) et la médiane ($\approx 444$ FCFA), très proches, indiquent une répartition assez symétrique autour du prix central, tandis que l'écart-type élevé révèle un marché contrasté ;
\item le passage du \cle{calcul au graphique} : l'histogramme visualise la répartition (pic central marqué) et le polygone des ECC permet de retrouver la médiane, ce qui renforce la cohérence des résultats ;
\item la \cle{rédaction d'une conclusion justifiée} : les élèves apprennent qu'un indicateur statistique n'a de valeur que s'il est interprété et argumenté dans une décision concrète (fixation d'un prix d'achat groupé).
\end{itemize}
La mise en commune a montré que des groupes disposant des mêmes données peuvent rédiger des conclusions légèrement différentes : c'est la \emph{qualité de la justification} qui fait la valeur du rapport, compétence essentielle au Probatoire et au Baccalauréat technique.

\newpage

\coursec{Méthodes et exercices résolus}

\coursec{Séries statistiques regroupées en classes de même amplitude}

\courssub{Regroupement en classes, effectifs, effectifs cumulés et fréquences}

\begin{methode}[Dresser le tableau statistique complet d'une série regroupée en classes]
Pour traiter une série regroupée en classes de même amplitude $[a_i\,;\,a_{i+1}[$ :
\begin{enumerate}
\item Vérifier que les classes sont \textbf{disjointes} et \textbf{recouvrent} toutes les données ; calculer l'\cle{amplitude} commune $a = a_{i+1}-a_i$.
\item Compter l'\cle{effectif} $n_i$ de chaque classe et vérifier que $\sum n_i = N$ (effectif total).
\item Calculer les \cle{effectifs cumulés croissants} : $N_1=n_1$, puis $N_k = N_{k-1}+n_k$. Le dernier doit valoir $N$.
\item Calculer les \cle{fréquences} : $f_i=\dfrac{n_i}{N}$ (en décimal ou en pourcentage). Vérifier que $\sum f_i = 1$ (ou $100\,\%$).
\item Calculer les \cle{fréquences cumulées croissantes} : $F_k = F_{k-1}+f_k$. La dernière vaut $1$ (ou $100\,\%$).
\end{enumerate}
\textbf{Réflexe :} toujours ajouter une colonne « Total » et vérifier les trois contrôles : $\sum n_i = N$, $\sum f_i = 1$, dernier cumul $=N$ (ou $1$).
\end{methode}

\begin{exoresolu}[Contrôle de pièces dans un atelier de Douala]
Un atelier d'usinage contrôle le diamètre (en mm) de $50$ pièces. Les résultats sont regroupés ainsi :
\begin{center}
\begin{tabular}{|l|ccccc|}
\hline
\rowcolor{popblueL} Diamètre (mm) & $[19,5\,;\,19,7[$ & $[19,7\,;\,19,9[$ & $[19,9\,;\,20,1[$ & $[20,1\,;\,20,3[$ & $[20,3\,;\,20,5[$ \\
\hline
Effectif & 4 & 9 & 20 & 12 & 5 \\
\hline
\end{tabular}
\end{center}
Dresser le tableau complet avec effectifs cumulés croissants, fréquences et fréquences cumulées croissantes (en \%).
\tcblower
\textbf{Étape 1 : effectif total.} $N = 4+9+20+12+5 = 50$. Le contrôle est bon.
\smallskip

\textbf{Étape 2 : effectifs cumulés croissants.}
\begin{align*}
N_1 &= 4 \\
N_2 &= 4+9 = 13 \\
N_3 &= 13+20 = 33 \\
N_4 &= 33+12 = 45 \\
N_5 &= 45+5 = 50 = N \quad \text{(contrôle validé)}
\end{align*}

\textbf{Étape 3 : fréquences} $f_i=\dfrac{n_i}{50}$ :
\[ f_1=\frac{4}{50}=0{,}08 \;;\quad f_2=\frac{9}{50}=0{,}18 \;;\quad f_3=\frac{20}{50}=0{,}40 \;;\quad f_4=\frac{12}{50}=0{,}24 \;;\quad f_5=\frac{5}{50}=0{,}10. \]
Contrôle : $0{,}08+0{,}18+0{,}40+0{,}24+0{,}10 = 1$. 

\textbf{Étape 4 : fréquences cumulées} (en \%) : $8$ ; $8+18=26$ ; $26+40=66$ ; $66+24=90$ ; $90+10=100$.

\textbf{Tableau final :}
\begin{center}
\begin{tabular}{|l|c|c|c|c|}
\hline
\rowcolor{popblueL} Classe & $n_i$ & $N_i$ cumulés & $f_i$ (\%) & $F_i$ cumulées (\%) \\
\hline
$[19,5\,;\,19,7[$ & 4 & 4 & 8 & 8 \\
\hline
$[19,7\,;\,19,9[$ & 9 & 13 & 18 & 26 \\
\hline
$[19,9\,;\,20,1[$ & 20 & 33 & 40 & 66 \\
\hline
$[20,1\,;\,20,3[$ & 12 & 45 & 24 & 90 \\
\hline
$[20,3\,;\,20,5[$ & 5 & 50 & 10 & 100 \\
\hline
\rowcolor{popcream} Total & 50 & --- & 100 & --- \\
\hline
\end{tabular}
\end{center}
\textbf{Conclusion :} le tableau est cohérent (cumuls terminaux égaux à $N=50$ et $100\,\%$). On lit par exemple que $66\,\%$ des pièces ont un diamètre inférieur à $20{,}1$ mm.
\end{exoresolu}

\courssub{Caractéristiques de position : moyenne, classe modale, médiane}

\begin{methode}[Moyenne d'une série regroupée en classes]
\begin{enumerate}
\item Calculer le \cle{centre} de chaque classe : $c_i=\dfrac{a_i+a_{i+1}}{2}$.
\item Appliquer la formule de la moyenne :
\[ \bar{x}=\dfrac{\sum n_i c_i}{N}=\dfrac{n_1c_1+n_2c_2+\cdots+n_kc_k}{N}. \]
\item Conclure par une phrase avec l'\textbf{unité} de la grandeur.
\end{enumerate}
\textbf{Réflexe :} dresser une colonne $c_i$ et une colonne $n_i c_i$ dans le tableau, puis sommer la dernière.
\end{methode}

\begin{exoresolu}[Chiffre d'affaires journalier d'une quincaillerie]
Une quincaillerie de Yaoundé relève son chiffre d'affaires journalier (en milliers de FCFA) sur $40$ jours :
\begin{center}
\begin{tabular}{|l|cccc|}
\hline
\rowcolor{popblueL} Classe & $[50\,;\,70[$ & $[70\,;\,90[$ & $[90\,;\,110[$ & $[110\,;\,130[$ \\
\hline
Effectif $n_i$ & 6 & 14 & 12 & 8 \\
\hline
\end{tabular}
\end{center}
Calculer le chiffre d'affaires journalier moyen.
\tcblower
On complète le tableau avec les centres $c_i$ et les produits $n_ic_i$ :
\begin{center}
\begin{tabular}{|l|c|c|c|}
\hline
\rowcolor{popblueL} Classe & $n_i$ & $c_i$ & $n_i c_i$ \\
\hline
$[50\,;\,70[$ & 6 & 60 & 360 \\
\hline
$[70\,;\,90[$ & 14 & 80 & 1120 \\
\hline
$[90\,;\,110[$ & 12 & 100 & 1200 \\
\hline
$[110\,;\,130[$ & 8 & 120 & 960 \\
\hline
\rowcolor{popcream} Total & 40 & --- & 3640 \\
\hline
\end{tabular}
\end{center}
Détail des centres : $c_1=\dfrac{50+70}{2}=60$ ; $c_2=\dfrac{70+90}{2}=80$ ; $c_3=\dfrac{90+110}{2}=100$ ; $c_4=\dfrac{110+130}{2}=120$.
\smallskip

Application de la formule :
\[ \bar{x}=\dfrac{\sum n_ic_i}{N}=\dfrac{3640}{40}=91. \]
\textbf{Conclusion :} le chiffre d'affaires journalier moyen de la quincaillerie est de $91$ milliers de FCFA, soit $\num{91000}$ FCFA.
\end{exoresolu}

\begin{methode}[Déterminer la classe modale et la médiane]
\textbf{Classe modale :}
\begin{enumerate}
\item Repérer l'effectif $n_i$ le plus grand : la classe correspondante est la \cle{classe modale}.
\item Rédiger : « la classe modale est $[a\,;\,b[$ » (on ne donne pas une valeur unique mais une classe).
\end{enumerate}
\textbf{Médiane (par interpolation linéaire) :}
\begin{enumerate}
\item Calculer $\dfrac{N}{2}$ et repérer, dans la colonne des effectifs cumulés croissants, la \cle{classe médiane} $[a\,;\,b[$ : la première classe dont l'effectif cumulé atteint ou dépasse $\dfrac{N}{2}$.
\item Noter $N_{\text{avant}}$ l'effectif cumulé de la classe précédente, $n$ l'effectif de la classe médiane, $a$ son amplitude.
\item Appliquer la formule d'\cle{interpolation linéaire} :
\[ M_e = a_{\text{borne inf}} + \dfrac{\frac{N}{2}-N_{\text{avant}}}{n}\times a. \]
\end{enumerate}
\end{methode}

\begin{exoresolu}[Âges des apprentis d'un centre de formation]
Les $60$ apprentis d'un centre de formation technique de Bafoussam sont répartis selon l'âge :
\begin{center}
\begin{tabular}{|l|cccc|}
\hline
\rowcolor{popblueL} Âge (ans) & $[15\,;\,17[$ & $[17\,;\,19[$ & $[19\,;\,21[$ & $[21\,;\,23[$ \\
\hline
Effectif & 10 & 26 & 16 & 8 \\
\hline
\end{tabular}
\end{center}
\begin{enumerate}
\item Déterminer la classe modale.
\item Déterminer la médiane par interpolation linéaire et interpréter le résultat.
\end{enumerate}
\tcblower
\textbf{1. Classe modale.} L'effectif maximal est $26$, atteint pour la classe $[17\,;\,19[$. La classe modale est donc $[17\,;\,19[$ : la tranche d'âge la plus fréquente est celle des 17 à 19 ans.
\smallskip

\textbf{2. Médiane.} On calcule d'abord les effectifs cumulés croissants :
\begin{center}
\begin{tabular}{|l|c|c|}
\hline
\rowcolor{popblueL} Classe & $n_i$ & $N_i$ cumulés \\
\hline
$[15\,;\,17[$ & 10 & 10 \\
\hline
$[17\,;\,19[$ & 26 & 36 \\
\hline
$[19\,;\,21[$ & 16 & 52 \\
\hline
$[21\,;\,23[$ & 8 & 60 \\
\hline
\end{tabular}
\end{center}
On a $\dfrac{N}{2}=\dfrac{60}{2}=30$. Le premier effectif cumulé qui atteint ou dépasse $30$ est $36$ : la classe médiane est donc $[17\,;\,19[$.
\smallskip

On applique l'interpolation linéaire avec : borne inférieure $17$, $N_{\text{avant}}=10$, effectif de la classe $n=26$, amplitude $a=2$ :
\[ M_e = 17 + \dfrac{30-10}{26}\times 2 = 17 + \dfrac{40}{26} \approx 17 + 1{,}54 \approx 18{,}54. \]
\textbf{Conclusion :} la médiane vaut environ $18{,}5$ ans. Interprétation : au moins la moitié des apprentis du centre ont moins de $18{,}5$ ans environ.
\end{exoresolu}

\courssub{Caractéristiques de dispersion : étendue, variance et écart-type}

\begin{methode}[Calculer variance et écart-type d'une série en classes]
\begin{enumerate}
\item Calculer les centres $c_i$ et la moyenne $\bar{x}$.
\item Pour chaque classe, calculer l'écart $c_i-\bar{x}$, puis son carré $(c_i-\bar{x})^2$, puis le produit $n_i(c_i-\bar{x})^2$.
\item La \cle{variance} est :
\[ V=\dfrac{\sum n_i(c_i-\bar{x})^2}{N}. \]
\item L'\cle{écart-type} est $\sigma=\sqrt{V}$. Ne pas oublier la racine carrée et l'unité.
\item L'\cle{étendue} se lit directement : $e = $ borne supérieure de la dernière classe $-$ borne inférieure de la première classe.
\end{enumerate}
\textbf{Interprétation :} plus $\sigma$ est petit, plus la série est homogène (valeurs concentrées autour de la moyenne).
\end{methode}

\begin{exoresolu}[Comparaison de deux machines en atelier]
Deux machines A et B produisent des tiges de longueur nominale $100$ mm. On relève les longueurs (en mm) de $40$ tiges par machine. Pour la machine A :
\begin{center}
\begin{tabular}{|l|cccc|}
\hline
\rowcolor{popblueL} Classe & $[98\,;\,99[$ & $[99\,;\,100[$ & $[100\,;\,101[$ & $[101\,;\,102[$ \\
\hline
Effectif & 4 & 12 & 16 & 8 \\
\hline
\end{tabular}
\end{center}
\begin{enumerate}
\item Calculer la moyenne, la variance et l'écart-type de la série A.
\item Sachant que la machine B donne $\bar{x}_B=100{,}05$ mm et $\sigma_B=1{,}9$ mm, quelle machine est la plus régulière ? Justifier.
\end{enumerate}
\tcblower
\textbf{1.} Tableau de calcul (centres, produits pour la moyenne, écarts au carré pondérés) :
\begin{center}
\begin{tabular}{|l|c|c|c|c|c|}
\hline
\rowcolor{popblueL} Classe & $n_i$ & $c_i$ & $n_ic_i$ & $(c_i-\bar{x})^2$ & $n_i(c_i-\bar{x})^2$ \\
\hline
$[98\,;\,99[$ & 4 & 98{,}5 & 394 & 2{,}89 & 11{,}56 \\
\hline
$[99\,;\,100[$ & 12 & 99{,}5 & 1194 & 0{,}49 & 5{,}88 \\
\hline
$[100\,;\,101[$ & 16 & 100{,}5 & 1608 & 0{,}09 & 1{,}44 \\
\hline
$[101\,;\,102[$ & 8 & 101{,}5 & 812 & 1{,}69 & 13{,}52 \\
\hline
\rowcolor{popcream} Total & 40 & --- & 4008 & --- & 32{,}40 \\
\hline
\end{tabular}
\end{center}
Moyenne : $\bar{x}=\dfrac{4008}{40}=100{,}2$ mm.
\smallskip

Détail des écarts à la moyenne :
$c_1-\bar{x}=98{,}5-100{,}2=-1{,}7 \Rightarrow (-1{,}7)^2=2{,}89 \Rightarrow 4\times 2{,}89=11{,}56$ ;
$c_2-\bar{x}=99{,}5-100{,}2=-0{,}7 \Rightarrow 0{,}49 \Rightarrow 12\times 0{,}49=5{,}88$ ;
$c_3-\bar{x}=100{,}5-100{,}2=0{,}3 \Rightarrow 0{,}09 \Rightarrow 16\times 0{,}09=1{,}44$ ;
$c_4-\bar{x}=101{,}5-100{,}2=1{,}3 \Rightarrow 1{,}69 \Rightarrow 8\times 1{,}69=13{,}52$.
\smallskip

Somme des écarts au carré pondérés : $11{,}56+5{,}88+1{,}44+13{,}52=32{,}40$.
D'où :
\[ V=\dfrac{32{,}40}{40}=0{,}81 \qquad\text{et}\qquad \sigma=\sqrt{0{,}81}=0{,}9\ \text{mm}. \]
\textbf{2. Comparaison.} Les deux moyennes sont proches de $100$ mm ($\bar{x}_A=100{,}2$, $\bar{x}_B=100{,}05$), mais $\sigma_A=0{,}9 < \sigma_B=1{,}9$ mm.
\smallskip

\textbf{Conclusion :} la machine A est la plus régulière car son écart-type est le plus petit : ses productions sont plus concentrées autour de la moyenne, donc plus homogènes.
\end{exoresolu}

\courssub{Histogramme et polygone des fréquences}

\begin{methode}[Tracer un histogramme (classes de même amplitude)]
\begin{enumerate}
\item Les classes étant de \textbf{même amplitude}, la hauteur de chaque rectangle est proportionnelle à l'\textbf{effectif} (ou à la fréquence) de la classe.
\item Choisir les échelles : en abscisse les bornes des classes, en ordonnée les effectifs.
\item Tracer des rectangles \textbf{juxtaposés} (sans espace), de bases égales, de hauteurs $n_i$.
\item Titrer les axes avec les \textbf{unités}.
\end{enumerate}
\textbf{Lecture :} l'histogramme permet de repérer visuellement la classe modale (rectangle le plus haut) et la symétrie éventuelle de la distribution.
\end{methode}

\begin{methode}[Tracer le polygone des fréquences]
\begin{enumerate}
\item Sur le même repère que l'histogramme (ou à part), repérer les centres des classes $c_i$ en abscisse et les effectifs (ou fréquences) $n_i$ en ordonnée.
\item Tracer les points $M_i(c_i\,;\,n_i)$.
\item Relier les points consécutifs par des segments.
\item Pour fermer le polygone, ajouter une classe fictive de part et d'autre avec effectif nul (ou relier aux axes aux bornes extrêmes selon la convention).
\end{enumerate}
\end{methode}

\begin{exoresolu}[Histogramme et polygone des salaires dans une PME]
Les salaires mensuels (en milliers de FCFA) des $30$ employés d'une PME de Buea sont :
\begin{center}
\begin{tabular}{|l|cccc|}
\hline
\rowcolor{popblueL} Salaire & $[50\,;\,100[$ & $[100\,;\,150[$ & $[150\,;\,200[$ & $[200\,;\,250[$ \\
\hline
Effectif & 5 & 12 & 9 & 4 \\
\hline
\end{tabular}
\end{center}
\begin{enumerate}
\item Construire l'histogramme et le polygone des fréquences de cette série.
\item Quel pourcentage d'employés gagne moins de $\num{150000}$ FCFA ?
\end{enumerate}
\tcblower
\textbf{1.} Les classes ont toutes la même amplitude ($50$ milliers de FCFA) : les hauteurs des rectangles sont les effectifs $5$, $12$, $9$, $4$. Les centres des classes sont $75$, $125$, $175$, $225$.
\begin{popfigure}
\begin{center}
\begin{tikzpicture}[x=0.022cm,y=0.28cm]
% Axes
\draw[->,thick] (40,0) -- (270,0) node[right]{\small salaire (milliers FCFA)};
\draw[->,thick] (50,0) -- (50,15) node[above]{\small effectif};
% Graduations abscisses
\foreach \x in {50,100,150,200,250} \draw (\x,0) node[below]{\small \x};
% Graduations ordonnées
\foreach \y in {4,8,12} \draw (50,\y) node[left]{\small \y};
% Histogramme (rectangles)
\foreach \x/\h in {50/5,100/12,150/9,200/4}
  \fill[popblue!45,draw=popblue,thick] (\x,0) rectangle (\x+50,\h);
% Polygone des fréquences (centres : 75, 125, 175, 225)
\draw[poporange, thick, mark=*, mark size=2pt] plot coordinates {(75,5) (125,12) (175,9) (225,4)};
% Classes fictives pour fermer le polygone (optionnel, ici on s'arrête aux centres)
\end{tikzpicture}
\end{center}
\legende{Histogramme et polygone des fréquences des salaires mensuels (30 employés)}
\end{popfigure}

\textbf{2.} Les employés gagnant moins de $\num{150000}$ FCFA sont ceux des deux premières classes : $5+12=17$ employés.
\[ f=\dfrac{17}{30}\approx 0{,}567 \quad\text{soit environ } 56{,}7\,\%. \]
\textbf{Conclusion :} environ $56{,}7\,\%$ des employés de la PME gagnent moins de $\num{150000}$ FCFA par mois.
\end{exoresolu}

\courssub{Exploiter les cumuls pour répondre à une question concrète}

\begin{methode}[Répondre à une question du type « combien de... au moins / moins de... »]
\begin{enumerate}
\item Traduire la question en termes de \textbf{borne de classe} : « moins de $b$ » $\to$ effectif cumulé croissant jusqu'à $b$ ; « au moins $a$ » $\to$ $N$ moins l'effectif cumulé jusqu'à $a$.
\item Lire directement dans la colonne des effectifs (ou fréquences) cumulés.
\item Conclure par une phrase en contexte (métier, unité).
\end{enumerate}
\textbf{Réflexe :} si la borne demandée tombe à l'intérieur d'une classe, on utilise l'\cle{interpolation linéaire}, comme pour la médiane.
\end{methode}

\begin{exoresolu}[Notes au Probatoire blanc d'un lycée technique]
Les notes de mathématiques de $80$ élèves de Première technique sont regroupées ainsi :
\begin{center}
\begin{tabular}{|l|ccccc|}
\hline
\rowcolor{popblueL} Note & $[0\,;\,4[$ & $[4\,;\,8[$ & $[8\,;\,12[$ & $[12\,;\,16[$ & $[16\,;\,20]$ \\
\hline
Effectif & 6 & 18 & 30 & 20 & 6 \\
\hline
\end{tabular}
\end{center}
\begin{enumerate}
\item Combien d'élèves ont obtenu une note inférieure à $8$ ? Quelle proportion cela représente-t-il ?
\item Estimer, par interpolation linéaire, le nombre d'élèves ayant obtenu la moyenne (note $\geq 10$).
\end{enumerate}
\tcblower
\textbf{1.} Effectifs cumulés croissants : $6$ ; $6+18=24$ ; $24+30=54$ ; $54+20=74$ ; $74+6=80$.
\smallskip

Les élèves de note inférieure à $8$ sont dans les deux premières classes : effectif cumulé $=24$ élèves.
\[ f=\dfrac{24}{80}=0{,}30 \quad\text{soit } 30\,\%. \]
\textbf{2.} La note $10$ se situe dans la classe $[8\,;\,12[$, d'effectif $30$ et d'amplitude $4$. Par interpolation linéaire, l'effectif cumulé jusqu'à la note $10$ est :
\[ N(10) = 24 + \dfrac{10-8}{4}\times 30 = 24 + \dfrac{2}{4}\times 30 = 24+15 = 39. \]
Le nombre d'élèves ayant au moins $10$ est donc :
\[ 80 - 39 = 41. \]
\textbf{Conclusion :} $24$ élèves ($30\,\%$) ont une note inférieure à $8$, et environ $41$ élèves, soit un peu plus de la moitié de l'effectif, ont obtenu la moyenne au Probatoire blanc.
\end{exoresolu}

\begin{aretenir}[Synthèse : formules et repères pour le Probatoire]
\textbf{Tableau de base :} Classes $[a_i;a_{i+1}[$ de même amplitude $a$ | Effectifs $n_i$ | Centres $c_i=\frac{a_i+a_{i+1}}{2}$ | $n_ic_i$ | Effectifs cumulés $N_i$ | Fréquences $f_i=n_i/N$ | Fréquences cumulées $F_i$.
\smallskip

\textbf{Moyenne :} $\bar{x}=\frac{\sum n_ic_i}{N}$.
\smallskip

\textbf{Classe modale :} classe d'effectif maximal.
\smallskip

\textbf{Médiane (interpolation) :} $M_e = a_{\text{inf}} + \frac{\frac{N}{2}-N_{\text{avant}}}{n_{\text{médiane}}}\times a$.
\smallskip

\textbf{Variance :} $V=\frac{\sum n_i(c_i-\bar{x})^2}{N}$ \quad Écart-type : $\sigma=\sqrt{V}$ \quad Étendue : $e = \text{max} - \text{min}$.
\smallskip

\textbf{Histogramme :} rectangles juxtaposés, hauteur = effectif (ou fréquence), base = amplitude.
\textbf{Polygone :} points $(c_i; n_i)$ reliés.
\smallskip

\textbf{Interrogation des cumuls :} « $< b$ » $\to$ lire $N$ ou $F$ à la borne $b$ ; « $\ge a$ » $\to$ $N - N(a)$ ou $1-F(a)$.
\end{aretenir}

\begin{attention}[Les 5 erreurs qui coûtent des points au Probatoire]
\begin{enumerate}
\item \textbf{Oublier les centres de classes.} Pour la moyenne et la variance d'une série en classes, on utilise les centres $c_i=\dfrac{a_i+a_{i+1}}{2}$, jamais les bornes. Calculer $\bar{x}$ avec les bornes inférieures est une erreur éliminatoire.
\item \textbf{Confondre effectif et effectif cumulé.} La médiane et les questions « moins de... » se lisent sur les \emph{cumuls}. Vérifier systématiquement que le dernier effectif cumulé égale $N$ et que la dernière fréquence cumulée vaut $1$ (ou $100\,\%$).
\item \textbf{Donner l'écart-type sans racine carrée (ou la variance à sa place).} $\sigma=\sqrt{V}$ : annoncer $\sigma=V$ est faux. Et toujours préciser l'unité de l'écart-type (celle des données).
\item \textbf{Espacer les rectangles de l'histogramme.} Pour des classes de même amplitude, les rectangles sont \emph{juxtaposés}, sans espace, et les axes doivent être gradués et titrés avec les unités. Un diagramme en barres espacées n'est pas un histogramme.
\item \textbf{Conclure sans phrase ni unité.} Une moyenne de $91$ ne veut rien dire seule : il faut écrire « le chiffre d'affaires moyen est de $\num{91000}$ FCFA ». Toute réponse à une situation concrète doit se terminer par une phrase d'interprétation en contexte.
\end{enumerate}
\end{attention}

\newpage

\coursec{Exercices}

\courssub{Je vérifie mes connaissances}

\begin{exercice}[QCM : Vocabulaire et formules de base]
On considère une série statistique regroupée en classes de même amplitude $a$. Parmi les affirmations suivantes, une seule est exacte. Laquelle ?
\begin{enumerate}
    \item L'amplitude d'une classe $[a_i; a_{i+1}[$ est $a_{i+1} - a_i + 1$.
    \item Le centre (ou milieu) de la classe $[a_i; a_{i+1}[$ est $\dfrac{a_i + a_{i+1}}{2}$.
    \item La fréquence cumulée croissante de la dernière classe est toujours égale à $0$.
    \item L'écart-type d'une série regroupée se calcule avec la formule exacte $\sqrt{\frac{\sum n_i (x_i - \bar{x})^2}{\sum n_i}}$ où $x_i$ sont les valeurs exactes des données.
\end{enumerate}
\end{exercice}

\begin{exercice}[Vrai ou Faux ? Justifiez]
Déterminez si les affirmations suivantes sont vraies ou fausses. Justifiez votre réponse par un calcul ou une propriété du cours.
\begin{enumerate}
    \item Dans un histogramme représentant une série regroupée en classes de même amplitude, la surface de chaque rectangle est proportionnelle à l'effectif de la classe correspondante.
    \item La médiane d'une série regroupée appartient toujours à la classe médiane.
    \item La formule de König-Huygens $V = \frac{\sum n_i c_i^2}{\sum n_i} - \bar{x}^2$ donne la variance exacte de la série initiale (non regroupée).
    \item Si l'effectif total $N$ est pair, la médiane est la moyenne des deux valeurs centrales de la série ordonnée.
\end{enumerate}
\end{exercice}

\begin{exercice}[Définitions et repérage]
On donne le tableau suivant des hauteurs (en cm) de 40 plants de maïs dans une exploitation agricole de la région de l'Ouest :

\begin{center}
\begin{tabular}{|c|c|c|c|c|c|}\hline
Classes (cm) & $[100;120[$ & $[120;140[$ & $[140;160[$ & $[160;180[$ & $[180;200[$ \\ \hline
Effectifs $n_i$ & 5 & 12 & 15 & 6 & 2 \\ \hline
\end{tabular}
\end{center}

\begin{enumerate}
    \item Quelle est l'amplitude des classes ? Donnez le centre de la classe $[140;160[$.
    \item Identifiez la classe modale. Quelle est sa fréquence (en \%) ?
    \item Déterminez la classe médiane (celle qui contient la médiane).
    \item Écrivez la formule de l'interpolation linéaire pour estimer la médiane dans la classe médiane. Précisez le sens de chaque lettre.
\end{enumerate}
\end{exercice}

\begin{exercice}[Calculs directs sur effectifs et fréquences]
Une équipe de maintenance relève les durées d'intervention (en minutes) sur 50 pannes électriques. Le tableau incomplet ci-dessous résulte de l'enquête.

\begin{center}
\begin{tabular}{|c|c|c|c|c|c|}\hline
Classes (min) & $[10;20[$ & $[20;30[$ & $[30;40[$ & $[40;50[$ & $[50;60[$ \\ \hline
Effectifs $n_i$ & 8 & 15 & $a$ & 10 & $b$ \\ \hline
Effectifs cumulés croissants & 8 & 23 & 35 & $c$ & 50 \\ \hline
Fréquences $f_i$ (en \%) & 16 & $d$ & 24 & 20 & $e$ \\ \hline
\end{tabular}
\end{center}

\begin{enumerate}
    \item Retrouvez les valeurs manquantes $a, b, c, d, e$.
    \item Calculez la moyenne approximative de cette série (arrondie au dixième).
    \item Calculez l'écart-type approximatif (arrondi au centième).
\end{enumerate}
\end{exercice}

\courssub{Je m'entraîne}

\begin{exercice}[Sacs de maïs au marché de Yaoundé]
Un contrôleur mesure la masse (en kg) de 80 sacs de maïs. Les résultats sont regroupés dans le tableau ci-dessous.

\begin{center}
\begin{tabular}{|c|c|c|c|c|c|c|}\hline
Classes (kg) & $[40;45[$ & $[45;50[$ & $[50;55[$ & $[55;60[$ & $[60;65[$ & $[65;70[$ \\ \hline
Effectifs $n_i$ & 5 & 12 & 23 & 20 & 15 & 5 \\ \hline
\end{tabular}
\end{center}

\begin{enumerate}
    \item Complétez le tableau en ajoutant les lignes : effectifs cumulés croissants, fréquences (en \%), fréquences cumulées croissantes (en \%), centres de classes $c_i$.
    \item Calculez la moyenne $\bar{x}$ et l'écart-type $\sigma$ de cette série (arrondis au centième).
    \item Déterminez la médiane par interpolation linéaire (arrondie au centième).
    \item Construisez l'histogramme de cette série (échelle : 1 cm pour 5 kg en abscisse, 1 cm pour 2 sacs en ordonnée).
\end{enumerate}
\end{exercice}

\begin{exercice}[Histogramme incomplet : résistances de câbles]
Dans un atelier d'électricité, on teste la résistance (en ohms) de 100 câbles. L'histogramme ci-dessous est partiellement tracé. Les classes ont une amplitude de 2 ohms. Les deux derniers rectangles (classes $[12;14[$ et $[14;16[$) sont manquants.
On sait que la moyenne de la série est $\bar{x} = 9,8$ ohms.

\begin{center}
\begin{popfigure}
\begin{tikzpicture}[scale=0.7]
\begin{axis}[
    width=\linewidth, height=8cm,
    xlabel={Résistance (ohms)}, ylabel={Effectifs},
    xmin=0, xmax=16, ymin=0, ymax=25,
    xtick={0,2,4,6,8,10,12,14,16},
    ytick={0,5,10,15,20,25},
    ymajorgrids=true, grid style=dashed,
    axis lines=middle,
    enlargelimits={abs=0.5},
]
% Rectangles connus
\addplot[popblue!50, fill=popblue!30, draw=popblue] coordinates {(2,0) (2,10) (4,10) (4,0)} \closedcycle;
\addplot[popblue!50, fill=popblue!30, draw=popblue] coordinates {(4,0) (4,18) (6,18) (6,0)} \closedcycle;
\addplot[popblue!50, fill=popblue!30, draw=popblue] coordinates {(6,0) (6,22) (8,22) (8,0)} \closedcycle;
\addplot[popblue!50, fill=popblue!30, draw=popblue] coordinates {(8,0) (8,20) (10,20) (10,0)} \closedcycle;
\addplot[popblue!50, fill=popblue!30, draw=popblue] coordinates {(10,0) (10,12) (12,12) (12,0)} \closedcycle;
% Rectangles manquants (contours pointillés)
\draw[poporange, dashed, thick] (12,0) -- (12,0) -- (14,0) -- (14,0) -- cycle; % Placeholder
\draw[poporange, dashed, thick] (14,0) -- (14,0) -- (16,0) -- (16,0) -- cycle; % Placeholder
\node[poporange] at (13,2) {?};
\node[poporange] at (15,2) {?};
\end{axis}
\end{tikzpicture}
\legende{Histogramme des résistances (deux dernières classes manquantes)}
\end{popfigure}
\end{center}

\begin{enumerate}
    \item En déduisant les effectifs des classes connues, calculez la somme des effectifs des deux classes manquantes.
    \item Soit $x$ l'effectif de la classe $[12;14[$ et $y$ celui de $[14;16[$. Établissez un système de deux équations à deux inconnues en utilisant l'effectif total et la moyenne.
    \item Résolvez le système et donnez les effectifs $x$ et $y$.
    \item Calculez alors l'écart-type de la série complète.
\end{enumerate}
\end{exercice}

\begin{exercice}[Comparaison de deux marques d'ampoules — Entreprise de Douala]
Une entreprise de Douala souhaite choisir un fournisseur pour des ampoules. Elle teste la durée de vie (en centaines d'heures) de deux marques A et B. Résultats regroupés :

\emph{Marque A (échantillon de 60 ampoules) :}
\begin{center}
\begin{tabular}{|c|c|c|c|c|c|}\hline
Classes & $[5;7[$ & $[7;9[$ & $[9;11[$ & $[11;13[$ & $[13;15[$ \\ \hline
Effectifs & 5 & 15 & 20 & 15 & 5 \\ \hline
\end{tabular}
\end{center}

\emph{Marque B (échantillon de 60 ampoules) :}
\begin{center}
\begin{tabular}{|c|c|c|c|c|c|}\hline
Classes & $[5;7[$ & $[7;9[$ & $[9;11[$ & $[11;13[$ & $[13;15[$ \\ \hline
Effectifs & 10 & 10 & 20 & 10 & 10 \\ \hline
\end{tabular}
\end{center}

\begin{enumerate}
    \item Pour chaque marque, calculez la moyenne et l'écart-type (arrondis au centième).
    \item Tracez les deux histogrammes sur le même repère (utilisez des hachures ou couleurs différentes).
    \item Comparez la dispersion des deux séries. Laquelle est la plus homogène ?
    \item Rédigez une conclusion argumentée (3 à 4 lignes) pour conseiller l'entreprise dans son choix de fournisseur.
\end{enumerate}
\end{exercice}

\begin{exercice}[Polygone des fréquences cumulées — Notes au Probatoire blanc]
Les notes sur 20 d'une classe de 50 élèves de Première Technique à un devoir de mathématiques sont regroupées ainsi :

\begin{center}
\begin{tabular}{|c|c|c|c|c|c|c|}\hline
Classes & $[0;4[$ & $[4;8[$ & $[8;12[$ & $[12;16[$ & $[16;20[$ \\ \hline
Effectifs & 5 & 10 & 18 & 12 & 5 \\ \hline
\end{tabular}
\end{center}

\begin{enumerate}
    \item Construisez le tableau des effectifs cumulés croissants et des fréquences cumulées croissantes (en \%). Donnez les coordonnées des points du polygone des fréquences cumulées croissantes.
    \item Tracez le polygone des fréquences cumulées croissantes (échelle : 1 cm pour 2 notes en abscisse, 1 cm pour 10\% en ordonnée).
    \item Par lecture graphique, estimez la médiane, le premier quartile $Q_1$ et le troisième quartile $Q_3$.
    \item Vérifiez ces valeurs par interpolation linéaire (formules). Commentez la précision de la lecture graphique.
\end{enumerate}
\end{exercice}

\courssub{Je me prépare au Probatoire}

\begin{exercice}[Étude complète : Motos-taxis de Bafoussam — Type Probatoire]
Une enquête est menée auprès de 50 conducteurs de motos-taxis à Bafoussam sur la distance quotidienne parcourue (en km). Les résultats sont les suivants :

\begin{center}
\begin{tabular}{|c|c|c|c|c|c|c|c|}\hline
Classes (km) & $[50;70[$ & $[70;90[$ & $[90;110[$ & $[110;130[$ & $[130;150[$ & $[150;170[$ \\ \hline
Effectifs & 4 & 8 & 15 & 12 & 7 & 4 \\ \hline
\end{tabular}
\end{center}

\begin{enumerate}
    \item Complétez le tableau statistique : centres de classes $c_i$, effectifs cumulés croissants, fréquences (en \%), fréquences cumulées croissantes (en \%).
    \item Calculez la moyenne $\bar{x}$ et l'écart-type $\sigma$ de cette série (arrondis au dixième de km).
    \item Déterminez la médiane par interpolation linéaire. Interprétez ce résultat dans le contexte.
    \item Calculez le premier quartile $Q_1$ et le troisième quartile $Q_3$. En déduire l'intervalle interquartile $IIQ$.
    \item Construisez l'histogramme de cette série (échelle : 1 cm pour 20 km en abscisse, 1 cm pour 2 conducteurs en ordonnée).
    \item \textbf{Rédaction :} Rédigez un rapport statistique de 5 lignes maximum à l'intention de la mairie de Bafoussam, synthétisant les résultats principaux (moyenne, dispersion, médiane) et proposant une recommandation pour l'aménagement des stations de repos.
\end{enumerate}
\end{exercice}

\begin{exercice}[Utilisation de la formule de König-Huygens — Pièces mécaniques]
Dans une usine de mécanique de précision, on contrôle le diamètre (en mm) de 200 pièces. La série est regroupée en classes de même amplitude. On note $c_i$ les centres de classes et $n_i$ les effectifs.
On connaît les résultats partiels suivants :
\[
\sum n_i = 200, \quad \sum n_i c_i = 5000, \quad \sum n_i c_i^2 = 125\,400, \quad \text{Variance } V = 4,5.
\]

\begin{enumerate}
    \item Vérifiez que la moyenne $\bar{x}$ est égale à 25 mm.
    \item En utilisant la formule de König-Huygens ($V = \frac{\sum n_i c_i^2}{\sum n_i} - \bar{x}^2$), vérifiez la cohérence de la variance donnée. Si elle n'est pas cohérente, calculez la variance réelle.
    \item Sachant que l'amplitude des classes est $a = 2$ mm, déterminez les bornes de la première classe si la dernière classe est $[29;31[$ et qu'il y a 6 classes au total.
    \item On suppose que la série suit approximativement une distribution normale. Quel pourcentage de pièces a un diamètre compris entre $\bar{x} - \sigma$ et $\bar{x} + \sigma$ ? (Arrondi à l'entier).
\end{enumerate}
\end{exercice}

\begin{exercice}[Synthèse : Salaires dans une PME de Bamenda — Décision de gestion]
La direction d'une PME de Bamenda étudie les salaires mensuels (en milliers de FCFA) de ses 80 employés pour préparer le budget d'augmentation. Le tableau ci-dessous donne la répartition :

\begin{center}
\begin{tabular}{|c|c|c|c|c|c|c|}\hline
Classes (kFCFA) & $[150;200[$ & $[200;250[$ & $[250;300[$ & $[300;350[$ & $[350;400[$ & $[400;450[$ \\ \hline
Effectifs & 10 & 18 & 22 & 15 & 10 & 5 \\ \hline
\end{tabular}
\end{center}

\begin{enumerate}
    \item Déterminez la classe modale. Estimez le mode par la formule d'interpolation (ou graphiquement) : $M_o \approx x_{m} + \frac{n_m - n_{m-1}}{(n_m - n_{m-1}) + (n_m - n_{m+1})} \times a$. Interprétez.
    \item Calculez la moyenne $\bar{x}$ et l'écart-type $\sigma$ (en kFCFA, arrondis au centième).
    \item Déterminez la médiane par interpolation linéaire. Comparez la médiane et la moyenne. Que pouvez-vous dire de l'asymétrie de la distribution ?
    \item Calculez le coefficient de variation $CV = \frac{\sigma}{\bar{x}} \times 100$. La dispersion est-elle forte ou faible ?
    \item La direction dispose d'une enveloppe de 1\,200 kFCFA pour les augmentations. Elle souhaite augmenter de 5\% les salaires inférieurs à la médiane et de 2\% les salaires supérieurs à la médiane. En utilisant la médiane calculée, estimez le coût total de cette mesure. L'enveloppe est-elle suffisante ? Justifiez par un calcul.
\end{enumerate}
\end{exercice}

\newpage

\coursec{Corrigés des exercices}

\coursec{Séries statistiques regroupées en classes de même amplitude}

\courssub{Regroupement en classes et effectifs}

\begin{definition}[Classe, amplitude, centre]
Pour une série statistique quantitative continue ou discrète à effectif total $N$ important, on regroupe les valeurs en $k$ \cle{classes} d'\cle{amplitude} constante $a$ :
\[ [a_1;a_2[,\ [a_2;a_3[,\ \dots,\ [a_k;a_{k+1}[ \]
avec $a = a_{i+1}-a_i$ (identique pour tout $i$).
Le \cle{centre de la classe} $i$ (ou \cle{valeur représentative}) est :
\[ c_i = \frac{a_i + a_{i+1}}{2}. \]
L'\cle{effectif} $n_i$ est le nombre de données appartenant à la classe $i$. On a $\sum_{i=1}^k n_i = N$.
\end{definition}

\begin{propriete}[Choix des classes]
\begin{itemize}
    \item Les classes sont emboîtées, exhaustives et mutuellement exclusives.
    \item L'amplitude $a$ se choisit souvent arrondie (5, 10, 20, \dots) pour faciliter la lecture.
    \item Le nombre de classes $k$ est généralement compris entre 5 et 15 (règle de Sturges : $k \approx 1 + 3,3 \log_{10} N$).
\end{itemize}
\end{propriete}

\begin{exemplebox}[Construction d'un tableau de regroupement]
On mesure le diamètre (en mm) de 50 pièces tournées. Valeurs min 19,2 ; max 30,8.
On choisit $a = 2$ mm $\Rightarrow$ 6 classes : $[19;21[$, $[21;23[$, $[23;25[$, $[25;27[$, $[27;29[$, $[29;31[$.
On compte les effectifs $n_i$ pour chaque classe.
\end{exemplebox}

\courssub{Effectifs cumulés et fréquences}

\begin{definition}[Effectifs cumulés, fréquences]
\begin{itemize}
    \item \cle{Eff. cumulé croissant} $N_i = \sum_{j=1}^i n_j$ (nb de valeurs $< a_{i+1}$).
    \item \cle{Eff. cumulé décroissant} $N'_i = \sum_{j=i}^k n_j$ (nb de valeurs $\ge a_i$).
    \item \cle{Fréquence} $f_i = \dfrac{n_i}{N}$ (souvent en \%). $\sum f_i = 1$.
    \item \cle{Fréquence cumulée croissante} $F_i = \dfrac{N_i}{N}$.
\end{itemize}
\end{definition}

\begin{propriete}
La dernière fréquence cumulée croissante vaut toujours $1$ (ou $100\,\%$).
\end{propriete}

\begin{methode}[Remplir un tableau statistique complet]
\begin{enumerate}
    \item Déterminer $a$, les bornes $a_i$ et les centres $c_i$.
    \item Compter les effectifs $n_i$ (recherche des valeurs manquantes par soustraction si besoin).
    \item Calculer les effectifs cumulés croissants $N_i$ (addition successive).
    \item Calculer les fréquences $f_i = n_i/N$ et fréquences cumulées $F_i = N_i/N$.
\end{enumerate}
\end{methode}

\courssub{Caractéristiques de position}

\begin{propriete}[Moyenne approchée]
On remplace chaque donnée par le centre de sa classe :
\[ \bar{x} = \frac{\sum_{i=1}^k n_i c_i}{\sum_{i=1}^k n_i} = \frac{\sum n_i c_i}{N}. \]
\end{propriete}

\begin{propriete}[Médiane par interpolation linéaire]
\emph{Exigible au Probatoire.}
\begin{enumerate}
    \item Trouver la \cle{classe médiane} : première classe dont l'effectif cumulé croissant $N_m \ge \dfrac{N}{2}$.
    \item Soit $x_m$ sa borne inférieure, $n_m$ son effectif, $N_{m-1}$ le cumulé précédent.
    \item Médiane estimée :
    \[ M_e \approx x_m + \frac{\dfrac{N}{2} - N_{m-1}}{n_m} \times a. \]
\end{enumerate}
\end{propriete}

\begin{propriete}[Mode par interpolation linéaire (classe modale)]
La \cle{classe modale} est celle de plus grand effectif $n_{mo}$.
Soient $d_1 = n_{mo} - n_{mo-1}$ et $d_2 = n_{mo} - n_{mo+1}$ (effectifs des classes voisines, $0$ si classe extrême).
\[ M_o \approx x_{mo} + \frac{d_1}{d_1+d_2} \times a. \]
\end{propriete}

\begin{propriete}[Quartiles par interpolation linéaire]
Même principe que la médiane :
\begin{itemize}
    \item $Q_1$ (premier quartile) : classe où $N_i \ge \dfrac{N}{4}$.
    \item $Q_3$ (troisième quartile) : classe où $N_i \ge \dfrac{3N}{4}$.
    \item \cle{Écart interquartile} $IIQ = Q_3 - Q_1$ (dispersion du centre de la série).
\end{itemize}
\end{propriete}

\begin{attention}[Pièges fréquents]
\begin{itemize}
    \item Ne pas confondre \textbf{valeur exacte} (série non regroupée) et \textbf{estimation par centres/interpolation} (série regroupée).
    \item Pour la médiane : vérifier la condition « \emph{première classe dont l'effectif cumulé atteint ou dépasse $N/2$} ».
    \item L'interpolation suppose une répartition \emph{uniforme} des données à l'intérieur de la classe.
\end{itemize}
\end{attention}

\courssub{Caractéristiques de dispersion}

\begin{propriete}[Étendue]
\[ E = \text{borne sup max} - \text{borne inf min} = a_{k+1} - a_1. \]
\end{propriete}

\begin{propriete}[Variance et écart-type (formule de König-Huygens)]
\emph{Formule à connaître et à écrire explicitement.}
\[ V = \frac{\sum n_i c_i^2}{N} - \bar{x}^2 \qquad \sigma = \sqrt{V}. \]
\emph{Remarque :} comme on utilise les centres $c_i$, $V$ et $\sigma$ sont des \emph{valeurs approchées} de la variance et de l'écart-type réels.
\end{propriete}

\begin{propriete}[Coefficient de variation]
\[ CV = \frac{\sigma}{\bar{x}} \times 100 \quad (\text{en \%}). \]
Permet de comparer la dispersion relative de deux séries de moyennes différentes.
\begin{itemize}
    \item $CV < 15\,\%$ : dispersion faible (série homogène).
    \item $15\,\% \le CV \le 30\,\%$ : dispersion modérée.
    \item $CV > 30\,\%$ : dispersion forte (série hétérogène).
\end{itemize}
\end{propriete}

\courssub{Histogramme et polygone des fréquences cumulées}

\begin{propriete}[Histogramme (classes même amplitude)]
\begin{itemize}
    \item Axe des abscisses : bornes des classes (échelle réelle).
    \item Axe des ordonnées : effectifs $n_i$ (ou fréquences $f_i$).
    \item Largeur des rectangles = amplitude $a$.
    \item Hauteur = effectif $n_i$ (car $a$ constant $\Rightarrow$ aire $\propto n_i$).
\end{itemize}
\end{propriete}

\begin{propriete}[Polygone des fréquences cumulées croissantes]
On trace les points $\big(a_{i+1},\ F_i\big)$ pour $i=1\dots k$, en ajoutant l'origine $\big(a_1,\ 0\big)$.
On relie les points par des segments.
\emph{Utilité :} lecture graphique approchée de $M_e$ (ordonnée $0,5$), $Q_1$ ($0,25$), $Q_3$ ($0,75$).
\end{propriete}

\begin{experience}[Construction d'un histogramme avec TikZ]
\begin{popfigure}
\begin{tikzpicture}[scale=0.8]
\begin{axis}[
    width=\linewidth, height=6cm,
    xlabel={Variable $X$}, ylabel={Effectifs},
    xmin=0, xmax=30, ymin=0, ymax=16,
    xtick={0,5,10,15,20,25,30}, ytick={0,5,10,15},
    ymajorgrids=true, grid style=dashed, axis lines=middle,
]
\addplot[fill=popblue!30, draw=popblue] coordinates {(0,0) (0,5) (5,5) (5,0)} \closedcycle;
\addplot[fill=popblue!30, draw=popblue] coordinates {(5,0) (5,12) (10,12) (10,0)} \closedcycle;
\addplot[fill=popblue!30, draw=popblue] coordinates {(10,0) (10,14) (15,14) (15,0)} \closedcycle;
\addplot[fill=popblue!30, draw=popblue] coordinates {(15,0) (15,8) (20,8) (20,0)} \closedcycle;
\addplot[fill=popblue!30, draw=popblue] coordinates {(20,0) (20,4) (25,4) (25,0)} \closedcycle;
\addplot[fill=popblue!30, draw=popblue] coordinates {(25,0) (25,2) (30,2) (30,0)} \closedcycle;
\end{axis}
\end{tikzpicture}
\legende{Exemple d'histogramme (classes d'amplitude 5)}
\end{popfigure}
\end{experience}

\courssub{Synthèse : démarche complète type Probatoire}

\begin{methode}[Résolution d'un problème complet]
\begin{enumerate}
    \item \textbf{Tableau :} Construire le tableau complet (classes, $n_i$, $c_i$, $N_i$, $f_i$, $F_i$). Vérifier $\sum n_i = N$, $\sum f_i = 1$.
    \item \textbf{Moyenne :} $\bar{x} = \frac{\sum n_i c_i}{N}$ (écrire la somme explicitement).
    \item \textbf{Écart-type :} $V = \frac{\sum n_i c_i^2}{N} - \bar{x}^2$, $\sigma = \sqrt{V}$ (écrire la formule de König-Huygens).
    \item \textbf{Médiane / Quartiles :} Identifier la classe (cumulé), écrire la formule d'interpolation, calculer, \textbf{interpréter en contexte}.
    \item \textbf{Mode :} Identifier classe modale, interpolation si demandée.
    \item \textbf{Représentation :} Histogramme (règle : aire = effectif) et/ou polygone cumulé.
    \item \textbf{Conclusion :} Réponse argumentée à la question posée (choix, décision, comparaison).
\end{enumerate}
\end{methode}

\begin{aretenir}[Vocabulaire Probatoire]
\begin{itemize}
    \item « \textbf{Justifier} » : écrire la formule, substituer les nombres, conclure.
    \item « \textbf{Interpréter} » : donner du sens au chiffre dans le contexte de l'exercice (ex: « la moitié des conducteurs parcourt moins de 107 km »).
    \item « \textbf{Estimer} » : calculer une valeur approchée (moyenne, médiane, mode) à partir des centres de classes.
\end{itemize}
\end{aretenir}

% ============================================================
% EXERCICES
% ============================================================

\courssub{Je vérifie mes connaissances}

\begin{corrige}[Exercice 1 — QCM : Vocabulaire et formules de base]
La bonne réponse est l'affirmation \textbf{b}.
\begin{itemize}
    \item \textbf{a) Fausse.} L'amplitude de la classe $[a_i;a_{i+1}[$ est $a = a_{i+1} - a_i$ (sans ajouter 1). Exemple : la classe $[100;120[$ a pour amplitude $120-100=20$.
    \item \textbf{b) Exacte.} Le centre de classe est $c_i = \dfrac{a_i+a_{i+1}}{2}$ : c'est la valeur qui remplace toutes les données de la classe dans les calculs approchés.
    \item \textbf{c) Fausse.} La fréquence cumulée croissante de la dernière classe vaut toujours $1$ (ou $100\,\%$), jamais $0$.
    \item \textbf{d) Fausse.} Dans une série regroupée, on ne connaît pas les valeurs exactes : on utilise les \cle{centres de classes} $c_i$. La formule $\sigma = \sqrt{\dfrac{\sum n_i(c_i-\bar{x})^2}{\sum n_i}}$ ne donne donc qu'une valeur \emph{approchée} de l'écart-type.
\end{itemize}
\end{corrige}

\begin{corrige}[Exercice 2 — Vrai ou Faux ? Justifiez]
\begin{enumerate}
    \item \textbf{Vrai.} L'amplitude $a$ étant la même pour toutes les classes, la hauteur de chaque rectangle est proportionnelle à l'effectif $n_i$ ; l'aire $a \times h_i$ est donc aussi proportionnelle à $n_i$.
    \item \textbf{Vrai.} Par définition, la \cle{classe médiane} est la première classe dont l'effectif cumulé croissant atteint ou dépasse $\dfrac{N}{2}$ : la médiane (valeur qui partage la série en deux effectifs égaux) appartient à cette classe.
    \item \textbf{Faux.} La formule de König-Huygens appliquée avec les centres $c_i$ ne donne que la variance \emph{approchée} de la série regroupée, car on a remplacé chaque donnée par le centre de sa classe. Elle ne coïncide avec la variance exacte que si toutes les données d'une classe sont égales au centre.
    \item \textbf{Faux} dans le cadre de cette leçon. Cette règle (moyenne des deux valeurs centrales) ne s'applique qu'à une série \emph{non regroupée}, dont on connaît les valeurs individuelles. Pour une série regroupée en classes, on estime la médiane par \cle{interpolation linéaire} dans la classe médiane.
\end{enumerate}
\end{corrige}

\begin{corrige}[Exercice 3 — Définitions et repérage]
\begin{enumerate}
    \item Amplitude : $a = 120-100 = 20$ cm (identique pour toutes les classes). Centre de $[140;160[$ : $c_3 = \dfrac{140+160}{2} = 150$ cm.
    \item La \cle{classe modale} est celle du plus grand effectif : $[140;160[$ avec $n_3 = 15$. Sa fréquence est $f_3 = \dfrac{15}{40} = 0{,}375$, soit $\mathbf{37{,}5\,\%}$.
    \item Effectifs cumulés croissants : $5$ ; $17$ ; $32$ ; $38$ ; $40$. On a $\dfrac{N}{2} = \dfrac{40}{2} = 20$. La première classe dont l'effectif cumulé atteint 20 est $[140;160[$ (cumulé 32) : c'est la \cle{classe médiane}.
    \item Formule d'interpolation linéaire :
    \[ M_e \approx x_m + \frac{\dfrac{N}{2} - N_{m-1}}{n_m}\times a \]
    où $x_m$ est la borne inférieure de la classe médiane ($140$), $N$ l'effectif total ($40$), $N_{m-1}$ l'effectif cumulé de la classe précédente ($17$), $n_m$ l'effectif de la classe médiane ($15$) et $a$ l'amplitude ($20$).
\end{enumerate}
\end{corrige}

\begin{corrige}[Exercice 4 — Calculs directs sur effectifs et fréquences]
\begin{enumerate}
    \item \textbf{Recherche des valeurs manquantes.}
    \begin{itemize}
        \item $a = 35 - 23 = 12$ (effectif cumulé moins cumulé précédent) ;
        \item $c = 35 + 10 = 45$ ;
        \item $b = 50 - 45 = 5$ ;
        \item $d = \dfrac{15}{50}\times 100 = 30\,\%$ ;
        \item $e = \dfrac{5}{50}\times 100 = 10\,\%$.
    \end{itemize}
    Vérification : $16+30+24+20+10 = 100\,\%$ et $8+15+12+10+5 = 50$. 
    \item \textbf{Moyenne.} Centres : $15, 25, 35, 45, 55$.
    \[ \sum n_i c_i = 8(15)+15(25)+12(35)+10(45)+5(55) = 120+375+420+450+275 = 1640. \]
    \[ \bar{x} = \frac{1640}{50} = \mathbf{32{,}8 \text{ min}}. \]
    \item \textbf{Écart-type} (formule de König-Huygens) :
    \[ \sum n_i c_i^2 = 8(225)+15(625)+12(1225)+10(2025)+5(3025) = 1800+9375+14700+20250+15125 = 61250. \]
    \[ V = \frac{61250}{50} - 32{,}8^2 = 1225 - 1075{,}84 = 149{,}16 \quad\Longrightarrow\quad \sigma = \sqrt{149{,}16} \approx \mathbf{12{,}21 \text{ min}}. \]
\end{enumerate}
\end{corrige}

\courssub{Je m'entraîne}

\begin{corrige}[Exercice 5 — Sacs de maïs au marché de Yaoundé]
\begin{enumerate}
    \item Tableau complété ($N = 80$) :
    \begin{center}
    \begin{tabular}{|c|c|c|c|c|c|c|}\hline
    Classes (kg) & $[40;45[$ & $[45;50[$ & $[50;55[$ & $[55;60[$ & $[60;65[$ & $[65;70[$ \\ \hline
    $n_i$ & 5 & 12 & 23 & 20 & 15 & 5 \\ \hline
    Cumulés croissants & 5 & 17 & 40 & 60 & 75 & 80 \\ \hline
    $f_i$ (\%) & 6,25 & 15 & 28,75 & 25 & 18,75 & 6,25 \\ \hline
    $f_i$ cumulées (\%) & 6,25 & 21,25 & 50 & 75 & 93,75 & 100 \\ \hline
    Centres $c_i$ & 42,5 & 47,5 & 52,5 & 57,5 & 62,5 & 67,5 \\ \hline
    \end{tabular}
    \end{center}
    \item \textbf{Moyenne et écart-type.}
    \begin{align*}
    \sum n_i c_i &= 5(42{,}5)+12(47{,}5)+23(52{,}5)+20(57{,}5)+15(62{,}5)+5(67{,}5)\\
    &= 212{,}5+570+1207{,}5+1150+937{,}5+337{,}5 = 4415.
    \end{align*}
    \[ \bar{x} = \frac{4415}{80} = 55{,}1875 \approx \mathbf{55{,}19 \text{ kg}}. \]
    \begin{align*}
    \sum n_i c_i^2 &= 9031{,}25+27075+63393{,}75+66125+58593{,}75+22781{,}25 = 247000.\\
    V &= \frac{247000}{80} - 55{,}1875^2 = 3087{,}5 - 3045{,}66 = 41{,}84.
    \end{align*}
    \[ \sigma = \sqrt{41{,}84} \approx \mathbf{6{,}47 \text{ kg}}. \]
    \item \textbf{Médiane.} $\dfrac{N}{2} = 40$ : classe médiane $[50;55[$ (cumulé 17 $\to$ 40).
    \[ M_e = 50 + \frac{40-17}{23}\times 5 = 50 + \frac{115}{23} = 50 + 5 = \mathbf{55{,}00 \text{ kg}}. \]
    La moitié des sacs pèse moins de 55 kg.
    \item \textbf{Histogramme :}
    \begin{center}
    \begin{popfigure}
    \begin{tikzpicture}[scale=0.85]
    \begin{axis}[
        width=\linewidth, height=7cm,
        xlabel={Masse (kg)}, ylabel={Effectifs},
        xmin=38, xmax=72, ymin=0, ymax=26,
        xtick={40,45,50,55,60,65,70},
        ytick={0,5,10,15,20,25},
        ymajorgrids=true, grid style=dashed,
        axis lines=middle,
    ]
    \addplot[fill=popblue!30, draw=popblue] coordinates {(40,0) (40,5) (45,5) (45,0)} \closedcycle;
    \addplot[fill=popblue!30, draw=popblue] coordinates {(45,0) (45,12) (50,12) (50,0)} \closedcycle;
    \addplot[fill=popblue!30, draw=popblue] coordinates {(50,0) (50,23) (55,23) (55,0)} \closedcycle;
    \addplot[fill=popblue!30, draw=popblue] coordinates {(55,0) (55,20) (60,20) (60,0)} \closedcycle;
    \addplot[fill=popblue!30, draw=popblue] coordinates {(60,0) (60,15) (65,15) (65,0)} \closedcycle;
    \addplot[fill=popblue!30, draw=popblue] coordinates {(65,0) (65,5) (70,5) (70,0)} \closedcycle;
    \end{axis}
    \end{tikzpicture}
    \legende{Histogramme des masses des 80 sacs de maïs}
    \end{popfigure}
    \end{center}
\end{enumerate}
\end{corrige}

\begin{corrige}[Exercice 6 — Histogramme incomplet : résistances de câbles]
\begin{enumerate}
    \item Effectifs lus sur l'histogramme : $10, 18, 22, 20, 12$, soit $82$ câbles. Il manque donc $100 - 82 = \mathbf{18}$ câbles répartis sur les deux dernières classes.
    \item Centres des classes : $3, 5, 7, 9, 11, 13, 15$. Somme connue :
    \[ \sum n_i c_i = 10(3)+18(5)+22(7)+20(9)+12(11) = 30+90+154+180+132 = 586. \]
    Avec $\bar{x} = 9{,}8$, il faudrait $\sum n_i c_i = 980$, d'où le système :
    \[ \begin{cases} x + y = 18 \\ 13x + 15y = 980 - 586 = 394 \end{cases} \]
    \item \textbf{Point de vigilance.} Ce système est \emph{impossible} : avec $x+y=18$, on a $13x+15y \leq 15\times 18 = 270 < 394$. Les données de l'énoncé (moyenne $9{,}8$) sont donc incohérentes avec l'histogramme. On corrige en prenant $\bar{x} = 8{,}4$ ohms (valeur compatible). Alors $\sum n_i c_i = 840$ et :
    \[ \begin{cases} x + y = 18 \\ 13x + 15y = 840 - 586 = 254 \end{cases} \]
    De la première équation $x = 18 - y$ ; en substituant : $13(18-y)+15y = 254 \Rightarrow 234 + 2y = 254 \Rightarrow y = 10$, puis $x = 8$.
    \[ \boxed{x = 8 \text{ et } y = 10} \]
    \item \textbf{Écart-type} (série complétée, $\bar{x} = 8{,}4$) :
    \begin{align*}
    \sum n_i c_i^2 &= 10(9)+18(25)+22(49)+20(81)+12(121)+8(169)+10(225)\\
    &= 90+450+1078+1620+1452+1352+2250 = 8292.
    \end{align*}
    \[ V = \frac{8292}{100} - 8{,}4^2 = 82{,}92 - 70{,}56 = 12{,}36 \quad\Longrightarrow\quad \sigma \approx \mathbf{3{,}52 \text{ ohms}}. \]
\end{enumerate}
\end{corrige}

\begin{corrige}[Exercice 7 — Comparaison de deux marques d'ampoules]
Centres de classes : $6, 8, 10, 12, 14$ (en centaines d'heures).
\begin{enumerate}
    \item \textbf{Marque A :}
    \[ \sum n_i c_i = 5(6)+15(8)+20(10)+15(12)+5(14) = 30+120+200+180+70 = 600 \Rightarrow \bar{x}_A = \frac{600}{60} = \mathbf{10{,}00}. \]
    \[ \sum n_i c_i^2 = 180+960+2000+2160+980 = 6280 \Rightarrow V_A = \frac{6280}{60}-100 = 4{,}67 \Rightarrow \sigma_A \approx \mathbf{2{,}16}. \]
    \textbf{Marque B :}
    \[ \sum n_i c_i = 10(6)+10(8)+20(10)+10(12)+10(14) = 60+80+200+120+140 = 600 \Rightarrow \bar{x}_B = \mathbf{10{,}00}. \]
    \[ \sum n_i c_i^2 = 360+640+2000+1440+1960 = 6400 \Rightarrow V_B = \frac{6400}{60}-100 = 6{,}67 \Rightarrow \sigma_B \approx \mathbf{2{,}58}. \]
    \item Histogrammes superposés :
    \begin{center}
    \begin{popfigure}
    \begin{tikzpicture}[scale=0.85]
    \begin{axis}[
        width=\linewidth, height=7cm,
        xlabel={Durée de vie (centaines d'heures)}, ylabel={Effectifs},
        xmin=4, xmax=16, ymin=0, ymax=24,
        xtick={5,7,9,11,13,15}, ytick={0,5,10,15,20},
        ymajorgrids=true, grid style=dashed, axis lines=middle,
    ]
    \addplot[fill=popblue!40, draw=popblue] coordinates {(5,0) (5,5) (7,5) (7,0)} \closedcycle;
    \addplot[fill=popblue!40, draw=popblue] coordinates {(7,0) (7,15) (9,15) (9,0)} \closedcycle;
    \addplot[fill=popblue!40, draw=popblue] coordinates {(9,0) (9,20) (11,20) (11,0)} \closedcycle;
    \addplot[fill=popblue!40, draw=popblue] coordinates {(11,0) (11,15) (13,15) (13,0)} \closedcycle;
    \addplot[fill=popblue!40, draw=popblue] coordinates {(13,0) (13,5) (15,5) (15,0)} \closedcycle;
    \addplot[fill=poporange!45, draw=poporange, opacity=0.7] coordinates {(5,0) (5,10) (7,10) (7,0)} \closedcycle;
    \addplot[fill=poporange!45, draw=poporange, opacity=0.7] coordinates {(7,0) (7,10) (9,10) (9,0)} \closedcycle;
    \addplot[fill=poporange!45, draw=poporange, opacity=0.7] coordinates {(13,0) (13,10) (15,10) (15,0)} \closedcycle;
    \node[popblue] at (6.2,20) {A};
    \node[poporange] at (14.5,17) {B};
    \end{axis}
    \end{tikzpicture}
    \legende{Histogrammes des marques A (bleu) et B (orange)}
    \end{popfigure}
    \end{center}
    \item Les deux marques ont la \emph{même moyenne} (1000 h), mais $\sigma_A = 2{,}16 < \sigma_B = 2{,}58$ : la marque A est plus concentrée autour de la moyenne, donc \textbf{plus homogène}.
    \item \textbf{Conclusion.} Les deux marques offrent la même durée de vie moyenne (environ 1000 heures). Cependant, la marque A présente un écart-type plus faible (2,16 contre 2,58) : ses ampoules ont des durées de vie plus régulières et plus prévisibles, avec moins de risque de tomber sur une ampoule de très courte durée. L'entreprise a donc intérêt à choisir le \textbf{fournisseur A}, à prix égal.
\end{enumerate}
\end{corrige}

\begin{corrige}[Exercice 8 — Polygone des fréquences cumulées — Notes au Probatoire blanc]
\begin{enumerate}
    \item Tableau des cumuls ($N = 50$) :
    \begin{center}
    \begin{tabular}{|c|c|c|c|c|c|}\hline
    Classes & $[0;4[$ & $[4;8[$ & $[8;12[$ & $[12;16[$ & $[16;20[$ \\ \hline
    $n_i$ & 5 & 10 & 18 & 12 & 5 \\ \hline
    Cumulés croissants & 5 & 15 & 33 & 45 & 50 \\ \hline
    Fréq. cumulées (\%) & 10 & 30 & 66 & 90 & 100 \\ \hline
    \end{tabular}
    \end{center}
    Points du polygone (borne supérieure de classe, fréquence cumulée) :
    \[ (0;0),\ (4;10),\ (8;30),\ (12;66),\ (16;90),\ (20;100). \]
    \item Polygone des fréquences cumulées croissantes :
    \begin{center}
    \begin{popfigure}
    \begin{tikzpicture}[scale=0.85]
    \begin{axis}[
        width=\linewidth, height=7.5cm,
        xlabel={Notes sur 20}, ylabel={Fréquences cumulées (\%)},
        xmin=0, xmax=20, ymin=0, ymax=105,
        xtick={0,2,4,6,8,10,12,14,16,18,20}, ytick={0,10,25,30,50,66,75,90,100},
        grid=both, grid style=dashed, axis lines=middle,
    ]
    \addplot[popcurve, mark=*, mark size=1.5pt] coordinates {(0,0) (4,10) (8,30) (12,66) (16,90) (20,100)};
    \addplot[poporange, dashed] coordinates {(0,50) (10.2,50)} ;
    \addplot[poporange, dashed] coordinates {(10.2,50) (10.2,0)} ;
    \addplot[popgreen, dashed] coordinates {(0,25) (7,25)} ;
    \addplot[popgreen, dashed] coordinates {(7,25) (7,0)} ;
    \addplot[poppurple, dashed] coordinates {(0,75) (13.5,75)} ;
    \addplot[poppurple, dashed] coordinates {(13.5,75) (13.5,0)} ;
    \end{axis}
    \end{tikzpicture}
    \legende{Polygone des fréquences cumulées croissantes et lectures graphiques}
    \end{popfigure}
    \end{center}
    \item \textbf{Lecture graphique :} la médiane est l'abscisse du point d'ordonnée 50\,\% : $M_e \approx 10$ ; $Q_1$ (ordonnée 25\,\%) $\approx 7$ ; $Q_3$ (ordonnée 75\,\%) $\approx 13{,}5$.
    \item \textbf{Vérification par interpolation linéaire.}
    \begin{itemize}
        \item Médiane : $\dfrac{N}{2}=25$, classe médiane $[8;12[$ (cumulé 15 $\to$ 33) :
        \[ M_e = 8 + \frac{25-15}{18}\times 4 = 8 + \frac{40}{18} \approx \mathbf{10{,}22}. \]
        \item $Q_1$ : $\dfrac{N}{4} = 12{,}5$, classe $[4;8[$ (cumulé 5 $\to$ 15) :
        \[ Q_1 = 4 + \frac{12{,}5-5}{10}\times 4 = 4 + 3 = \mathbf{7{,}00}. \]
        \item $Q_3$ : $\dfrac{3N}{4} = 37{,}5$, classe $[12;16[$ (cumulé 33 $\to$ 45) :
        \[ Q_3 = 12 + \frac{37{,}5-33}{12}\times 4 = 12 + 1{,}5 = \mathbf{13{,}50}. \]
    \end{itemize}
    \textbf{Commentaire.} Les lectures graphiques ($10$ ; $7$ ; $13{,}5$) sont très proches des valeurs calculées ($10{,}22$ ; $7$ ; $13{,}5$) : l'écart est inférieur à 0,3 point. La lecture graphique donne donc une bonne estimation, mais seule l'interpolation linéaire fournit une valeur précise et justifiée, exigée à l'écrit.
\end{enumerate}
\end{corrige}

\courssub{Je me prépare au Probatoire}

\begin{corrige}[Exercice 9 — Motos-taxis de Bafoussam — Type Probatoire]
\begin{enumerate}
    \item Tableau complété ($N = 50$) :
    \begin{center}
    \begin{tabular}{|c|c|c|c|c|c|c|}\hline
    Classes (km) & $[50;70[$ & $[70;90[$ & $[90;110[$ & $[110;130[$ & $[130;150[$ & $[150;170[$ \\ \hline
    $n_i$ & 4 & 8 & 15 & 12 & 7 & 4 \\ \hline
    Centres $c_i$ & 60 & 80 & 100 & 120 & 140 & 160 \\ \hline
    Cumulés croissants & 4 & 12 & 27 & 39 & 46 & 50 \\ \hline
    $f_i$ (\%) & 8 & 16 & 30 & 24 & 14 & 8 \\ \hline
    $f_i$ cumulées (\%) & 8 & 24 & 54 & 78 & 92 & 100 \\ \hline
    \end{tabular}
    \end{center}
    \item \textbf{Moyenne et écart-type.}
    \begin{align*}
    \sum n_i c_i &= 4(60)+8(80)+15(100)+12(120)+7(140)+4(160)\\
    &= 240+640+1500+1440+980+640 = 5440.
    \end{align*}
    \[ \bar{x} = \frac{5440}{50} = \mathbf{108{,}8 \text{ km}}. \]
    \begin{align*}
    \sum n_i c_i^2 &= 14400+51200+150000+172800+137200+102400 = 628000.\\
    V &= \frac{628000}{50} - 108{,}8^2 = 12560 - 11837{,}44 = 722{,}56.
    \end{align*}
    \[ \sigma = \sqrt{722{,}56} \approx \mathbf{26{,}9 \text{ km}}. \]
    \item \textbf{Médiane.} $\dfrac{N}{2} = 25$ : classe médiane $[90;110[$ (cumulé 12 $\to$ 27).
    \[ M_e = 90 + \frac{25-12}{15}\times 20 = 90 + \frac{260}{15} \approx \mathbf{107{,}3 \text{ km}}. \]
    \emph{Interprétation :} la moitié des conducteurs parcourt moins d'environ 107 km par jour, l'autre moitié plus de 107 km.
    \item \textbf{Quartiles.}
    \[ Q_1 : \frac{N}{4} = 12{,}5 \Rightarrow [70;90[ : \quad Q_1 = 70 + \frac{12{,}5-4}{8}\times 20 = 70 + 21{,}25 = \mathbf{91{,}25 \text{ km}}. \]
    \[ Q_3 : \frac{3N}{4} = 37{,}5 \Rightarrow [110;130[ : \quad Q_3 = 110 + \frac{37{,}5-27}{12}\times 20 = 110 + 17{,}5 = \mathbf{127{,}5 \text{ km}}. \]
    \[ IIQ = Q_3 - Q_1 = 127{,}5 - 91{,}25 = \mathbf{36{,}25 \text{ km}}. \]
    La moitié centrale des conducteurs parcourt entre 91,25 km et 127,5 km par jour.
    \item \textbf{Histogramme :}
    \begin{center}
    \begin{popfigure}
    \begin{tikzpicture}[scale=0.85]
    \begin{axis}[
        width=\linewidth, height=7cm,
        xlabel={Distance quotidienne (km)}, ylabel={Effectifs},
        xmin=45, xmax=175, ymin=0, ymax=17,
        xtick={50,70,90,110,130,150,170}, ytick={0,2,4,6,8,10,12,14,16},
        ymajorgrids=true, grid style=dashed, axis lines=middle,
    ]
    \addplot[fill=popgreen!30, draw=popgreen] coordinates {(50,0) (50,4) (70,4) (70,0)} \closedcycle;
    \addplot[fill=popgreen!30, draw=popgreen] coordinates {(70,0) (70,8) (90,8) (90,0)} \closedcycle;
    \addplot[fill=popgreen!30, draw=popgreen] coordinates {(90,0) (90,15) (110,15) (110,0)} \closedcycle;
    \addplot[fill=popgreen!30, draw=popgreen] coordinates {(110,0) (110,12) (130,12) (130,0)} \closedcycle;
    \addplot[fill=popgreen!30, draw=popgreen] coordinates {(130,0) (130,7) (150,7) (150,0)} \closedcycle;
    \addplot[fill=popgreen!30, draw=popgreen] coordinates {(150,0) (150,4) (170,4) (170,0)} \closedcycle;
    \end{axis}
    \end{tikzpicture}
    \legende{Histogramme des distances quotidiennes parcourues}
    \end{popfigure}
    \end{center}
    \item \textbf{Rapport (exemple de rédaction).} L'enquête menée auprès de 50 motos-taxis montre qu'un conducteur parcourt en moyenne 108,8 km par jour, avec un écart-type de 26,9 km : les distances sont donc assez dispersées. La médiane (107,3 km) confirme qu'un conducteur sur deux dépasse 107 km quotidiennement, et la moitié des conducteurs roule entre 91 et 128 km. La mairie a intérêt à implanter des stations de repos le long des axes les plus fréquentés, en visant une couverture adaptée à des journées de 90 à 130 km.
\end{enumerate}
\textbf{Barème indicatif (sur 10) :} tableau complété 2 pts ; moyenne 1 pt ; écart-type 1,5 pt ; médiane + interprétation 1,5 pt ; quartiles et $IIQ$ 1,5 pt ; histogramme 1,5 pt ; rapport rédigé 1 pt.\\
\textbf{Points de vigilance :} utiliser les \emph{centres de classes} dans tous les calculs ; citer la condition « première classe dont l'effectif cumulé atteint $\frac{N}{2}$ » pour la classe médiane ; ne pas oublier l'interprétation en contexte, exigée par le savoir-faire « rédiger et justifier ».
\end{corrige}

\begin{corrige}[Exercice 10 — Formule de König-Huygens — Pièces mécaniques]
\begin{enumerate}
    \item $\bar{x} = \dfrac{\sum n_i c_i}{\sum n_i} = \dfrac{5000}{200} = \mathbf{25 \text{ mm}}$. La moyenne annoncée est correcte.
    \item Formule de König-Huygens :
    \[ V = \frac{\sum n_i c_i^2}{\sum n_i} - \bar{x}^2 = \frac{125400}{200} - 25^2 = 627 - 625 = 2. \]
    La variance réelle est $V = 2$ (mm$^2$), et non $4{,}5$ : la valeur donnée est \textbf{incohérente}. L'écart-type réel est $\sigma = \sqrt{2} \approx 1{,}41$ mm.
    \item Six classes d'amplitude 2 mm, la dernière étant $[29;31[$. En remontant : $[27;29[$, $[25;27[$, $[23;25[$, $[21;23[$, $[19;21[$. La première classe est donc $\mathbf{[19;21[}$. (Cohérent avec la moyenne 25 mm, centre de la classe médiane $[23;25[$.)
    \item Pour une distribution approximativement normale, environ $68\,\%$ des valeurs appartiennent à l'intervalle $[\bar{x}-\sigma\,;\,\bar{x}+\sigma] = [23{,}59\,;\,26{,}41]$ mm. On estime donc qu'environ $\mathbf{68\,\%}$ des pièces (soit environ 136 pièces sur 200) ont un diamètre dans cet intervalle.
\end{enumerate}
\textbf{Barème indicatif (sur 5) :} moyenne 1 pt ; application de König-Huygens et conclusion 2 pts ; bornes des classes 1 pt ; pourcentage 1 pt.\\
\textbf{Points de vigilance :} bien écrire la formule avant de substituer ; ne pas confondre variance ($V = 2$) et écart-type ($\sigma = \sqrt{2}$) ; justifier la cohérence par un calcul, pas par une affirmation.
\end{corrige}

\begin{corrige}[Exercice 11 — Salaires dans une PME de Bamenda — Décision de gestion]
\begin{enumerate}
    \item \textbf{Classe modale :} $[250;300[$ (effectif maximal 22). Mode par interpolation :
    \[ M_o \approx 250 + \frac{22-18}{(22-18)+(22-15)}\times 50 = 250 + \frac{4}{11}\times 50 \approx 250 + 18{,}18 = \mathbf{268{,}18 \text{ kFCFA}}. \]
    \emph{Interprétation :} le salaire le plus fréquent est d'environ 268\,000 FCFA.
    \item \textbf{Moyenne et écart-type.} Centres : $175, 225, 275, 325, 375, 425$.
    \begin{align*}
    \sum n_i c_i &= 10(175)+18(225)+22(275)+15(325)+10(375)+5(425)\\
    &= 1750+4050+6050+4875+3750+2125 = 22600.
    \end{align*}
    \[ \bar{x} = \frac{22600}{80} = \mathbf{282{,}50 \text{ kFCFA}}. \]
    \begin{align*}
    \sum n_i c_i^2 &= 306250+911250+1663750+1584375+1406250+903125 = 6775000.\\
    V &= \frac{6775000}{80} - 282{,}5^2 = 84687{,}5 - 79806{,}25 = 4881{,}25.
    \end{align*}
    \[ \sigma = \sqrt{4881{,}25} \approx \mathbf{69{,}87 \text{ kFCFA}}. \]
    \item \textbf{Médiane.} Cumulés : $10, 28, 50, 65, 75, 80$. $\dfrac{N}{2} = 40$ : classe médiane $[250;300[$.
    \[ M_e = 250 + \frac{40-28}{22}\times 50 = 250 + \frac{600}{22} \approx \mathbf{277{,}27 \text{ kFCFA}}. \]
    On a $M_o \approx 268{,}18 < M_e \approx 277{,}27 < \bar{x} = 282{,}5$ : la moyenne est supérieure à la médiane, ce qui traduit une distribution \textbf{étalée vers les salaires élevés} (asymétrie à droite) : quelques salaires élevés tirent la moyenne vers le haut.
    \item \textbf{Coefficient de variation :}
    \[ CV = \frac{\sigma}{\bar{x}}\times 100 = \frac{69{,}87}{282{,}5}\times 100 \approx \mathbf{24{,}7\,\%}. \]
    La dispersion est \textbf{modérée} (ni faible ni forte) : les salaires sont assez hétérogènes autour de la moyenne.
    \item \textbf{Coût de la mesure.} On raisonne avec les centres de classes.
    \begin{itemize}
        \item \emph{Salaires inférieurs à la médiane} (277,27 kFCFA) : les classes $[150;200[$ et $[200;250[$ entières (28 employés), plus la fraction $\dfrac{277{,}27-250}{50} \approx 0{,}545$ de la classe $[250;300[$, soit environ $0{,}545 \times 22 \approx 12$ employés.
        Masse salariale concernée : $1750 + 4050 + 12 \times 275 = 9100$ kFCFA.
        Coût à 5\,\% : $9100 \times 0{,}05 = 455$ kFCFA.
        \item \emph{Salaires supérieurs à la médiane} : les 10 employés restants de $[250;300[$ et les classes suivantes :
        $10 \times 275 + 4875 + 3750 + 2125 = 2750 + 10750 = 13500$ kFCFA.
        Coût à 2\,\% : $13500 \times 0{,}02 = 270$ kFCFA.
        \item \emph{Coût total estimé :} $455 + 270 = \mathbf{725 \text{ kFCFA}}$.
    \end{itemize}
    Comme $725 < 1200$, \textbf{l'enveloppe de 1\,200 kFCFA est suffisante}, avec une marge d'environ 475 kFCFA.
\end{enumerate}
\textbf{Barème indicatif (sur 10) :} classe modale + mode 1,5 pt ; moyenne 1 pt ; écart-type 1,5 pt ; médiane + comparaison 2 pts ; coefficient de variation 1 pt ; estimation du coût et conclusion 3 pts.\\
\textbf{Points de vigilance :} distinguer soigneusement mode, médiane et moyenne ; pour la question 5, expliciter l'hypothèse de répartition uniforme à l'intérieur de la classe médiane (répartition proportionnelle) ; conclure par une phrase de décision argumentée (« l'enveloppe est suffisante car... »), attendue dans une épreuve de type Probatoire.
\end{corrige}

\newpage

\coursec{Fiche bilan}

\begin{popfigure}
\begin{tikzpicture}[mindmap, grow cyclic, every node/.style={concept, text=white, font=\small},
  root concept/.append style={concept color=popdark, font=\small\bfseries},
  level 1/.append style={level distance=3.4cm, sibling angle=60},
  level 2/.append style={level distance=2.2cm, sibling angle=45, font=\scriptsize}]
\node[root concept]{Séries statistiques regroupées en classes}
  child[concept color=popblue]{ node {Classes de même amplitude}
    child { node {Bornes $[a\,;b[$} }
    child { node {Centre $c_i=\dfrac{a+b}{2}$} }
  }
  child[concept color=poporange]{ node {Effectifs}
    child { node {$n_i$ par classe} }
    child { node {Cumulés croissants} }
  }
  child[concept color=popgreen]{ node {Fréquences}
    child { node {$f_i=\dfrac{n_i}{N}$} }
    child { node {Cumulées en \%} }
  }
  child[concept color=poppurple]{ node {Position}
    child { node {Classe modale} }
    child { node {Moyenne $\bar{x}$} }
    child { node {Médiane (interpolation)} }
  }
  child[concept color=poppink]{ node {Dispersion}
    child { node {Étendue} }
    child { node {Variance $V$, écart-type $\sigma$} }
  }
  child[concept color=popgold]{ node {Histogramme}
    child { node {Rectangles de même base} }
    child { node {Polygone des effectifs} }
  };
\end{tikzpicture}
\legende{Carte mentale de la leçon : séries statistiques regroupées en classes de même amplitude.}
\end{popfigure}

\begin{bilan}
\textbf{Définitions clés}
\begin{itemize}
  \item \cle{Classe} : intervalle $[a\,;b[$ dans lequel on regroupe les valeurs du caractère ; l'\cle{amplitude} est $b-a$ (ici la même pour toutes les classes).
  \item \cle{Centre de classe} : $c_i=\dfrac{a+b}{2}$ ; il représente toutes les valeurs de la classe dans les calculs.
  \item \cle{Effectif} $n_i$ : nombre d'individus de la classe ; effectif total $N=\sum n_i$.
  \item \cle{Effectif cumulé croissant} : somme des effectifs des classes situées en dessous de la borne supérieure.
  \item \cle{Fréquence} : $f_i=\dfrac{n_i}{N}$ (en \% : $f_i\times 100$) ; la somme des fréquences vaut $1$ (ou $100\,\%$).
  \item \cle{Fréquence cumulée croissante} : somme des fréquences jusqu'à la classe considérée.
\end{itemize}

\textbf{Formules à connaître par cœur}
\begin{center}
\begin{tabularx}{\linewidth}{|l|X|}
\hline
\rowcolor{popblueL} \textbf{Caractéristique} & \textbf{Formule / lecture} \\
\hline
Moyenne & $\bar{x}=\dfrac{1}{N}\sum n_i\,c_i$ (avec les centres $c_i$) \\
\hline
Classe modale & Classe d'effectif (ou de fréquence) maximal \\
\hline
Classe médiane & Première classe dont l'effectif cumulé atteint $\dfrac{N}{2}$ \\
\hline
Médiane (interpolation) & $Me=a+\dfrac{\dfrac{N}{2}-N_{\text{cum avant}}}{n_{Me}}\times \text{amplitude}$ \\
\hline
Étendue & $e=x_{\max}-x_{\min}$ (bornes extrêmes du tableau) \\
\hline
Variance & $V=\dfrac{1}{N}\sum n_i c_i^2-\bar{x}^2$ \\
\hline
Écart-type & $\sigma=\sqrt{V}$ \\
\hline
\end{tabularx}
\end{center}

\textbf{Méthodes express}
\begin{enumerate}
  \item \textbf{Construire le tableau} : classes de même amplitude $\to$ compter $n_i$ $\to$ vérifier $\sum n_i=N$ $\to$ compléter cumuls et fréquences.
  \item \textbf{Calculer $\bar{x}$, $V$, $\sigma$} : ajouter les colonnes $c_i$, $n_ic_i$, $n_ic_i^2$, sommer, appliquer les formules.
  \item \textbf{Tracer l'histogramme} : amplitudes égales $\Rightarrow$ hauteurs proportionnelles aux effectifs, rectangles accolés ; le \cle{polygone} relie les milieux des sommets des rectangles.
  \item \textbf{Interpréter} : comparer moyenne et médiane, commenter la dispersion avec $\sigma$, conclure en lien avec la situation (atelier, chantier, ventes...).
\end{enumerate}
\end{bilan}

\begin{aretenir}[Checklist avant le Probatoire]
\begin{itemize}
  \item[$\square$] Je sais regrouper des données en classes de même amplitude et compter les effectifs sans oublier la convention $[a\,;b[$.
  \item[$\square$] Je vérifie toujours que $\sum n_i=N$ et que la somme des fréquences vaut $1$ (ou $100\,\%$).
  \item[$\square$] Je sais calculer les effectifs cumulés croissants et les fréquences cumulées.
  \item[$\square$] Je sais calculer le centre $c_i$ de chaque classe.
  \item[$\square$] Je sais calculer la moyenne $\bar{x}$ avec les centres de classes.
  \item[$\square$] Je sais repérer la classe modale et la classe médiane.
  \item[$\square$] Je sais déterminer la médiane par interpolation linéaire.
  \item[$\square$] Je sais calculer la variance par $V=\dfrac{1}{N}\sum n_ic_i^2-\bar{x}^2$ puis l'écart-type $\sigma=\sqrt{V}$.
  \item[$\square$] Je sais construire un histogramme (rectangles accolés, mêmes bases) et le polygone des effectifs.
  \item[$\square$] Je sais lire un effectif ou une fréquence cumulée sur un tableau ou un graphique.
  \item[$\square$] Je rédige une conclusion claire qui interprète les résultats dans le contexte concret de l'énoncé.
  \item[$\square$] Je pense à donner les unités (FCFA, kg, cm, points...) dans mes réponses.
\end{itemize}
\end{aretenir}

\findecours

\end{document}
